% Generated mechanically from the frozen defos.tex target.
% Do not edit this generated file; edit the bound locale source instead.

% OK, start here.
%

\setcounter{chapter}{90}
\chapter{形变理论}



\phantomsection
\label{defos-section-phantom}


\section{引言}
\label{defos-section-introduction}

\noindent
本章旨在较为平缓地介绍模、态射等对象的形变理论。本章讨论能够利用
朴素余切复形证明的那些结果。在关于余切复形的一章中，我们还会略微推广
这些结果。希望深入了解的读者可以参阅 Illusie 关于这一主题的专著，见
\cite{cotangent}。





\section{环的形变与朴素余切复形}
\label{defos-section-deformations}

\noindent
本节利用朴素余切复形来处理一些形变理论问题。先取一个满环同态
$A' \to A$，其核为平方零理想 $I$。此外，设已给定环同态
$A \to B$、$B$-模 $N$ 以及 $A$-模同态 $c : I \to N$。
我们要问，能否找到一个对象来替代下图中的问号，使图表交换：
\begin{equation}
\label{defos-equation-to-solve}
\vcenter{
\xymatrix{
0 \ar[r] & N \ar[r] & {?} \ar[r] & B \ar[r] & 0 \\
0 \ar[r] & I \ar[u]^c \ar[r] & A' \ar[u] \ar[r] & A \ar[u] \ar[r] & 0
}
}
\end{equation}
并且，如果解存在，它在何种意义下是唯一的。更确切地说，我们寻找一个
满的 $A'$-代数同态 $B' \to B$，其核是与 $N$ 等同的平方零理想，
并且 $A' \to B'$ 诱导出给定同态 $c$。我们称 $B'$ 为
(\ref{defos-equation-to-solve}) 的一个{\it 解}。

\begin{lemma}
\label{defos-lemma-huge-diagram}
给定交换图
$$
\xymatrix{
& 0 \ar[r] & N_2 \ar[r] & B'_2 \ar[r] & B_2 \ar[r] & 0 \\
& 0 \ar[r]|\hole & I_2 \ar[u]_{c_2} \ar[r] &
A'_2 \ar[u] \ar[r]|\hole & A_2 \ar[u] \ar[r] & 0 \\
0 \ar[r] & N_1 \ar[ruu] \ar[r] & B'_1 \ar[r] & B_1 \ar[ruu] \ar[r] & 0 \\
0 \ar[r] & I_1 \ar[ruu]|\hole \ar[u]^{c_1} \ar[r] &
A'_1 \ar[ruu]|\hole \ar[u] \ar[r] & A_1 \ar[ruu]|\hole \ar[u] \ar[r] & 0
}
$$
其中前后两面均为 (\ref{defos-equation-to-solve}) 的解，则有：
\begin{enumerate}
\item 在
$\Ext^1_{B_1}(\NL_{B_1/A_1}, N_2)$
中存在一个典范元；该元为零，当且仅当存在环同态
$B'_1 \to B'_2$ 使上述图表交换。
\item 如果存在使图表交换的同态 $B'_1 \to B'_2$，那么所有这类
同态组成的集合是 $\Hom_{B_1}(\Omega_{B_1/A_1}, N_2)$ 作用下的
主齐性空间。
\end{enumerate}
\end{lemma}

\begin{proof}
令 $E$ 为 $B_1$ 的底层集合。考虑满同态 $A_1[E] \to B_1$，记其核为
$J$；朴素余切复形正是借助该同态按下式定义的：
$$
\NL_{B_1/A_1} = (J/J^2 \to \Omega_{A_1[E]/A_1} \otimes_{A_1[E]} B_1)
$$
参见《交换代数》第 \ref{algebra-section-netherlander} 节。由于
$\Omega_{A_1[E]/A_1} \otimes B_1$ 是自由 $B_1$-模，故
$$
\Ext^1_{B_1}(\NL_{B_1/A_1}, N_2) =
\frac{\Hom_{B_1}(J/J^2, N_2)}
{\Hom_{B_1}(\Omega_{A_1[E]/A_1} \otimes B_1, N_2)}
$$
我们将在右端的模中构造一个障碍。令
$J' = \Ker(A'_1[E] \to B_1)$。注意，存在满同态 $J' \to J$，
其核为 $I_1A'_1[E]$。对每个 $e \in E$，以 $x_e \in A_1[E]$
表示相应的变量。选取 $x_e$ 在 $B_1$ 中的像于 $B'_1$ 中的一个提升
$y_e$，并选取 $x_e$ 在 $B_2$ 中的像于 $B'_2$ 中的一个提升 $z_e$。
这些选择确定了 $A'_1$-代数同态
$$
A'_1[E] \to B'_1 \quad\text{且}\quad A'_1[E] \to B'_2
$$
其中第一个给出同态 $J' \to N_1,\ f' \mapsto f'(y_e)$，第二个给出
同态 $J' \to N_2,\ f' \mapsto f'(z_e)$。计算表明，这两个同态均在
$(J')^2$ 上为零。由于图表的左方块（涉及 $c_1$ 与 $c_2$）交换，
把它们视为取值于 $N_2$ 的同态时，它们在 $I_1A'_1[E]$ 上一致。
注意，$B'_1$ 是 $J' \to A'_1[E]$ 与 $J' \to N_1$ 的推出。因此，
若同态 $J' \to N_1 \to N_2$ 与 $J' \to N_2$ 一致，就得到一个
使图表交换的同态 $B'_1 \to B'_2$。于是，把下列同态的类定义为障碍：
$$
J/J^2 \to N_2,\quad f \mapsto f'(z_e) - \nu(f'(y_e))
$$
这里 $\nu : N_1 \to N_2$ 是给定同态，而 $f' \in J'$ 是 $f$ 的一个
提升。由上述讨论，该定义是良定的。注意，对某些
$\delta_{i,e} \in N_i$，我们可以把所选的 $z_e$ 改为
$z_e + \delta_{2,e}$，把 $y_e$ 改为 $y_e + \delta_{1,e}$。
这样一来，上述同态变为
% P11-E0001: frozen source formula retained; see control/ERRATA_CANDIDATES.jsonl.
$$
f \mapsto f'(z_e + \delta_{2, e}) - \nu(f'(y_e + \delta_{1, e})) =
f'(z_e) - \nu(f'(z_e)) +
\sum (\delta_{2, e} - \nu(\delta_{1, e}))\frac{\partial f}{\partial x_e}
$$
这恰好意味着，我们用复合同态
$J/J^2 \to \Omega_{A_1[E]/A_1} \otimes B_1 \to N_2$
来修正同态 $J/J^2 \to N_2$，其中第二个同态把 $\text{d}x_e$ 映为
$\delta_{2,e} - \nu(\delta_{1,e})$。因此，我们的障碍是良定义的，
并且它为零当且仅当存在提升。

\medskip\noindent
第 (2) 点来自如下观察：给定两个使图表交换的同态
$\varphi, \psi : B'_1 \to B'_2$，差 $\varphi - \psi$ 分解为一个同态
$D : B_1 \to N_2$，而它是一个 $A_1$-导子：
\begin{align*}
D(fg) & = \varphi(f'g') - \psi(f'g') \\
& =
\varphi(f')\varphi(g') - \psi(f')\psi(g') \\
& =
(\varphi(f') - \psi(f'))\varphi(g') + \psi(f')(\varphi(g') - \psi(g')) \\
& =
gD(f) + fD(g)
\end{align*}
因此，$D$ 对应于唯一的 $B_1$-线性同态
$\Omega_{B_1/A_1} \to N_2$。反过来，给定这样的线性同态便得到一个
导子 $D$；再给定一个使图表交换的环同态
$\psi : B'_1 \to B'_2$，则 $\psi + D$ 是另一个使图表交换的环同态。
\end{proof}

\begin{lemma}
\label{defos-lemma-choices}
如果 (\ref{defos-equation-to-solve}) 存在解，那么解的同构类集合是
$\Ext^1_B(\NL_{B/A}, N)$ 作用下的主齐性空间。
\end{lemma}

\begin{proof}
先注意，给定 (\ref{defos-equation-to-solve}) 的两个解 $B'_1$ 与 $B'_2$，
由引理 \ref{defos-lemma-huge-diagram} 可得一个障碍元
$o(B'_1, B'_2) \in \Ext^1_B(\NL_{B/A}, N)$；它阻碍同态
$B'_1 \to B'_2$ 的存在。显然，该元也是同构存在性的障碍，因而能够区分
不同的同构类。因此，为完成证明，只需说明：给定一个解 $B'$ 和一个元素
$\xi \in \Ext^1_B(\NL_{B/A}, N)$，可以找到另一个解 $B'_\xi$，使得
$o(B', B'_\xi) = \xi$。

\medskip\noindent
令 $E$ 为 $B$ 的底层集合。考虑满同态 $A[E] \to B$，记其核为 $J$；
朴素余切复形正是借助该同态按下式定义的：
$$
\NL_{B/A} = (J/J^2 \to \Omega_{A[E]/A} \otimes_{A[E]} B)
$$
参见《交换代数》第 \ref{algebra-section-netherlander} 节。由于
$\Omega_{A[E]/A} \otimes B$ 是自由 $B$-模，故
$$
\Ext^1_B(\NL_{B/A}, N) =
\frac{\Hom_B(J/J^2, N)}
{\Hom_B(\Omega_{A[E]/A} \otimes B, N)}
$$
因此，可以用一个同态 $\delta : J/J^2 \to N$ 的类来表示 $\xi$。

\medskip\noindent
对每个 $e \in E$，以 $x_e \in A[E]$ 表示相应的变量。选取 $x_e$ 在
$B$ 中的像于 $B'$ 中的一个提升 $y_e$。这些选择确定了一个
$A'$-代数同态 $\varphi : A'[E] \to B'$。令
$J' = \Ker(A'[E] \to B)$。注意，$\varphi$ 诱导同态
$\varphi|_{J'} : J' \to N$，且 $B'$ 是下图所示的推出：
$$
\xymatrix{
0 \ar[r] & N \ar[r] & B' \ar[r] & B \ar[r] & 0 \\
0 \ar[r] & J' \ar[u]^{\varphi|_{J'}} \ar[r] & A'[E] \ar[u] \ar[r] &
B \ar[u]_{=} \ar[r] & 0
}
$$
令 $\psi : J' \to N$ 为同态 $\varphi|_{J'}$ 与下列复合同态之和：
$$
J' \to J'/(J')^2 \to J/J^2 \xrightarrow{\delta} N.
$$
于是，沿 $\psi$ 所作的推出给出另一个环扩张 $B'_\xi$，它可置于上述
形式的图表中。计算表明 $o(B', B'_\xi) = \xi$，正如所需。
\end{proof}

\begin{lemma}
\label{defos-lemma-extensions-of-algebras}
设 $A$ 为环，$B$ 为 $A$-代数，$N$ 为 $B$-模。所有如下
$A$-代数扩张的同构类所成的集合
$$
0 \to N \to B' \to B \to 0
$$
（其中 $N$ 是平方零理想）与 $\Ext^1_B(\NL_{B/A}, N)$ 典范双射。
\end{lemma}

\begin{proof}
为证明这一点，把前述结果应用于 (\ref{defos-equation-to-solve}) 由下图给出的
情形：
$$
\xymatrix{
0 \ar[r] & N \ar[r] & {?} \ar[r] & B \ar[r] & 0 \\
0 \ar[r] & 0 \ar[u] \ar[r] & A \ar[u] \ar[r]^{\text{id}} & A \ar[u] \ar[r] & 0
}
$$
于是，本引理由引理 \ref{defos-lemma-choices} 以及解的存在性推出；一个解就是
$N \oplus B$。（该双射的直接构造见下述评注。）
\end{proof}

\begin{remark}
\label{defos-remark-extensions-of-algebras}
设 $A \to B$ 与 $N$ 如引理 \ref{defos-lemma-extensions-of-algebras} 所述。
令 $\alpha : P \to B$ 为 $B$ 在 $A$ 上的一个呈示，见《交换代数》
第 \ref{algebra-section-netherlander} 节。记 $J = \Ker(\alpha)$，
则与 $\alpha$ 相伴的朴素余切复形 $\NL(\alpha)$ 为复形
$J/J^2 \to \Omega_{P/A} \otimes_P B$。我们有
$$
\Ext^1_B(\NL(\alpha), N) =
\Coker\left(\Hom_B(\Omega_{P/A} \otimes_P B, N) \to \Hom_B(J/J^2, N)\right)
$$
因为 $\Omega_{P/A}$ 是自由模。考虑引理中所述的一个扩张
$0 \to N \to B' \to B \to 0$。由于 $P$ 是 $A$ 上的多项式代数，
可以把 $\alpha$ 提升为 $A$-代数同态 $\alpha' : P' \to B'$。
由于 $N$ 在 $B'$ 中平方为零，同态 $\alpha'|_J : J \to N$ 分解为
$J \to J/J^2 \to N$。该引理把这个扩张映到上述余核中此同态
$J/J^2 \to N$ 的类。
\end{remark}

\begin{lemma}
\label{defos-lemma-extensions-of-algebras-functorial}
给定环同态 $A \to B \to C$、$B$-模 $M$、$C$-模 $N$、$B$-线性同态
$c : M \to N$，以及如下具有平方零核的 $A$-代数扩张：
\begin{enumerate}
\item[(a)] 与 $\xi \in \Ext^1_B(\NL_{B/A}, M)$ 对应的
$0 \to M \to B' \to B \to 0$；
\item[(b)] 与 $\zeta \in \Ext^1_C(\NL_{C/A}, N)$ 对应的
$0 \to N \to C' \to C \to 0$。
\end{enumerate}
参见引理 \ref{defos-lemma-extensions-of-algebras}。那么，存在与 $B \to C$
及 $c$ 相容的 $A$-代数同态 $B' \to C'$，当且仅当 $\xi$ 与 $\zeta$
在 $\Ext^1_B(\NL_{B/A}, N)$ 中的像相同。
\end{lemma}

\begin{proof}
这个陈述是有意义的，因为我们有同态
$$
\Ext^1_B(\NL_{B/A}, M) \to \Ext^1_B(\NL_{B/A}, N)
$$
它由同态 $M \to N$ 诱导；还具有复合同态
$$
\Ext^1_C(\NL_{C/A}, N) \to \Ext^1_B(\NL_{C/A}, N) \to \Ext^1_B(\NL_{B/A}, N)
$$
其中第一个箭头使用限制同态 $D(C) \to D(B)$，第二个箭头使用复形的
典范同态 $\NL_{B/A} \to \NL_{C/A}$。把引理
\ref{defos-lemma-huge-diagram} 应用于下图，即可推出本引理的陈述：
$$
\xymatrix{
& 0 \ar[r] & N \ar[r] & C' \ar[r] & C \ar[r] & 0 \\
& 0 \ar[r]|\hole & 0 \ar[u] \ar[r] &
A \ar[u] \ar[r]|\hole & A \ar[u] \ar[r] & 0 \\
0 \ar[r] & M \ar[ruu] \ar[r] & B' \ar[r] & B \ar[ruu] \ar[r] & 0 \\
0 \ar[r] & 0 \ar[ruu]|\hole \ar[u] \ar[r] &
A \ar[ruu]|\hole \ar[u] \ar[r] & A \ar[ruu]|\hole \ar[u] \ar[r] & 0
}
$$
同时还要使用引理 \ref{defos-lemma-extensions-of-algebras} 与
\ref{defos-lemma-huge-diagram} 的证明中各项构造之间的一种相容性；其陈述和
证明从略。（直接论证见下述评注。）
\end{proof}

\begin{remark}
\label{defos-remark-extensions-of-algebras-functorial}
设 $A \to B \to C$、$M$、$N$、$c : M \to N$、
$0 \to M \to B' \to B \to 0$、$\xi \in \Ext^1_B(\NL_{B/A}, M)$、
$0 \to N \to C' \to C \to 0$ 以及
$\zeta \in \Ext^1_C(\NL_{C/A}, N)$ 如引理
\ref{defos-lemma-extensions-of-algebras-functorial} 所述。沿
$c : M \to N$ 作推出，可构造扩张
$$
\xymatrix{
0 \ar[r] & N \ar[r] & B'_1 \ar[r] & B \ar[r] & 0 \\
0 \ar[r] & M \ar[u]^c \ar[r] & B' \ar[u] \ar[r] & B \ar[u] \ar[r] & 0
}
$$
其中令 $B'_1 = (N \times B')/M$，而 $M$ 反对角嵌入。沿 $B \to C$
作拉回，可构造扩张
$$
\xymatrix{
0 \ar[r] & N \ar[r] & C' \ar[r] & C \ar[r] & 0 \\
0 \ar[r] & N \ar[u] \ar[r] & B'_2 \ar[u] \ar[r] & B \ar[u] \ar[r] & 0
}
$$
其中令 $B'_2 = C' \times_C B$（环的纤维积）。简单的追图表明，存在
与 $B \to C$ 及 $c$ 相容的 $A$-代数同态 $B' \to C'$，当且仅当把
$B'_1$ 与 $B'_2$ 都看作 $B$ 被 $N$ 扩张所得的 $A$-代数时二者同构。
因此，为证明引理 \ref{defos-lemma-extensions-of-algebras-functorial}，
只需说明：经引理 \ref{defos-lemma-extensions-of-algebras} 的双射，$B'_1$
对应于 $\xi$ 在同态
$\Ext^1_B(\NL_{B/A}, M) \to \Ext^1_B(\NL_{B/A}, N)$
下的像，并且 $B'_2$ 对应于 $\zeta$ 在同态
$\Ext^1_C(\NL_{C/A}, N) \to \Ext^1_B(\NL_{B/A}, N)$.
第一个陈述由评注 \ref{defos-remark-extensions-of-algebras} 中类的构造立即得到。
为证明第二个陈述，选取交换图
$$
\xymatrix{
Q \ar[r]_\beta & C \\
P \ar[u]^\varphi \ar[r]^\alpha & B \ar[u]
}
$$
其中各项都是 $A$-代数，$\alpha$ 是 $B$ 在 $A$ 上的一个呈示，而
$\beta$ 是 $C$ 在 $A$ 上的一个呈示。参见评注
\ref{defos-remark-extensions-of-algebras} 及其中的文献。令
$J = \Ker(\alpha)$、$K = \Ker(\beta)$。同态 $\varphi$ 诱导复形同态
$\NL(\alpha) \to \NL(\beta)$，特别地诱导
$\bar\varphi : J/J^2 \to K/K^2$。选取提升 $\beta$ 的 $A$-代数同态
$\beta' : Q \to C'$。于是
$\alpha' = (\beta' \circ \varphi, \alpha) : P \to B'_2 = C' \times_C B$
是 $\alpha$ 的一个提升。在这些选择下，由 $\beta'$ 诱导的同态
$K/K^2 \to N$ 与 $\bar\varphi : J/J^2 \to K/K^2$ 的复合，正是
$\alpha'$ 在 $J/J^2$ 上的限制。展开评注
\ref{defos-remark-extensions-of-algebras} 中各个类的构造，这确实表明
$B'_2$ 对应于 $\zeta$ 在同态
$\Ext^1_C(\NL_{C/A}, N) \to \Ext^1_B(\NL_{B/A}, N)$.
下的像。
\end{remark}

\begin{lemma}
\label{defos-lemma-parametrize-solutions}
设 $0 \to I \to A' \to A \to 0$、$A \to B$ 以及 $c : I \to N$
如 (\ref{defos-equation-to-solve}) 所述。记
$\xi \in \Ext^1_A(\NL_{A/A'}, I)$ 为经引理
\ref{defos-lemma-extensions-of-algebras} 与 $A$ 被 $I$ 扩张所得的 $A'$
相对应的元素。解的同构类集合与同态
$$
\Ext^1_B(\NL_{B/A'}, N) \to \Ext^1_A(\NL_{A/A'}, N)
$$
在 $\xi$ 的像上方的纤维典范双射。
\end{lemma}

\begin{proof}
把引理 \ref{defos-lemma-extensions-of-algebras} 应用于 $A' \to B$ 和
$B$-模 $N$，可见 $\Ext^1_B(\NL_{B/A'}, N)$ 中的元素 $\zeta$
参数化 $A'$-代数扩张 $0 \to N \to B' \to B \to 0$。把引理
\ref{defos-lemma-extensions-of-algebras-functorial} 应用于
$A' \to A \to B$ 和 $c : I \to N$，可见存在与 $c$ 及 $A \to B$
相容的 $A'$-代数同态 $A' \to B'$，当且仅当 $\zeta$ 映到 $\xi$。
这当然等价于说 $B'$ 是 (\ref{defos-equation-to-solve}) 的一个解。
\end{proof}

\begin{remark}
\label{defos-remark-parametrize-solutions}
注意，在引理 \ref{defos-lemma-parametrize-solutions} 的情形中有
$$
\Ext^1_A(\NL_{A/A'}, N) =
\Ext^1_B(\NL_{A/A'} \otimes_A^\mathbf{L} B, N) =
\Ext^1_B(\NL_{A/A'} \otimes_A B, N)
$$
第一个等号来自《更多交换代数》引理
\ref{more-algebra-lemma-tensor-hom-adjoint}，第二个等号来自
《更多交换代数》引理 \ref{more-algebra-lemma-tensor-NL}。我们有复形同态
$$
\NL_{A/A'} \otimes_A B \to \NL_{B/A'} \to \NL_{B/A}
$$
它们几乎构成一个可区别三角，见《交换代数》引理
\ref{algebra-lemma-exact-sequence-NL}。如果它确实是可区别三角，便可得出：
$\xi$ 在 $\Ext^2_B(\NL_{B/A}, N)$ 中的像，就是
(\ref{defos-equation-to-solve}) 存在解的障碍。
\end{remark}

\noindent
如果环同态 $A \to B$ 是局部完全交同态，那么解存在。这是一类提升结果；
注意，对于 syntomic 环同态，我们已在《环同态的光滑化》命题
\ref{smoothing-proposition-lift-smooth} 中证明了一个相当强的提升结果。

\begin{lemma}
\label{defos-lemma-existence-lci}
如果 $A \to B$ 是局部完全交环同态，那么
(\ref{defos-equation-to-solve}) 存在解。
\end{lemma}

\begin{proof}[第一种证明]
写成 $B = A[x_1, \ldots, x_n]/J$。由《更多交换代数》定义
\ref{more-algebra-definition-local-complete-intersection}，理想 $J$ 是
Koszul 正则的。这蕴含 $J$ 是 $H_1$-正则且拟正则的，
见《更多交换代数》第 \ref{more-algebra-section-ideals} 节。令
$J' \subset A'[x_1, \ldots, x_n]$ 为 $J$ 的原像。以
$I[x_1, \ldots, x_n]$ 表示同态
$A'[x_1, \ldots, x_n] \to A[x_1, \ldots, x_n]$ 的核。由
《更多交换代数》引理
\ref{more-algebra-lemma-conormal-sequence-H1-regular-ideal}，有
$I[x_1, \ldots, x_n] \cap (J')^2 = J'I[x_1, \ldots, x_n] =
JI[x_1, \ldots, x_n]$。因此得到短正合列
$$
0 \to I \otimes_A B \to J'/(J')^2 \to J/J^2 \to 0
$$
由于 $J/J^2$ 是投射模（《更多交换代数》引理
\ref{more-algebra-lemma-quasi-regular-ideal-finite-projective}），
可以选取该正合列的一个分裂
$$
J'/(J')^2 = I \otimes_A B \oplus J/J^2 
$$
令 $(J')^2 \subset J'' \subset J'$ 为映到上述分解中第二个直和项的那些
元素。于是
$$
0 \to I \otimes_A B \to A'[x_1, \ldots, x_n]/J'' \to B \to 0
$$
是 $N = I \otimes_A B$ 时 (\ref{defos-equation-to-solve}) 的一个解。
沿给定同态 $I \otimes_A B \to N$ 作推出，即得一般情形。
\end{proof}

\begin{proof}[第二种证明]
请先阅读评注 \ref{defos-remark-parametrize-solutions}。由《更多交换代数》引理
\ref{more-algebra-lemma-transitive-lci-at-end}，同态
$\NL_{A/A'} \otimes_A B \to \NL_{B/A'} \to \NL_{B/A}$
在 $D(B)$ 中确实构成一个可区别三角。因此，只需证明
$\Ext^2_B(\NL_{B/A}, N)$ 为零。由《更多交换代数》引理
\ref{more-algebra-lemma-lci-NL}，复形 $\NL_{B/A}$ 是 Tor 振幅位于
$[-1,0]$ 的完美复形。例如，由《更多交换代数》引理
\ref{more-algebra-lemma-splitting-unique} 第 (1) 点，这蕴含上述
$\Ext^2$ 为零。
\end{proof}







\section{环化空间的增厚}
\label{defos-section-thickenings-spaces}

\noindent
在下面几节中，我们将使用以下概念：
\begin{enumerate}
\item 环化空间 $(X', \mathcal{O}_{X'})$ 上的理想层
$\mathcal{I} \subset \mathcal{O}_{X'}$ 称为{\it 局部幂零的}，如果
$\mathcal{I}$ 的任意局部截面都是局部幂零的。与《交换代数》第
\ref{algebra-item-ideal-locally-nilpotent} 项比较。
\item 环化空间的一个{\it 增厚}是环化空间的态射
$i : (X, \mathcal{O}_X) \to (X', \mathcal{O}_{X'})$，满足：
\begin{enumerate}
\item $i$ 诱导同胚 $X \to X'$；
\item 同态 $i^\sharp : \mathcal{O}_{X'} \to i_*\mathcal{O}_X$ 是满的；
\item $i^\sharp$ 的核是局部幂零理想层。
\end{enumerate}
\item 环化空间的一个{\it 一阶增厚}是环化空间的增厚
$i : (X, \mathcal{O}_X) \to (X', \mathcal{O}_{X'})$，且
$\Ker(i^\sharp)$ 平方为零。
\item 显然可以定义{\it 增厚的态射}、{\it 基环化空间上增厚的态射}等概念。
\end{enumerate}
若 $i : (X, \mathcal{O}_X) \to (X', \mathcal{O}_{X'})$ 是环化空间的
一个增厚，我们便等同其底层拓扑空间，并把 $\mathcal{O}_X$、
$\mathcal{O}_{X'}$ 以及 $\mathcal{I} = \Ker(i^\sharp)$ 都视为
$X = X'$ 上的层。由此得到 $\mathcal{O}_{X'}$-模的短正合列
$$
0 \to \mathcal{I} \to \mathcal{O}_{X'} \to \mathcal{O}_X \to 0
$$
由《模层》引理 \ref{modules-lemma-i-star-equivalence}，
$\mathcal{O}_X$-模范畴等价于被 $\mathcal{I}$ 零化的
$\mathcal{O}_{X'}$-模范畴。特别地，如果 $i$ 是一阶增厚，那么
$\mathcal{I}$ 是一个 $\mathcal{O}_X$-模。

\begin{situation}
\label{defos-situation-morphism-thickenings}
增厚的一个态射 $(f,f')$ 由如下交换图给出：
\begin{equation}
\label{defos-equation-morphism-thickenings}
\vcenter{
\xymatrix{
(X, \mathcal{O}_X) \ar[r]_i \ar[d]_f & (X', \mathcal{O}_{X'}) \ar[d]^{f'} \\
(S, \mathcal{O}_S) \ar[r]^t & (S', \mathcal{O}_{S'})
}
}
\end{equation}
其中各项为环化空间，水平箭头为增厚。在这一情形中，令
$\mathcal{I} = \Ker(i^\sharp) \subset \mathcal{O}_{X'}$、
$\mathcal{J} = \Ker(t^\sharp) \subset \mathcal{O}_{S'}$。由于在底层
拓扑空间上 $f=f'$，我们等同（拓扑）逆像函子 $f^{-1}$ 与
$(f')^{-1}$。注意，
$(f')^\sharp : f^{-1}\mathcal{O}_{S'} \to \mathcal{O}_{X'}$
特别地诱导同态 $f^{-1}\mathcal{J} \to \mathcal{I}$，从而诱导
$\mathcal{O}_{X'}$-模同态
$$
(f')^*\mathcal{J} \longrightarrow \mathcal{I}
$$
如果 $i$ 与 $t$ 都是一阶增厚，那么
$(f')^*\mathcal{J}=f^*\mathcal{J}$，上述同态就成为
$f^*\mathcal{J} \to \mathcal{I}$。
\end{situation}

\begin{definition}
\label{defos-definition-strict-morphism-thickenings}
在情形 \ref{defos-situation-morphism-thickenings} 中，如果同态
$(f')^*\mathcal{J} \longrightarrow \mathcal{I}$ 是满的，则称
$(f,f')$ 是一个{\it 增厚的严格态射}。
\end{definition}

\noindent
下述引理特别说明：概形增厚的态射
$(f,f'):(X\subset X')\to(S\subset S')$ 是严格的，当且仅当
$X=S\times_{S'}X'$。

\begin{lemma}
\label{defos-lemma-strict-morphism-thickenings}
在情形 \ref{defos-situation-morphism-thickenings} 中，态射 $(f,f')$
是增厚的严格态射，当且仅当
(\ref{defos-equation-morphism-thickenings}) 在环化空间范畴中是笛卡儿的。
\end{lemma}

\begin{proof}
从略。
\end{proof}




\section{环化空间一阶增厚上的模}
\label{defos-section-modules-thickenings}

\noindent
本节讨论模的形变理论所需的一些预备知识。设
$i : (X,\mathcal{O}_X)\to(X',\mathcal{O}_{X'})$ 是环化空间的一阶增厚。
我们将随时使用第 \ref{defos-section-thickenings-spaces} 节引入的记号；
特别地，我们等同底层拓扑空间。本节考虑 $\mathcal{O}_{X'}$-模的短正合列
\begin{equation}
\label{defos-equation-extension}
0 \to \mathcal{K} \to \mathcal{F}' \to \mathcal{F} \to 0
\end{equation}
其中 $\mathcal{F}$、$\mathcal{K}$ 是 $\mathcal{O}_X$-模，而
$\mathcal{F}'$ 是 $\mathcal{O}_{X'}$-模。在这种情形中，有典范的
$\mathcal{O}_X$-模同态
$$
c_{\mathcal{F}'} :
\mathcal{I} \otimes_{\mathcal{O}_X} \mathcal{F}
\longrightarrow
\mathcal{K}
$$
其中 $\mathcal{I}=\Ker(i^\sharp)$。具体而言，给定 $\mathcal{I}$ 的
局部截面 $f$ 与 $\mathcal{F}$ 的局部截面 $s$，令
$c_{\mathcal{F}'}(f\otimes s)=fs'$，其中 $s'$ 是提升 $s$ 的
$\mathcal{F}'$ 的局部截面。

\begin{lemma}
\label{defos-lemma-inf-map}
设 $i : (X,\mathcal{O}_X)\to(X',\mathcal{O}_{X'})$ 是环化空间的一阶
增厚，并给定扩张
$$
0 \to \mathcal{K} \to \mathcal{F}' \to \mathcal{F} \to 0
\quad\text{且}\quad
0 \to \mathcal{L} \to \mathcal{G}' \to \mathcal{G} \to 0
$$
如 (\ref{defos-equation-extension}) 所述，以及同态
$\varphi:\mathcal{F}\to\mathcal{G}$ 与
$\psi:\mathcal{K}\to\mathcal{L}$。
\begin{enumerate}
\item 如果存在与 $\varphi$ 和 $\psi$ 相容的 $\mathcal{O}_{X'}$-模同态
$\varphi':\mathcal{F}'\to\mathcal{G}'$，那么图表
$$
\xymatrix{
\mathcal{I} \otimes_{\mathcal{O}_X} \mathcal{F}
\ar[r]_-{c_{\mathcal{F}'}} \ar[d]_{1 \otimes \varphi} &
\mathcal{K} \ar[d]^\psi \\
\mathcal{I} \otimes_{\mathcal{O}_X} \mathcal{G}
\ar[r]^-{c_{\mathcal{G}'}} &
\mathcal{L}
}
$$
交换。
\item 与 $\varphi$ 和 $\psi$ 相容的 $\mathcal{O}_{X'}$-模同态
$\varphi':\mathcal{F}'\to\mathcal{G}'$ 的集合，如果非空，就是
$\Hom_{\mathcal{O}_X}(\mathcal{F},\mathcal{L})$ 作用下的主齐性空间。
\end{enumerate}
\end{lemma}

\begin{proof}
第 (1) 点由各同态的描述立即得到。对于第 (2) 点，若 $\varphi'$ 与
$\varphi''$ 是两个与 $\varphi$ 和 $\psi$ 相容的同态
$\mathcal{F}'\to\mathcal{G}'$，那么 $\varphi'-\varphi''$ 分解为
$$
\mathcal{F}' \to \mathcal{F} \to \mathcal{L} \to \mathcal{G}'
$$
由《模层》引理 \ref{modules-lemma-i-star-equivalence}，中间的同态来自
$\Hom_{\mathcal{O}_X}(\mathcal{F},\mathcal{L})$ 中唯一的元素。反过来，
给定该群中的元素 $\alpha$，可以把上述复合（中间同态取 $\alpha$）
加到 $\varphi'$ 上。若干细节从略。
\end{proof}

\begin{lemma}
\label{defos-lemma-inf-obs-map}
设 $i:(X,\mathcal{O}_X)\to(X',\mathcal{O}_{X'})$ 是环化空间的一阶
增厚，并给定扩张
$$
0 \to \mathcal{K} \to \mathcal{F}' \to \mathcal{F} \to 0
\quad\text{且}\quad
0 \to \mathcal{L} \to \mathcal{G}' \to \mathcal{G} \to 0
$$
如 (\ref{defos-equation-extension}) 所述，以及同态
$\varphi:\mathcal{F}\to\mathcal{G}$ 与
$\psi:\mathcal{K}\to\mathcal{L}$。假设图表
$$
\xymatrix{
\mathcal{I} \otimes_{\mathcal{O}_X} \mathcal{F}
\ar[r]_-{c_{\mathcal{F}'}} \ar[d]_{1 \otimes \varphi} &
\mathcal{K} \ar[d]^\psi \\
\mathcal{I} \otimes_{\mathcal{O}_X} \mathcal{G}
\ar[r]^-{c_{\mathcal{G}'}} &
\mathcal{L}
}
$$
交换。那么存在元素
$$
o(\varphi, \psi) \in
\Ext^1_{\mathcal{O}_X}(\mathcal{F}, \mathcal{L})
$$
它为零，当且仅当存在与 $\varphi$ 和 $\psi$ 相容的同态
$\varphi':\mathcal{F}'\to\mathcal{G}'$。
\end{lemma}

\begin{proof}
可以显式构造扩张
$$
0 \to \mathcal{L} \to \mathcal{H} \to \mathcal{F} \to 0
$$
方法是令 $\mathcal{H}$ 为复形
$$
\mathcal{K}
\xrightarrow{1, - \psi}
\mathcal{F}' \oplus \mathcal{G}' \xrightarrow{\varphi, 1}
\mathcal{G}
$$
在中间项处的上同调层（记号含义显然）。利用引理中图表交换这一假设，
对局部截面作计算可知 $\mathcal{H}$ 被 $\mathcal{I}$ 零化。因此
$\mathcal{H}$ 在下述群中定义了一个类：
$$
\Ext^1_{\mathcal{O}_X}(\mathcal{F}, \mathcal{L})
\subset
\Ext^1_{\mathcal{O}_{X'}}(\mathcal{F}, \mathcal{L})
$$
最后，$\mathcal{H}$ 的类等于沿 $\psi$ 对扩张 $\mathcal{F}'$ 所作的
推出与沿 $\varphi$ 对扩张 $\mathcal{G}'$ 所作的拉回之差（计算从略）。
因此，$\mathcal{H}$ 的类为零，等价于存在交换图
$$
\xymatrix{
0 \ar[r] &
\mathcal{K} \ar[r] \ar[d]_{\psi} &
\mathcal{F}' \ar[r] \ar[d]_{\varphi'} &
\mathcal{F} \ar[r] \ar[d]_\varphi & 0\\
0 \ar[r] &
\mathcal{L} \ar[r] &
\mathcal{G}' \ar[r] &
\mathcal{G} \ar[r] & 0
}
$$
这正是所需。
\end{proof}

\begin{lemma}
\label{defos-lemma-inf-ext}
设 $i:(X,\mathcal{O}_X)\to(X',\mathcal{O}_{X'})$ 是环化空间的一阶
增厚。给定 $\mathcal{O}_X$-模 $\mathcal{F}$、$\mathcal{K}$ 以及
$\mathcal{O}_X$-线性同态
$c:\mathcal{I}\otimes_{\mathcal{O}_X}\mathcal{F}\to\mathcal{K}$。
如果存在满足 $c_{\mathcal{F}'}=c$ 的正合列
(\ref{defos-equation-extension})，那么这些扩张的同构类集合是
$\Ext^1_{\mathcal{O}_X}(\mathcal{F},\mathcal{K})$ 作用下的主齐性空间。
\end{lemma}

\begin{proof}
假设给定扩张
$$
0 \to \mathcal{K} \to \mathcal{F}'_1 \to \mathcal{F} \to 0
\quad\text{且}\quad
0 \to \mathcal{K} \to \mathcal{F}'_2 \to \mathcal{F} \to 0
$$
且 $c_{\mathcal{F}'_1}=c_{\mathcal{F}'_2}=c$。那么二者之差
（在扩张群中取差，见《同调代数》第 \ref{homology-section-extensions} 节）
是扩张
$$
0 \to \mathcal{K} \to \mathcal{E} \to \mathcal{F} \to 0
$$
其中 $\mathcal{E}$ 被 $\mathcal{I}$ 零化（局部计算从略）。所以该正合列
是 $\mathcal{O}_X$-模的扩张，见《模层》引理
\ref{modules-lemma-i-star-equivalence}。反过来，给定这样的扩张
$\mathcal{E}$，可将它加到 $\mathcal{O}_{X'}$-扩张 $\mathcal{F}'$ 上，
而不改变同态 $c_{\mathcal{F}'}$。若干细节从略。
\end{proof}

\begin{lemma}
\label{defos-lemma-inf-obs-ext}
设 $i:(X,\mathcal{O}_X)\to(X',\mathcal{O}_{X'})$ 是环化空间的一阶
增厚。给定 $\mathcal{O}_X$-模 $\mathcal{F}$、$\mathcal{K}$ 以及
$\mathcal{O}_X$-线性同态
$c:\mathcal{I}\otimes_{\mathcal{O}_X}\mathcal{F}\to\mathcal{K}$。
那么存在元素
$$
o(\mathcal{F}, \mathcal{K}, c) \in
\Ext^2_{\mathcal{O}_X}(\mathcal{F}, \mathcal{K})
$$
它为零，当且仅当存在满足 $c_{\mathcal{F}'}=c$ 的正合列
(\ref{defos-equation-extension})。
\end{lemma}

\begin{proof}
先证明：如果 $\mathcal{K}$ 是内射 $\mathcal{O}_X$-模，那么确实存在
满足 $c_{\mathcal{F}'}=c$ 的正合列 (\ref{defos-equation-extension})。
为此，选取平坦 $\mathcal{O}_{X'}$-模 $\mathcal{H}'$ 以及满同态
$\mathcal{H}'\to\mathcal{F}$（《模层》引理
\ref{modules-lemma-module-quotient-flat}）。令
$\mathcal{J}\subset\mathcal{H}'$ 为其核。由于 $\mathcal{H}'$ 平坦，故有
$$
\mathcal{I} \otimes_{\mathcal{O}_{X'}} \mathcal{H}' =
\mathcal{I}\mathcal{H}'
\subset \mathcal{J} \subset \mathcal{H}'
$$
注意，同态
$$
\mathcal{I}\mathcal{H}' =
\mathcal{I} \otimes_{\mathcal{O}_{X'}} \mathcal{H}'
\longrightarrow
\mathcal{I} \otimes_{\mathcal{O}_{X'}} \mathcal{F} =
\mathcal{I} \otimes_{\mathcal{O}_X} \mathcal{F}
$$
在 $\mathcal{I}\mathcal{J}$ 上为零。事实上，如果 $f$ 是
$\mathcal{I}$ 的局部截面、$s$ 是 $\mathcal{H}$ 的局部截面，那么
$fs$ 映到 $f\otimes\overline{s}$，其中 $\overline{s}$ 是 $s$ 在
$\mathcal{F}$ 中的像。因此得到如下 $\mathcal{O}_X$-模图表：
$$
\xymatrix{
\mathcal{I}\mathcal{H}'/\mathcal{I}\mathcal{J}
\ar@{^{(}->}[r] \ar[d] &
\mathcal{J}/\mathcal{I}\mathcal{J} \ar@{..>}[d]_\gamma \\
\mathcal{I} \otimes_{\mathcal{O}_X} \mathcal{F} \ar[r]^-c &
\mathcal{K}
}
$$
如果 $\mathcal{K}$ 是内射 $\mathcal{O}_X$-模，就得到虚线箭头。
以 $\gamma':\mathcal{J}\to\mathcal{K}$ 表示 $\gamma$ 与
$\mathcal{J}\to\mathcal{J}/\mathcal{I}\mathcal{J}$ 的复合。
局部计算表明，下列推出
$$
\xymatrix{
0 \ar[r] &
\mathcal{J} \ar[r] \ar[d]_{\gamma'} &
\mathcal{H}' \ar[r] \ar[d] &
\mathcal{F} \ar[r] \ar@{=}[d] &
0 \\
0 \ar[r] &
\mathcal{K} \ar[r] &
\mathcal{F}' \ar[r] &
\mathcal{F} \ar[r] &
0
}
$$
是本引理所提问题的一个解。

\medskip\noindent
一般情形。选取嵌入 $\mathcal{K}\subset\mathcal{K}'$，其中
$\mathcal{K}'$ 是内射 $\mathcal{O}_X$-模。令 $\mathcal{Q}$ 为商模，
于是有正合列
$$
0 \to \mathcal{K} \to \mathcal{K}' \to \mathcal{Q} \to 0
$$
以
$c':\mathcal{I}\otimes_{\mathcal{O}_X}\mathcal{F}\to\mathcal{K}'$
表示相应的复合。由上一段，存在正合列
$$
0 \to \mathcal{K}' \to \mathcal{E}' \to \mathcal{F} \to 0
$$
如 (\ref{defos-equation-extension}) 所述，且 $c_{\mathcal{E}'}=c'$。
注意，$c'$ 与同态 $\mathcal{K}'\to\mathcal{Q}$ 的复合为零；因此，
沿 $\mathcal{K}'\to\mathcal{Q}$ 对 $\mathcal{E}'$ 作推出所得的是扩张
$$
0 \to \mathcal{Q} \to \mathcal{D}' \to \mathcal{F} \to 0
$$
如 (\ref{defos-equation-extension}) 所述，且 $c_{\mathcal{D}'}=0$。
这恰好意味着 $\mathcal{D}'$ 被 $\mathcal{I}$ 零化；换言之，
$\mathcal{D}'$ 是 $\mathcal{O}_X$-模的扩张，从而定义元素
$$
o(\mathcal{F}, \mathcal{K}, c) \in
\Ext^1_{\mathcal{O}_X}(\mathcal{F}, \mathcal{Q}) =
\Ext^2_{\mathcal{O}_X}(\mathcal{F}, \mathcal{K})
$$
（等号来自与上述正合列相伴的长正合上同调列，以及取值于内射模
$\mathcal{K}'$ 的高次 Ext 群为零）。如果
$o(\mathcal{F},\mathcal{K},c)=0$，便可选取分裂
$s:\mathcal{F}\to\mathcal{D}'$，并令
$$
\mathcal{F}' = \Ker(\mathcal{E}' \to \mathcal{D}'/s(\mathcal{F}))
$$
于是得到下图：
$$
\xymatrix{
0 \ar[r] &
\mathcal{K} \ar[r] \ar[d] &
\mathcal{F}' \ar[r] \ar[d] &
\mathcal{F} \ar[r] \ar@{=}[d] &
0 \\
0 \ar[r] &
\mathcal{K}' \ar[r] &
\mathcal{E}' \ar[r] &
\mathcal{F} \ar[r] & 0
}
$$
两行均正合，这说明 $c_{\mathcal{F}'}=c$。反过来，如果
$\mathcal{F}'$ 存在，那么由引理 \ref{defos-lemma-inf-ext} 以及取值于内射模
$\mathcal{K}'$ 的高次 Ext 群为零，沿 $\mathcal{K}\to\mathcal{K}'$
对 $\mathcal{F}'$ 作推出所得的扩张同构于 $\mathcal{E}'$。由此得到
上述形式的图表，它蕴含 $\mathcal{D}'$ 作为扩张是分裂的；也就是说，
类 $o(\mathcal{F},\mathcal{K},c)$ 为零。
\end{proof}

\begin{remark}
\label{defos-remark-trivial-thickening}
设 $(X,\mathcal{O}_X)$ 为环化空间。如果存在环化空间的态射
$\pi:(X',\mathcal{O}_{X'})\to(X,\mathcal{O}_X)$，它是 $i$ 的左逆，
则称一阶增厚
$i:(X,\mathcal{O}_X)\to(X',\mathcal{O}_{X'})$ 是{\it 平凡的}。
这样的态射 $\pi$ 的一个选择称为该一阶增厚的一个{\it 平凡化}。
给定 $\pi$，利用 $\pi^\sharp$ 分裂满同态
$\mathcal{O}_{X'}\to\mathcal{O}_X$，便得到 $X$ 上代数层的分裂
\begin{equation}
\label{defos-equation-splitting}
\mathcal{O}_{X'} = \mathcal{O}_X \oplus \mathcal{I}
\end{equation}
反过来，这样的分裂确定一个态射 $\pi$。环化空间
$(X,\mathcal{O}_X)$ 的带平凡化一阶增厚所成的范畴，等价于
$\mathcal{O}_X$-模范畴。
\end{remark}

\begin{remark}
\label{defos-remark-trivial-extension}
设 $i:(X,\mathcal{O}_X)\to(X',\mathcal{O}_{X'})$ 是环化空间的平凡
一阶增厚，而 $\pi:(X',\mathcal{O}_{X'})\to(X,\mathcal{O}_X)$ 是一个
平凡化。给定任意三元组 $(\mathcal{F},\mathcal{K},c)$，其中
$\mathcal{F}$、$\mathcal{K}$ 是 $\mathcal{O}_X$-模，且
$c:\mathcal{I}\otimes_{\mathcal{O}_X}\mathcal{F}\to\mathcal{K}$，
可以令
$$
\mathcal{F}'_{c, triv} = \mathcal{F} \oplus \mathcal{K}
$$
再用与 $\pi$ 相伴的分裂 (\ref{defos-equation-splitting}) 和同态 $c$
定义 $\mathcal{O}_{X'}$-模结构，从而得到扩张
(\ref{defos-equation-extension})。我们称 $\mathcal{F}'_{c,triv}$ 为
由 $c$ 与平凡化 $\pi$ 所对应的 $\mathcal{F}$ 被 $\mathcal{K}$
扩张所得的{\it 平凡扩张}。给定 (\ref{defos-equation-extension}) 中的任意
扩张 $\mathcal{F}'$，可以利用
$\pi^\sharp:\mathcal{O}_X\to\mathcal{O}_{X'}$ 把 $\mathcal{F}'$
看成 $\mathcal{O}_X$-模的扩张，因而得到
$\Ext^1_{\mathcal{O}_X}(\mathcal{F},\mathcal{K})$ 中的类
$\xi_{\mathcal{F}'}$。引理 \ref{defos-lemma-inf-ext} 保证映射
$\mathcal{F}' \mapsto \xi_{\mathcal{F}'}$
诱导双射
$$
\left\{
\begin{matrix}
\text{扩张的同构类}\\
\mathcal{F}'\text{ 如 (\ref{defos-equation-extension}) 所述，且 }c=c_{\mathcal{F}'}
\end{matrix}
\right\}
\longrightarrow
\Ext^1_{\mathcal{O}_X}(\mathcal{F}, \mathcal{K})
$$
此外，平凡扩张 $\mathcal{F}'_{c,triv}$ 映到零类。
\end{remark}

\begin{remark}
\label{defos-remark-extension-functorial}
设 $(X,\mathcal{O}_X)$ 为环化空间，并设
$(X,\mathcal{O}_X)\to(X'_i,\mathcal{O}_{X'_i})$（$i=1,2$）为一阶
增厚，其理想层为 $\mathcal{I}_i$。设
$h:(X'_1,\mathcal{O}_{X'_1})\to(X'_2,\mathcal{O}_{X'_2})$
为 $(X,\mathcal{O}_X)$ 的一阶增厚之间的态射。图示如下：
$$
\xymatrix{
& (X, \mathcal{O}_X) \ar[ld] \ar[rd] & \\
(X'_1, \mathcal{O}_{X'_1}) \ar[rr]^h & & 
(X'_2, \mathcal{O}_{X'_2})
}
$$
注意，$h^\sharp:\mathcal{O}_{X'_2}\to\mathcal{O}_{X'_1}$ 特别地诱导
$\mathcal{O}_X$-模同态 $\mathcal{I}_2\to\mathcal{I}_1$。设
$\mathcal{F}$ 为 $\mathcal{O}_X$-模。对 $i=1,2$，设
$(\mathcal{K}_i,c_i)$ 由 $\mathcal{O}_X$-模 $\mathcal{K}_i$ 及同态
$c_i:\mathcal{I}_i\otimes_{\mathcal{O}_X}\mathcal{F}\to\mathcal{K}_i$
组成。再给定 $\mathcal{O}_X$-模同态
$\mathcal{K}_2\to\mathcal{K}_1$，使得
$$
\xymatrix{
\mathcal{I}_2 \otimes_{\mathcal{O}_X} \mathcal{F}
\ar[r]_-{c_2} \ar[d] &
\mathcal{K}_2 \ar[d] \\
\mathcal{I}_1 \otimes_{\mathcal{O}_X} \mathcal{F}
\ar[r]^-{c_1} &
\mathcal{K}_1
}
$$
交换。那么存在如下典范函子性：
$$
\left\{
\begin{matrix}
\mathcal{F}'_2\text{ 如 (\ref{defos-equation-extension}) 所述，且}\\
c_2=c_{\mathcal{F}'_2}\text{、}\mathcal{K}=\mathcal{K}_2
\end{matrix}
\right\}
\longrightarrow
\left\{
\begin{matrix}
\mathcal{F}'_1\text{ 如 (\ref{defos-equation-extension}) 所述，且}\\
c_1=c_{\mathcal{F}'_1}\text{、}\mathcal{K}=\mathcal{K}_1
\end{matrix}
\right\}
$$
具体而言，把 $\mathcal{O}_X$、$\mathcal{O}_{X'_i}$、$\mathcal{F}$、
$\mathcal{K}_i$ 等都视为 $X$ 上的层；给定 $\mathcal{F}'_2$ 后，
令层 $\mathcal{F}'_1$ 为推出，也就是说，它置于下列扩张图表中：
$$
\xymatrix{
0 \ar[r] &
\mathcal{K}_2 \ar[r] \ar[d] &
\mathcal{F}'_2 \ar[r] \ar[d] &
\mathcal{F} \ar@{=}[d] \ar[r] & 0 \\
0 \ar[r] &
\mathcal{K}_1 \ar[r] &
\mathcal{F}'_1 \ar[r] &
\mathcal{F} \ar[r] & 0
}
$$
推出上的 $\mathcal{O}_{X'_1}$-模结构之构造从略（该构造使用涉及
$c_1$ 与 $c_2$ 的图表的交换性）。
\end{remark}

\begin{remark}
\label{defos-remark-trivial-extension-functorial}
设 $(X,\mathcal{O}_X)$、
$(X,\mathcal{O}_X)\to(X'_i,\mathcal{O}_{X'_i})$、$\mathcal{I}_i$ 以及
$h:(X'_1,\mathcal{O}_{X'_1})\to(X'_2,\mathcal{O}_{X'_2})$ 如评注
\ref{defos-remark-extension-functorial} 所述。假设给定平凡化
$\pi_i:X'_i\to X$，满足 $\pi_1=h\circ\pi_2$。换言之，假设 $h$ 是
$(X,\mathcal{O}_X)$ 的带平凡化一阶增厚之间的态射。对 $i=1,2$，
设 $(\mathcal{K}_i,c_i)$ 由 $\mathcal{O}_X$-模 $\mathcal{K}_i$
及同态
$c_i:\mathcal{I}_i\otimes_{\mathcal{O}_X}\mathcal{F}\to\mathcal{K}_i$
组成。再给定 $\mathcal{O}_X$-模同态
$\mathcal{K}_2\to\mathcal{K}_1$，使得
$$
\xymatrix{
\mathcal{I}_2 \otimes_{\mathcal{O}_X} \mathcal{F}
\ar[r]_-{c_2} \ar[d] &
\mathcal{K}_2 \ar[d] \\
\mathcal{I}_1 \otimes_{\mathcal{O}_X} \mathcal{F}
\ar[r]^-{c_1} &
\mathcal{K}_1
}
$$
交换。在这种情形下，评注 \ref{defos-remark-trivial-extension} 的构造诱导
交换图
$$
\xymatrix{
\{\mathcal{F}'_2\text{ 如 (\ref{defos-equation-extension}) 所述，且 }
c_2=c_{\mathcal{F}'_2}\text{、}\mathcal{K}=\mathcal{K}_2\}
\ar[d] \ar[rr] & &
\Ext^1_{\mathcal{O}_X}(\mathcal{F}, \mathcal{K}_2) \ar[d] \\
\{\mathcal{F}'_1\text{ 如 (\ref{defos-equation-extension}) 所述，且 }
c_1=c_{\mathcal{F}'_1}\text{、}\mathcal{K}=\mathcal{K}_1\}
\ar[rr] & &
\Ext^1_{\mathcal{O}_X}(\mathcal{F}, \mathcal{K}_1)
}
$$
其中右侧竖直同态由 $\Ext$ 的函子性与同态
$\mathcal{K}_2\to\mathcal{K}_1$ 给出，左侧竖直映射则是评注
\ref{defos-remark-extension-functorial} 中的映射。
\end{remark}

\begin{remark}
\label{defos-remark-short-exact-sequence-thickenings}
设 $(X,\mathcal{O}_X)$ 为环化空间。对于 $(X,\mathcal{O}_X)$ 的一阶
增厚态射列
$$
(X'_1, \mathcal{O}_{X'_1}) \to
(X'_2, \mathcal{O}_{X'_2}) \to
(X'_3, \mathcal{O}_{X'_3})
$$
如果理想层 $\mathcal{I}_i$ 之间的相应同态给出
$\mathcal{O}_X$-模复形
$\mathcal{I}_3\to\mathcal{I}_2\to\mathcal{I}_1$
（即复合为零），则称该态射列为一个{\it 复形}。在这种情形下，复合
$(X'_1,\mathcal{O}_{X'_1})\to(X'_3,\mathcal{O}_{X'_3})$ 经
$(X,\mathcal{O}_X)\to(X'_3,\mathcal{O}_{X'_3})$ 分解；换言之，
$(X,\mathcal{O}_X)$ 的一阶增厚 $(X'_1,\mathcal{O}_{X'_1})$
是平凡的，并且带有典范平凡化
$\pi:(X'_1,\mathcal{O}_{X'_1})\to(X,\mathcal{O}_X)$。

\medskip\noindent
如果 $(X,\mathcal{O}_X)$ 的一阶增厚态射列
$$
(X'_1, \mathcal{O}_{X'_1}) \to
(X'_2, \mathcal{O}_{X'_2}) \to
(X'_3, \mathcal{O}_{X'_3})
$$
所对应的理想层同态构成 $\mathcal{O}_X$-模的短正合列
$$
0 \to \mathcal{I}_3 \to \mathcal{I}_2 \to \mathcal{I}_1 \to 0
$$
则称该态射列为一个{\it 短正合列}。
\end{remark}

\begin{remark}
\label{defos-remark-complex-thickenings-and-ses-modules}
设 $(X,\mathcal{O}_X)$ 为环化空间，$\mathcal{F}$ 为
$\mathcal{O}_X$-模。设
$$
(X'_1, \mathcal{O}_{X'_1}) \to
(X'_2, \mathcal{O}_{X'_2}) \to
(X'_3, \mathcal{O}_{X'_3})
$$
是 $(X,\mathcal{O}_X)$ 的一阶增厚所成的复形，见评注
\ref{defos-remark-short-exact-sequence-thickenings}。对 $i=1,2,3$，设
$(\mathcal{K}_i,c_i)$ 由 $\mathcal{O}_X$-模 $\mathcal{K}_i$ 与同态
$c_i:\mathcal{I}_i\otimes_{\mathcal{O}_X}\mathcal{F}\to\mathcal{K}_i$
组成。假设给定 $\mathcal{O}_X$-模的短正合列
$$
0 \to \mathcal{K}_3 \to \mathcal{K}_2 \to \mathcal{K}_1 \to 0
$$
使得图表
$$
\vcenter{
\xymatrix{
\mathcal{I}_2 \otimes_{\mathcal{O}_X} \mathcal{F}
\ar[r]_-{c_2} \ar[d] &
\mathcal{K}_2 \ar[d] \\
\mathcal{I}_1 \otimes_{\mathcal{O}_X} \mathcal{F}
\ar[r]^-{c_1} &
\mathcal{K}_1
}
}
\quad\text{且}\quad
\vcenter{
\xymatrix{
\mathcal{I}_3 \otimes_{\mathcal{O}_X} \mathcal{F}
\ar[r]_-{c_3} \ar[d] &
\mathcal{K}_3 \ar[d] \\
\mathcal{I}_2 \otimes_{\mathcal{O}_X} \mathcal{F}
\ar[r]^-{c_2} &
\mathcal{K}_2
}
}
$$
均交换。最后，假设给定 $\mathcal{O}_{X'_2}$-模的扩张
$$
0 \to \mathcal{K}_2 \to \mathcal{F}'_2 \to \mathcal{F} \to 0
$$
如 (\ref{defos-equation-extension}) 所述，其中
$\mathcal{K}=\mathcal{K}_2$ 且 $c_{\mathcal{F}'_2}=c_2$。在这种情形下，
应用评注 \ref{defos-remark-extension-functorial} 的函子性，可得 $X'_1$
上的扩张 $\mathcal{F}'_1$（稍后会描述这一特殊情形下的
$\mathcal{F}'_1$）。由评注 \ref{defos-remark-trivial-extension}，利用评注
\ref{defos-remark-short-exact-sequence-thickenings} 中的典范分裂
$\pi:(X'_1,\mathcal{O}_{X'_1})\to(X,\mathcal{O}_X)$，可得
$\xi_{\mathcal{F}'_1} \in
\Ext^1_{\mathcal{O}_X}(\mathcal{F}, \mathcal{K}_1)$.
最后，还有障碍
$$
o(\mathcal{F}, \mathcal{K}_3, c_3) \in
\Ext^2_{\mathcal{O}_X}(\mathcal{F}, \mathcal{K}_3)
$$
见引理 \ref{defos-lemma-inf-obs-ext}。在这种情形下，我们{\bf 断言}典范同态
$$
\partial :
\Ext^1_{\mathcal{O}_X}(\mathcal{F}, \mathcal{K}_1)
\longrightarrow
\Ext^2_{\mathcal{O}_X}(\mathcal{F}, \mathcal{K}_3)
$$
由短正合列
$0\to\mathcal{K}_3\to\mathcal{K}_2\to\mathcal{K}_1\to0$
诱导，并把 $\xi_{\mathcal{F}'_1}$ 映到障碍类
$o(\mathcal{F},\mathcal{K}_3,c_3)$。

\medskip\noindent
为证明该断言，选取嵌入 $j:\mathcal{K}_3\to\mathcal{K}$，其中
$\mathcal{K}$ 是内射 $\mathcal{O}_X$-模。可以把 $j$ 提升为同态
$j':\mathcal{K}_2\to\mathcal{K}$。令
$\mathcal{E}'_2=j'_*\mathcal{F}'_2$ 为沿 $j'$ 对
$\mathcal{F}'_2$ 所作的推出，于是
$c_{\mathcal{E}'_2}=j'\circ c_2$。图示如下：
$$
\xymatrix{
0 \ar[r] &
\mathcal{K}_2 \ar[r] \ar[d]_{j'} &
\mathcal{F}'_2 \ar[r] \ar[d] &
\mathcal{F} \ar[r] \ar[d] & 0 \\
0 \ar[r] &
\mathcal{K} \ar[r] &
\mathcal{E}'_2 \ar[r] &
\mathcal{F} \ar[r] & 0
}
$$
令 $\mathcal{E}'_3=\mathcal{E}'_2$，但经
$\mathcal{O}_{X'_3}\to\mathcal{O}_{X'_2}$ 把它视为
$\mathcal{O}_{X'_3}$-模。于是
$c_{\mathcal{E}'_3}=j\circ c_3$。引理 \ref{defos-lemma-inf-obs-ext} 的证明把
$o(\mathcal{F},\mathcal{K}_3,c_3)$ 构造为下列
$\mathcal{O}_X$-模扩张之类的边界：
$$
0 \to
\mathcal{K}/\mathcal{K}_3 \to
\mathcal{E}'_3/\mathcal{K}_3 \to
\mathcal{F} \to 0
$$
另一方面，注意 $\mathcal{F}'_1=\mathcal{F}'_2/\mathcal{K}_3$，故类
$\xi_{\mathcal{F}'_1}$ 是扩张
$$
0 \to \mathcal{K}_2/\mathcal{K}_3 \to \mathcal{F}'_2/\mathcal{K}_3
\to \mathcal{F} \to 0
$$
的类；这里利用 $\pi^\sharp$ 把该正合列视为 $\mathcal{O}_X$-模列，
而 $\pi:(X'_1,\mathcal{O}_{X'_1})\to(X,\mathcal{O}_X)$ 是典范分裂。
最后，该断言由下列交换图推出：
$$
\xymatrix{
0 \ar[r] &
\mathcal{K}_2/\mathcal{K}_3 \ar[r] \ar[d] &
\mathcal{F}'_2/\mathcal{K}_3 \ar[r] \ar[d] &
\mathcal{F} \ar[r] \ar[d] & 0 \\
0 \ar[r] &
\mathcal{K}/\mathcal{K}_3 \ar[r] &
\mathcal{E}'_3/\mathcal{K}_3 \ar[r] &
\mathcal{F} \ar[r] & 0
}
$$
该图对于上述 $\mathcal{O}_X$-模结构是 $\mathcal{O}_X$-线性的。
\end{remark}








\section{环化空间上模的无穷小形变}
\label{defos-section-deformation-modules}

\noindent
设 $i:(X,\mathcal{O}_X)\to(X',\mathcal{O}_{X'})$ 是环化空间的一阶
增厚。我们随时使用第 \ref{defos-section-thickenings-spaces} 节引入的记号。
设 $\mathcal{F}'$ 为 $\mathcal{O}_{X'}$-模，并令
$\mathcal{F}=i^*\mathcal{F}'$。在这种情形下，有
$\mathcal{O}_{X'}$-模的短正合列
$$
0 \to \mathcal{I}\mathcal{F}' \to \mathcal{F}' \to \mathcal{F} \to 0
$$
由于 $\mathcal{I}^2=0$，$\mathcal{I}\mathcal{F}'$ 上的
$\mathcal{O}_{X'}$-模结构来自唯一的 $\mathcal{O}_X$-模结构。因此，
上述正合列是 (\ref{defos-equation-extension}) 中的扩张。作为特例，如果
$\mathcal{F}'=\mathcal{O}_{X'}$，则
$i^*\mathcal{O}_{X'}=\mathcal{O}_X$ 且
$\mathcal{I}\mathcal{O}_{X'}=\mathcal{I}$，从而恢复结构层正合列
$$
0 \to \mathcal{I} \to \mathcal{O}_{X'} \to \mathcal{O}_X \to 0
$$

\begin{lemma}
\label{defos-lemma-inf-map-special}
设 $i:(X,\mathcal{O}_X)\to(X',\mathcal{O}_{X'})$ 是环化空间的一阶
增厚。设 $\mathcal{F}'$、$\mathcal{G}'$ 为
$\mathcal{O}_{X'}$-模。令
$\mathcal{F}=i^*\mathcal{F}'$、$\mathcal{G}=i^*\mathcal{G}'$，并设
$\varphi:\mathcal{F}\to\mathcal{G}$ 为 $\mathcal{O}_X$-线性同态。
$\varphi$ 到 $\mathcal{O}_{X'}$-线性同态
$\varphi':\mathcal{F}'\to\mathcal{G}'$ 的所有提升所成的集合，如果非空，
就是 $\Hom_{\mathcal{O}_X}(\mathcal{F},\mathcal{I}\mathcal{G}')$
作用下的主齐性空间。
\end{lemma}

\begin{proof}
这是引理 \ref{defos-lemma-inf-map} 的特例，不过我们也给出直接证明。我们有
模的短正合列
$$
0 \to \mathcal{I} \to \mathcal{O}_{X'} \to \mathcal{O}_X \to 0
\quad\text{且}\quad
0 \to \mathcal{I}\mathcal{G}' \to \mathcal{G}' \to \mathcal{G} \to 0
$$
$\mathcal{F}'$ 也有类似的正合列。由于 $\mathcal{I}$ 平方为零，
$\mathcal{I}$ 与 $\mathcal{I}\mathcal{G}'$ 上的
$\mathcal{O}_{X'}$-模结构均来自唯一的 $\mathcal{O}_X$-模结构。
因此
$$
\Hom_{\mathcal{O}_{X'}}(\mathcal{F}', \mathcal{I}\mathcal{G}') =
\Hom_{\mathcal{O}_X}(\mathcal{F}, \mathcal{I}\mathcal{G}')
\quad\text{且}\quad
\Hom_{\mathcal{O}_{X'}}(\mathcal{F}', \mathcal{G}) =
\Hom_{\mathcal{O}_X}(\mathcal{F}, \mathcal{G})
$$
本引理于是由正合列
$$
0 \to \Hom_{\mathcal{O}_{X'}}(\mathcal{F}', \mathcal{I}\mathcal{G}') \to
\Hom_{\mathcal{O}_{X'}}(\mathcal{F}', \mathcal{G}') \to
\Hom_{\mathcal{O}_{X'}}(\mathcal{F}', \mathcal{G})
$$
推出，见《同调代数》引理 \ref{homology-lemma-check-exactness}。
\end{proof}

\begin{lemma}
\label{defos-lemma-deform-module}
设 $(f,f')$ 是情形 \ref{defos-situation-morphism-thickenings} 中环化空间
一阶增厚的态射。设 $\mathcal{F}'$ 为 $\mathcal{O}_{X'}$-模，并令
$\mathcal{F}=i^*\mathcal{F}'$。假设 $\mathcal{F}$ 在 $S$ 上平坦，
且 $(f,f')$ 是增厚的严格态射（定义
\ref{defos-definition-strict-morphism-thickenings}）。那么下列条件等价：
\begin{enumerate}
\item $\mathcal{F}'$ 在 $S'$ 上平坦；
\item 典范同态
$f^*\mathcal{J} \otimes_{\mathcal{O}_X} \mathcal{F} \to
\mathcal{I}\mathcal{F}'$
是同构。
\end{enumerate}
此外，在这种情形下，同态
$$
f^*\mathcal{J} \otimes_{\mathcal{O}_X} \mathcal{F} \to
\mathcal{I} \otimes_{\mathcal{O}_X} \mathcal{F} \to
\mathcal{I}\mathcal{F}'
$$
均为同构。
\end{lemma}

\begin{proof}
由于 $(f,f')$ 是增厚的严格态射，同态
$f^*\mathcal{J}\to\mathcal{I}$ 是满的。因此，最后的陈述由 (2) 推出。

\medskip\noindent
证明 (1) 与 (2) 等价。可以在茎上检验这些条件。设
$x\in X\subset X'$，其像为 $s=f(x)\in S\subset S'$。令
$A'=\mathcal{O}_{S',s}$、$B'=\mathcal{O}_{X',x}$、
$A=\mathcal{O}_{S,s}$、$B=\mathcal{O}_{X,x}$。那么对某些平方零理想
$J$、$I$，有 $A=A'/J$、$B=B'/I$。由于 $(f,f')$ 是增厚的严格态射，
有 $I=JB'$。令 $M'=\mathcal{F}'_x$、$M=\mathcal{F}_x$。于是 $M'$
是 $B'$-模，$M$ 是 $B$-模。由 $\mathcal{F}=i^*\mathcal{F}'$ 可知，
满同态 $M'\to M$ 的核为 $IM'=JM'$。因而有短正合列
$$
0 \to JM' \to M' \to M \to 0
$$
利用《空间上的层》引理 \ref{sheaves-lemma-stalk-pullback-modules}
与《模层》引理 \ref{modules-lemma-stalk-tensor-product} 等同拉回与张量积
的茎，可见本引理中的典范同态在 $x$ 处的茎是同态
$$
(J \otimes_A B) \otimes_B M = J \otimes_A M = J \otimes_{A'} M'
\longrightarrow JM'
$$
假设 $\mathcal{F}$ 在 $S$ 上平坦，意味着 $M$ 是平坦 $A$-模。

\medskip\noindent
假设 (1) 成立。由《交换代数》引理
\ref{algebra-lemma-characterize-flat}，平坦性蕴含
$\text{Tor}_1^{A'}(M',A)=0$。由《交换代数》评注
\ref{algebra-remark-Tor-ring-mod-ideal}，这意味着
$J\otimes_{A'}M'\to M'$ 是单的。因此
$J\otimes_A M\to JM'$ 是同构。

\medskip\noindent
假设 (2) 成立。那么 $J\otimes_{A'}M'\to M'$ 是单的。由《交换代数》
评注 \ref{algebra-remark-Tor-ring-mod-ideal}，
$\text{Tor}_1^{A'}(M',A)=0$。再由《交换代数》引理
\ref{algebra-lemma-what-does-it-mean}，$M'$ 在 $A'$ 上平坦。
\end{proof}

\begin{lemma}
\label{defos-lemma-inf-map-rel}
设 $(f,f')$ 是情形 \ref{defos-situation-morphism-thickenings} 中一阶增厚的
态射。设 $\mathcal{F}'$、$\mathcal{G}'$ 为
$\mathcal{O}_{X'}$-模，并令
$\mathcal{F}=i^*\mathcal{F}'$、$\mathcal{G}=i^*\mathcal{G}'$。
设 $\varphi:\mathcal{F}\to\mathcal{G}$ 为 $\mathcal{O}_X$-线性同态。
假设 $\mathcal{G}'$ 在 $S'$ 上平坦，且 $(f,f')$ 是增厚的严格态射。
$\varphi$ 到 $\mathcal{O}_{X'}$-线性同态
$\varphi':\mathcal{F}'\to\mathcal{G}'$ 的所有提升所成的集合，如果非空，
就是下列群作用下的主齐性空间：
$$
\Hom_{\mathcal{O}_X}(\mathcal{F},
\mathcal{G} \otimes_{\mathcal{O}_X} f^*\mathcal{J})
$$
\end{lemma}

\begin{proof}
合并引理 \ref{defos-lemma-inf-map-special} 与 \ref{defos-lemma-deform-module} 即得。
\end{proof}

\begin{lemma}
\label{defos-lemma-inf-obs-map-special}
设 $i:(X,\mathcal{O}_X)\to(X',\mathcal{O}_{X'})$ 是环化空间的一阶
增厚。设 $\mathcal{F}'$、$\mathcal{G}'$ 为
$\mathcal{O}_{X'}$-模，并令
$\mathcal{F}=i^*\mathcal{F}'$、$\mathcal{G}=i^*\mathcal{G}'$。
设 $\varphi:\mathcal{F}\to\mathcal{G}$ 为 $\mathcal{O}_X$-线性同态。
存在元素
$$
o(\varphi) \in
\Ext^1_{\mathcal{O}_X}(Li^*\mathcal{F}',
\mathcal{I}\mathcal{G}')
$$
它为零，当且仅当 $\varphi$ 存在到 $\mathcal{O}_{X'}$-线性同态
$\varphi':\mathcal{F}'\to\mathcal{G}'$ 的提升。
\end{lemma}

\begin{proof}
由引理 \ref{defos-lemma-inf-map-special} 的证明可见，$\varphi$ 经同态
$$
\Hom_{\mathcal{O}_X}(\mathcal{F}, \mathcal{G}) =
\Hom_{\mathcal{O}_{X'}}(\mathcal{F}', \mathcal{G}) \longrightarrow
\Ext^1_{\mathcal{O}_{X'}}(\mathcal{F}', \mathcal{I}\mathcal{G}')
$$
所得的边界为零，当且仅当提升存在。又有
$$
\Ext^1_{\mathcal{O}_{X'}}(\mathcal{F}', \mathcal{I}\mathcal{G}') =
\Ext^1_{\mathcal{O}_X}(Li^*\mathcal{F}', \mathcal{I}\mathcal{G}')
$$
这是导出范畴上 $i_*=Ri_*$ 与 $Li^*$ 的伴随性（《层的上同调》引理
\ref{cohomology-lemma-adjoint}）。因此结论成立。
\end{proof}

\begin{lemma}
\label{defos-lemma-inf-obs-map-rel}
设 $(f,f')$ 是情形 \ref{defos-situation-morphism-thickenings} 中一阶增厚的
态射。设 $\mathcal{F}'$、$\mathcal{G}'$ 为
$\mathcal{O}_{X'}$-模，并令
$\mathcal{F}=i^*\mathcal{F}'$、$\mathcal{G}=i^*\mathcal{G}'$。
设 $\varphi:\mathcal{F}\to\mathcal{G}$ 为 $\mathcal{O}_X$-线性同态。
假设 $\mathcal{F}'$ 与 $\mathcal{G}'$ 均在 $S'$ 上平坦，且
$(f,f')$ 是增厚的严格态射。存在元素
$$
o(\varphi) \in  \Ext^1_{\mathcal{O}_X}(\mathcal{F},
\mathcal{G} \otimes_{\mathcal{O}_X} f^*\mathcal{J})
$$
它为零，当且仅当 $\varphi$ 存在到 $\mathcal{O}_{X'}$-线性同态
$\varphi':\mathcal{F}'\to\mathcal{G}'$ 的提升。
\end{lemma}

\begin{proof}[第一种证明]
这由引理 \ref{defos-lemma-inf-obs-map-special} 推出，因为我们断言在本引理的
假设下有
$$
\Ext^1_{\mathcal{O}_X}(Li^*\mathcal{F}',
\mathcal{I}\mathcal{G}') =
\Ext^1_{\mathcal{O}_X}(\mathcal{F},
\mathcal{G} \otimes_{\mathcal{O}_X} f^*\mathcal{J})
$$
事实上，由引理 \ref{defos-lemma-deform-module} 有
$\mathcal{I}\mathcal{G}' =
\mathcal{G} \otimes_{\mathcal{O}_X} f^*\mathcal{J}$
。另一方面，注意
$$
H^{-1}(Li^*\mathcal{F}') =
\text{Tor}_1^{\mathcal{O}_{X'}}(\mathcal{F}', \mathcal{O}_X)
$$
（局部计算从略）。利用短正合列
$$
0 \to \mathcal{I} \to \mathcal{O}_{X'} \to \mathcal{O}_X \to 0
$$
可见，该 $\text{Tor}_1$ 由同态
$\mathcal{I}\otimes_{\mathcal{O}_X}\mathcal{F}
\to\mathcal{I}\mathcal{F}'$ 的核计算；由引理
\ref{defos-lemma-deform-module} 的最后一个断言，该核为零。因此
$\tau_{\geq-1}Li^*\mathcal{F}'=\mathcal{F}$。另一方面，有
$$
\Ext^1_{\mathcal{O}_X}(Li^*\mathcal{F}',
\mathcal{I}\mathcal{G}') =
\Ext^1_{\mathcal{O}_X}(\tau_{\geq -1}Li^*\mathcal{F}',
\mathcal{I}\mathcal{G}')
$$
这是《导出范畴》引理 \ref{derived-lemma-negative-vanishing} 的对偶。
\end{proof}

\begin{proof}[第二种证明]
可以如下应用引理 \ref{defos-lemma-inf-obs-map}。注意，由引理
\ref{defos-lemma-deform-module} 有
$\mathcal{K}=\mathcal{I}\otimes_{\mathcal{O}_X}\mathcal{F}$、
$\mathcal{L}=\mathcal{I}\otimes_{\mathcal{O}_X}\mathcal{G}$
，并且 $c_{\mathcal{F}'}=1\otimes1$、
$c_{\mathcal{G}'}=1\otimes1$。取 $\psi=1\otimes\varphi$，引理中的
图表交换。因此可以取
$o(\varphi)=o(\varphi,1\otimes\varphi)$。
\end{proof}

\begin{lemma}
\label{defos-lemma-inf-ext-rel}
设 $(f,f')$ 是情形 \ref{defos-situation-morphism-thickenings} 中一阶增厚的
态射，$\mathcal{F}$ 为 $\mathcal{O}_X$-模。假设 $(f,f')$ 是增厚的
严格态射，且 $\mathcal{F}$ 在 $S$ 上平坦。如果存在二元组
$(\mathcal{F}',\alpha)$，其中 $\mathcal{F}'$ 是在 $S'$ 上平坦的
$\mathcal{O}_{X'}$-模，而
$\alpha:i^*\mathcal{F}'\to\mathcal{F}$ 是同构，那么这类二元组的
同构类集合是下列群作用下的主齐性空间：
$\Ext^1_{\mathcal{O}_X}(
\mathcal{F}, \mathcal{I} \otimes_{\mathcal{O}_X} \mathcal{F})$.
\end{lemma}

\begin{proof}
假设存在一个这样的模。那么由引理 \ref{defos-lemma-deform-module}，典范同态
$$
f^*\mathcal{J} \otimes_{\mathcal{O}_X} \mathcal{F} \to
\mathcal{I} \otimes_{\mathcal{O}_X} \mathcal{F}
$$
是同构。应用引理 \ref{defos-lemma-inf-ext}，取 $\mathcal{K}=
\mathcal{I} \otimes_{\mathcal{O}_X} \mathcal{F}$
、$c=1$。由引理 \ref{defos-lemma-deform-module}，相应的扩张
$\mathcal{F}'$ 均在 $S'$ 上平坦。
\end{proof}

\begin{lemma}
\label{defos-lemma-inf-obs-ext-rel}
设 $(f,f')$ 是情形 \ref{defos-situation-morphism-thickenings} 中一阶增厚的
态射，$\mathcal{F}$ 为 $\mathcal{O}_X$-模。假设 $(f,f')$ 是增厚的
严格态射，且 $\mathcal{F}$ 在 $S$ 上平坦。那么，存在在 $S'$ 上平坦的
$\mathcal{O}_{X'}$-模 $\mathcal{F}'$，满足
$i^*\mathcal{F}'\cong\mathcal{F}$，当且仅当：
\begin{enumerate}
\item 典范同态
$f^*\mathcal{J} \otimes_{\mathcal{O}_X} \mathcal{F} \to
\mathcal{I} \otimes_{\mathcal{O}_X} \mathcal{F}$
是同构；
\item 引理 \ref{defos-lemma-inf-obs-ext} 中的类
$o(\mathcal{F}, \mathcal{I} \otimes_{\mathcal{O}_X} \mathcal{F}, 1)
\in \Ext^2_{\mathcal{O}_X}(
\mathcal{F}, \mathcal{I} \otimes_{\mathcal{O}_X} \mathcal{F})$
为零。
\end{enumerate}
\end{lemma}

\begin{proof}
这立即由引理 \ref{defos-lemma-deform-module} 对在 $S'$ 上平坦的
$\mathcal{O}_{X'}$-模的刻画以及引理 \ref{defos-lemma-inf-obs-ext} 推出。
\end{proof}






\section{在环化空间平坦增厚上的平坦模中的应用}
\label{defos-section-flat}

\noindent
考虑交换图
$$
\xymatrix{
(X, \mathcal{O}_X) \ar[r]_i \ar[d]_f & (X', \mathcal{O}_{X'}) \ar[d]^{f'} \\
(S, \mathcal{O}_S) \ar[r]^t & (S', \mathcal{O}_{S'})
}
$$
其中各项为环化空间，水平箭头是情形
\ref{defos-situation-morphism-thickenings} 中的一阶增厚。令
$\mathcal{I}=\Ker(i^\sharp)\subset\mathcal{O}_{X'}$、
$\mathcal{J}=\Ker(t^\sharp)\subset\mathcal{O}_{S'}$。设
$\mathcal{F}$ 为 $\mathcal{O}_X$-模。假设：
\begin{enumerate}
\item $(f,f')$ 是增厚的严格态射；
\item $f'$ 平坦；
\item $\mathcal{F}$ 在 $S$ 上平坦。
\end{enumerate}
注意，(1) $+$ (2) 蕴含 $\mathcal{I}=f^*\mathcal{J}$（把引理
\ref{defos-lemma-deform-module} 应用于 $\mathcal{O}_{X'}$）。在这些假设下，
上一节的理论具有特别简洁的形式。下述引理汇总已经得到的结果。

\begin{lemma}
\label{defos-lemma-flat}
在上述情形中：
\begin{enumerate}
\item 存在在 $S'$ 上平坦的 $\mathcal{O}_{X'}$-模 $\mathcal{F}'$，
满足 $i^*\mathcal{F}'\cong\mathcal{F}$，当且仅当引理
\ref{defos-lemma-inf-obs-ext} 中的类
$o(\mathcal{F}, f^*\mathcal{J} \otimes_{\mathcal{O}_X} \mathcal{F}, 1)
\in \Ext^2_{\mathcal{O}_X}(
\mathcal{F}, f^*\mathcal{J} \otimes_{\mathcal{O}_X} \mathcal{F})$
为零。
\item 如果这样的模存在，那么提升的同构类集合是
$\Ext^1_{\mathcal{O}_X}(
\mathcal{F}, f^*\mathcal{J} \otimes_{\mathcal{O}_X} \mathcal{F})$
作用下的主齐性空间。
\item 给定提升 $\mathcal{F}'$，其拉回为
$\text{id}_{\mathcal{F}}$ 的 $\mathcal{F}'$ 自同构所成的集合，典范同构于
$\Ext^0_{\mathcal{O}_X}(
\mathcal{F}, f^*\mathcal{J} \otimes_{\mathcal{O}_X} \mathcal{F})$.
\end{enumerate}
\end{lemma}

\begin{proof}
由于上面已说明 $\mathcal{I}=f^*\mathcal{J}$，第 (1) 点由引理
\ref{defos-lemma-inf-obs-ext-rel} 推出。第 (2) 点由引理
\ref{defos-lemma-inf-ext-rel} 推出，第 (3) 点由引理
\ref{defos-lemma-inf-map-rel} 推出。
\end{proof}

\begin{situation}
\label{defos-situation-ses-flat-thickenings}
设 $f:(X,\mathcal{O}_X)\to(S,\mathcal{O}_S)$ 为环化空间的态射。
考虑交换图
$$
\xymatrix{
(X'_1, \mathcal{O}'_1) \ar[r]_h \ar[d]_{f'_1} &
(X'_2, \mathcal{O}'_2) \ar[r] \ar[d]_{f'_2} &
(X'_3, \mathcal{O}'_3) \ar[d]_{f'_3} \\
(S'_1, \mathcal{O}_{S'_1}) \ar[r] &
(S'_2, \mathcal{O}_{S'_2}) \ar[r] &
(S'_3, \mathcal{O}_{S'_3})
}
$$
其中：(a) 上行是 $X$ 的一阶增厚的短正合列；(b) 下行是 $S$ 的一阶
增厚的短正合列；(c) 每个 $f'_i$ 限制为 $f$；(d) 每个二元组
$(f,f'_i)$ 都是增厚的严格态射；(e) 每个 $f'_i$ 都平坦。最后，设
$\mathcal{F}'_2$ 为在 $S'_2$ 上平坦的 $\mathcal{O}'_2$-模，并令
$\mathcal{F}=\mathcal{F}'_2|_X$。设 $\pi:X'_1\to X$ 为典范分裂
（评注 \ref{defos-remark-short-exact-sequence-thickenings}）。
\end{situation}

\begin{lemma}
\label{defos-lemma-verify-iv}
在情形 \ref{defos-situation-ses-flat-thickenings} 中，模
$\pi^*\mathcal{F}$ 与 $h^*\mathcal{F}'_2$ 都是在 $S'_1$ 上平坦的
$\mathcal{O}'_1$-模，且在 $X$ 上限制为 $\mathcal{F}$。二者之差
（引理 \ref{defos-lemma-flat}）是
$\Ext^1_{\mathcal{O}_X}(
\mathcal{F}, f^*\mathcal{J}_1 \otimes_{\mathcal{O}_X} \mathcal{F})$
中的元素 $\theta$，它在
$\Ext^2_{\mathcal{O}_X}(
\mathcal{F}, f^*\mathcal{J}_3 \otimes_{\mathcal{O}_X} \mathcal{F})$
中的边界，等于把 $\mathcal{F}$ 提升为在 $S'_3$ 上平坦的
$\mathcal{O}'_3$-模的障碍（引理 \ref{defos-lemma-flat}）。
\end{lemma}

\begin{proof}
注意，$\pi^*\mathcal{F}$ 与 $h^*\mathcal{F}'_2$ 在 $X$ 上都限制为
$\mathcal{F}$，且同态 $\pi^*\mathcal{F}\to\mathcal{F}$ 与
$h^*\mathcal{F}'_2\to\mathcal{F}$ 的核均为
$f^*\mathcal{J}_1\otimes_{\mathcal{O}_X}\mathcal{F}$。因此，由引理
\ref{defos-lemma-deform-module} 得到平坦性。取边界是有意义的，因为模列
$$
0 \to f^*\mathcal{J}_3 \otimes_{\mathcal{O}_X} \mathcal{F} \to
f^*\mathcal{J}_2 \otimes_{\mathcal{O}_X} \mathcal{F} \to
f^*\mathcal{J}_1 \otimes_{\mathcal{O}_X} \mathcal{F} \to 0
$$
是短正合的；这是情形 \ref{defos-situation-ses-flat-thickenings} 中的假设以及
$\mathcal{F}$ 在 $S$ 上平坦的结果。关于障碍类的陈述，就是把评注
\ref{defos-remark-complex-thickenings-and-ses-modules} 的结果直接应用于
这一特殊情形。
\end{proof}






\section{环化空间的形变与朴素余切复形}
\label{defos-section-deformations-ringed-spaces}

\noindent
本节利用朴素余切复形来处理一些形变理论问题。先取环化空间的一阶增厚
$t:(S,\mathcal{O}_S)\to(S',\mathcal{O}_{S'})$。记
$\mathcal{J}=\Ker(t^\sharp)$，并等同 $S$ 与 $S'$ 的底层拓扑空间。
此外，设已给定环化空间的态射
$f:(X,\mathcal{O}_X)\to(S,\mathcal{O}_S)$、$\mathcal{O}_X$-模
$\mathcal{G}$ 以及模层的 $f$-映射 $c:\mathcal{J}\to\mathcal{G}$
（《空间上的层》定义 \ref{sheaves-definition-f-map} 及第
\ref{sheaves-section-ringed-spaces-functoriality-modules} 节）。
我们要问，能否找到一个对象来替代下图中的问号，使图表交换：
\begin{equation}
\label{defos-equation-to-solve-ringed-spaces}
\vcenter{
\xymatrix{
0 \ar[r] & \mathcal{G} \ar[r] & {?} \ar[r] & \mathcal{O}_X \ar[r] & 0 \\
0 \ar[r] & \mathcal{J} \ar[u]^c \ar[r] & \mathcal{O}_{S'} \ar[u] \ar[r] &
\mathcal{O}_S \ar[u] \ar[r] & 0
}
}
\end{equation}
（其中竖直箭头为 $f$-映射），并且，如果解存在，它在何种意义下是唯一的。
更确切地说，我们寻找一个一阶增厚
$i:(X,\mathcal{O}_X)\to(X',\mathcal{O}_{X'})$ 以及
(\ref{defos-equation-morphism-thickenings}) 中的增厚态射 $(f,f')$，其中
$\Ker(i^\sharp)$ 与 $\mathcal{G}$ 等同，且 $(f')^\sharp$ 诱导给定映射
$c$。我们称 $X'$ 为 (\ref{defos-equation-to-solve-ringed-spaces}) 的一个
{\it 解}。

\begin{lemma}
\label{defos-lemma-huge-diagram-ringed-spaces}
假设给定环化空间态射的交换图
\begin{equation}
\label{defos-equation-huge-1}
\vcenter{
\xymatrix{
& (X_2, \mathcal{O}_{X_2}) \ar[r]_{i_2} \ar[d]_{f_2} \ar[ddl]_g &
(X'_2, \mathcal{O}_{X'_2}) \ar[d]^{f'_2} \\
& (S_2, \mathcal{O}_{S_2}) \ar[r]^{t_2} \ar[ddl]|\hole &
(S'_2, \mathcal{O}_{S'_2}) \ar[ddl] \\
(X_1, \mathcal{O}_{X_1}) \ar[r]_{i_1} \ar[d]_{f_1} &
(X'_1, \mathcal{O}_{X'_1}) \ar[d]^{f'_1} \\
(S_1, \mathcal{O}_{S_1}) \ar[r]^{t_1} &
(S'_1, \mathcal{O}_{S'_1})
}
}
\end{equation}
其中水平箭头为一阶增厚。令
$\mathcal{G}_j=\Ker(i_j^\sharp)$，并假设给定模的 $g$-映射
$\nu:\mathcal{G}_1\to\mathcal{G}_2$，从而得到交换图
\begin{equation}
\label{defos-equation-huge-2}
\vcenter{
\xymatrix{
& 0 \ar[r] & \mathcal{G}_2 \ar[r] &
\mathcal{O}_{X'_2} \ar[r] &
\mathcal{O}_{X_2} \ar[r] & 0 \\
& 0 \ar[r]|\hole &
\mathcal{J}_2 \ar[u]_{c_2} \ar[r] &
\mathcal{O}_{S'_2} \ar[u] \ar[r]|\hole &
\mathcal{O}_{S_2} \ar[u] \ar[r] & 0 \\
0 \ar[r] & \mathcal{G}_1 \ar[ruu] \ar[r] &
\mathcal{O}_{X'_1} \ar[r] &
\mathcal{O}_{X_1} \ar[ruu] \ar[r] & 0 \\
0 \ar[r] & \mathcal{J}_1 \ar[ruu]|\hole \ar[u]^{c_1} \ar[r] &
\mathcal{O}_{S'_1} \ar[ruu]|\hole \ar[u] \ar[r] &
\mathcal{O}_{S_1} \ar[ruu]|\hole \ar[u] \ar[r] & 0
}
}
\end{equation}
其中前后两面均为 (\ref{defos-equation-to-solve-ringed-spaces}) 的解。
\begin{enumerate}
\item 在
$\Ext^1_{\mathcal{O}_{X_2}}(Lg^*\NL_{X_1/S_1}, \mathcal{G}_2)$
中存在一个典范元；该元为零，当且仅当存在与 $\nu$ 相容并可置入
(\ref{defos-equation-huge-1}) 的环化空间态射 $X'_2\to X'_1$。
\item 如果存在与 $\nu$ 相容并可置入 (\ref{defos-equation-huge-1}) 的态射
$X'_2\to X'_1$，那么所有这类态射组成的集合是下列群作用下的
主齐性空间：
$$
\Hom_{\mathcal{O}_{X_1}}(\Omega_{X_1/S_1}, g_*\mathcal{G}_2) =
\Hom_{\mathcal{O}_{X_2}}(g^*\Omega_{X_1/S_1}, \mathcal{G}_2) =
\Ext^0_{\mathcal{O}_{X_2}}(Lg^*\NL_{X_1/S_1}, \mathcal{G}_2).
$$
\end{enumerate}
\end{lemma}

\begin{proof}
朴素余切复形 $\NL_{X_1/S_1}$ 定义于《模层》定义
\ref{modules-definition-cotangent-complex-morphism-ringed-topoi}。
引理最后一个陈述中的等号来自以下事实：$g^*$ 左伴随于 $g_*$；
$H^0(\NL_{X_1/S_1})=\Omega_{X_1/S_1}$（由朴素余切复形的构造）；
$Lg^*$ 是 $g^*$ 的左导出函子。因此，在余下的证明中，我们使用群
$\Ext^k_{\mathcal{O}_{X_2}}(Lg^*\NL_{X_1/S_1}, \mathcal{G}_2)$,
其中 $k=0,1$。先说明，可以约化到本引理中所有环化空间的底层拓扑空间
都相同的情形。

\medskip\noindent
为此，注意 $g^{-1}\NL_{X_1/S_1}$ 等于环层同态
$g^{-1}f_1^{-1}\mathcal{O}_{S_1}\to g^{-1}\mathcal{O}_{X_1}$ 的
朴素余切复形，见《模层》引理 \ref{modules-lemma-pullback-NL}。
此外，$\NL_{X_1/S_1}$ 的 $0$ 次项是平坦
$\mathcal{O}_{X_1}$-模，因而典范同态
$$
Lg^*\NL_{X_1/S_1}
\longrightarrow
g^{-1}\NL_{X_1/S_1} \otimes_{g^{-1}\mathcal{O}_{X_1}} \mathcal{O}_{X_2}
$$
在 $0$ 次和 $-1$ 次上同调层上诱导同构。因此，可以把引理中的 Ext 群
替换为
$$
\Ext^k_{g^{-1}\mathcal{O}_{X_1}}(g^{-1}\NL_{X_1/S_1}, \mathcal{G}_2) =
\Ext^k_{g^{-1}\mathcal{O}_{X_1}}(
\NL_{g^{-1}\mathcal{O}_{X_1}/g^{-1}f_1^{-1}\mathcal{O}_{S_1}}, \mathcal{G}_2)
$$
与 $\nu$ 相容并可置入 (\ref{defos-equation-huge-1}) 的环化空间态射
$X'_2\to X'_1$ 所成的集合，与下列同态所成的集合一一对应：与
$f^\sharp$ 和 $\nu$ 相容的
$g^{-1}f_1^{-1}\mathcal{O}_{S'_1}$-代数同态
$g^{-1}\mathcal{O}_{X'_1}\to\mathcal{O}_{X'_2}$。由此可见，可以假设
我们有 $X$ 上层的图表 (\ref{defos-equation-huge-2})，并寻找可置入该图的
环层同态 $\mathcal{O}_{X'_1}\to\mathcal{O}_{X'_2}$。

\medskip\noindent
在本引理余下的证明中，假设所有底层拓扑空间都相同。也就是说，我们有
空间 $X$ 上层的图表 (\ref{defos-equation-huge-2})，并寻找可置入其中的环层
同态 $\mathcal{O}_{X'_1}\to\mathcal{O}_{X'_2}$。所用的 Ext 群为
$\Ext^k_{\mathcal{O}_{X_1}}(
\NL_{\mathcal{O}_{X_1}/\mathcal{O}_{S_1}}, \mathcal{G}_2)$, $k = 0, 1$.

\medskip\noindent
第 1 步：构造障碍类。考虑集合层
$$
\mathcal{E} = \mathcal{O}_{X'_1} \times_{\mathcal{O}_{X_2}} \mathcal{O}_{X'_2}
$$
它带有满映射 $\alpha:\mathcal{E}\to\mathcal{O}_{X_1}$，故可用
$\NL(\alpha)$ 代替
$\NL_{\mathcal{O}_{X_1}/\mathcal{O}_{S_1}}$，见《模层》引理
\ref{modules-lemma-NL-up-to-qis}。令
$$
\mathcal{I}' =
\Ker(\mathcal{O}_{S'_1}[\mathcal{E}] \to \mathcal{O}_{X_1})
\quad\text{且}\quad
\mathcal{I} =
\Ker(\mathcal{O}_{S_1}[\mathcal{E}] \to \mathcal{O}_{X_1})
$$
存在满同态 $\mathcal{I}'\to\mathcal{I}$，其核为
$\mathcal{J}_1\mathcal{O}_{S'_1}[\mathcal{E}]$。我们得到两个
$\mathcal{O}_{S'_1}$-代数同态
$$
a : \mathcal{O}_{S'_1}[\mathcal{E}] \to \mathcal{O}_{X'_1}
\quad\text{且}\quad
b : \mathcal{O}_{S'_1}[\mathcal{E}] \to \mathcal{O}_{X'_2}
$$
它们分别诱导同态
$a|_{\mathcal{I}'}:\mathcal{I}'\to\mathcal{G}_1$ 与
$b|_{\mathcal{I}'}:\mathcal{I}'\to\mathcal{G}_2$。$a$ 与 $b$ 均在
$(\mathcal{I}')^2$ 上为零。此外，由于 (\ref{defos-equation-huge-2}) 的左方块
交换，把 $a$ 与 $b$ 视为取值于 $\mathcal{G}_2$ 的同态时，它们在
$\mathcal{J}_1\mathcal{O}_{S'_1}[\mathcal{E}]$ 上一致。因此，差
$b|_{\mathcal{I}'} - \nu \circ a|_{\mathcal{I}'}$
诱导良定义的 $\mathcal{O}_{X_1}$-线性同态
$$
\xi : \mathcal{I}/\mathcal{I}^2 \longrightarrow \mathcal{G}_2
$$
它把 $\mathcal{I}$ 的局部截面 $f$ 的类映到
$\nu(a(f'))-b(f')$，其中 $f'$ 是 $f$ 到 $\mathcal{I}'$ 的局部截面的
一个提升。令
% P11-E0002: frozen source sign convention retained; see control/ERRATA_CANDIDATES.jsonl.
$[\xi] \in \Ext^1_{\mathcal{O}_{X_1}}(\NL(\alpha), \mathcal{G}_2)$
为其像（见下文）。

\medskip\noindent
第 2 步：$[\xi]$ 为零是必要条件。记
$\Omega = \Omega_{\mathcal{O}_{S_1}[\mathcal{E}]/\mathcal{O}_{S_1}}
\otimes_{\mathcal{O}_{S_1}[\mathcal{E}]} \mathcal{O}_{X_1}$.
注意，$\NL(\alpha)=(\mathcal{I}/\mathcal{I}^2\to\Omega)$ 可置于可区别三角
$$
\Omega[0] \to
\NL(\alpha) \to
\mathcal{I}/\mathcal{I}^2[1] \to
\Omega[1]
$$
因此可见，$[\xi]$ 为零，当且仅当对某个同态
$\Omega\to\mathcal{G}_2$，$\xi$ 是复合
$\mathcal{I}/\mathcal{I}^2\to\Omega\to\mathcal{G}_2$。假设存在可置入
(\ref{defos-equation-huge-2}) 的环层同态
$\varphi:\mathcal{O}_{X'_1}\to\mathcal{O}_{X'_2}$。此时考虑同态
$\mathcal{O}_{S'_1}[\mathcal{E}] \to \mathcal{G}_2$,
$f'\mapsto b(f')-\varphi(a(f'))$。计算表明，该同态在
$\mathcal{J}_1\mathcal{O}_{S'_1}[\mathcal{E}]$ 上为零，并诱导导子
$\mathcal{O}_{S_1}[\mathcal{E}]\to\mathcal{G}_2$。所得的线性同态
$\Omega\to\mathcal{G}_2$ 说明此时 $[\xi]=0$。

\medskip\noindent
第 3 步：$[\xi]$ 为零是充分条件。设
$\theta:\Omega\to\mathcal{G}_2$ 为 $\mathcal{O}_{X_1}$-线性同态，
使 $\xi=\theta\circ(\mathcal{I}/\mathcal{I}^2\to\Omega)$。
计算表明
$$
b + \theta \circ d : \mathcal{O}_{S'_1}[\mathcal{E}] \to \mathcal{O}_{X'_2}
$$
在 $\mathcal{I}'$ 上的限制与
$\nu\circ a:\mathcal{I}'\to\mathcal{G}_2$ 一致。由于
$\mathcal{O}_{X'_1}$ 是 $\mathcal{I}'\to
\mathcal{O}_{S'_1}[\mathcal{E}]$ 与
$\mathcal{I}'\to\mathcal{G}_1$ 的推出，同态
$b+\theta\circ d$ 与 $a$ 定义一个可置入
(\ref{defos-equation-huge-2}) 的同态
$\mathcal{O}_{X'_1}\to\mathcal{O}_{X'_2}$。

\medskip\noindent
上述特殊情形下第 (2) 点的证明从略。提示：这与引理
\ref{defos-lemma-huge-diagram} 第 (2) 点的证明完全相同。
\end{proof}

\begin{lemma}
\label{defos-lemma-NL-represent-ext-class}
设 $X$ 为拓扑空间，$\mathcal{A}\to\mathcal{B}$ 为环层同态，
$\mathcal{G}$ 为 $\mathcal{B}$-模。设
$\xi\in\Ext^1_\mathcal{B}(\NL_{\mathcal{B}/\mathcal{A}},\mathcal{G})$。
存在集合层的映射 $\alpha:\mathcal{E}\to\mathcal{B}$，使得把
$\xi\in\Ext^1_\mathcal{B}(\NL(\alpha),\mathcal{G})$ 表示为同态
$\mathcal{I}/\mathcal{I}^2\to\mathcal{G}$ 的类（记号见证明）。
\end{lemma}

\begin{proof}
回顾：给定 $\alpha:\mathcal{E}\to\mathcal{B}$，若
$\mathcal{A}[\mathcal{E}]\to\mathcal{B}$ 是满的且核为
$\mathcal{I}$，则复形
$\NL(\alpha) = (\mathcal{I}/\mathcal{I}^2 \to 
\Omega_{\mathcal{A}[\mathcal{E}]/\mathcal{A}}
\otimes_{\mathcal{A}[\mathcal{E}]} \mathcal{B})$
典范同构于 $\NL_{\mathcal{B}/\mathcal{A}}$，见《模层》引理
\ref{modules-lemma-NL-up-to-qis}。此外，注意
$\Omega = \Omega_{\mathcal{A}[\mathcal{E}]/\mathcal{A}}
\otimes_{\mathcal{A}[\mathcal{E}]} \mathcal{B}$
是与预层
$U \mapsto \bigoplus_{e \in \mathcal{E}(U)} \mathcal{B}(U)$.
相伴的层。换言之，$\Omega$ 是由集合层 $\mathcal{E}$ 生成的自由
$\mathcal{B}$-模；特别地，存在典范映射 $\mathcal{E}\to\Omega$。

\medskip\noindent
据此，选取某个 $\mathcal{E}$（例如像朴素余切复形定义中那样取
$\mathcal{E}=\mathcal{B}$）。把 $\xi$ 写成某个同态
$\mathcal{I}/\mathcal{I}^2\to\mathcal{G}$ 的类之障碍，是
$\Ext^1_\mathcal{B}(\Omega,\mathcal{G})$ 中的元素。设该元素由
$\mathcal{B}$-模扩张
$0\to\mathcal{G}\to\mathcal{H}\to\Omega\to0$ 表示。考虑集合层
$\mathcal{E}' = \mathcal{E} \times_\Omega \mathcal{H}$
它带有诱导映射 $\alpha':\mathcal{E}'\to\mathcal{B}$。令
$\mathcal{I}'=\Ker(\mathcal{A}[\mathcal{E}']\to\mathcal{B})$，
并令 $\Omega'=\Omega_{\mathcal{A}[\mathcal{E}']/\mathcal{A}}
\otimes_{\mathcal{A}[\mathcal{E}']} \mathcal{B}$.
在拟同构 $\NL(\alpha')\to\NL(\alpha)$ 下拉回 $\xi$，其在
$\Ext^1_\mathcal{B}(\Omega',\mathcal{G})$ 中的像为零。事实上，沿
$\Omega'\to\Omega$ 拉回扩张 $\mathcal{H}$ 所得的扩张是分裂的：
$\Omega'$ 是由集合层 $\mathcal{E}'$ 生成的自由 $\mathcal{B}$-模，
且按构造有交换图
$$
\xymatrix{
\mathcal{E}' \ar[r] \ar[d] & \mathcal{E} \ar[d] \\
\mathcal{H} \ar[r] & \Omega
}
$$
证明完成。
\end{proof}

\begin{lemma}
\label{defos-lemma-choices-ringed-spaces}
如果 (\ref{defos-equation-to-solve-ringed-spaces}) 存在解，那么解的同构类
集合是 $\Ext^1_{\mathcal{O}_X}(\NL_{X/S},\mathcal{G})$ 作用下的
主齐性空间。
\end{lemma}

\begin{proof}
先注意，给定 (\ref{defos-equation-to-solve-ringed-spaces}) 的两个解
$X'_1$ 与 $X'_2$，由引理 \ref{defos-lemma-huge-diagram-ringed-spaces}
可得障碍元
$o(X'_1,X'_2)\in\Ext^1_{\mathcal{O}_X}(\NL_{X/S},\mathcal{G})$；
它阻碍态射 $X'_1\to X'_2$ 的存在。显然，该元也是同构存在性的障碍，
因而能够区分不同的同构类。因此，为完成证明，只需说明：给定一个解
$X'$ 和元素
$\xi\in\Ext^1_{\mathcal{O}_X}(\NL_{X/S},\mathcal{G})$，可以找到
另一个解 $X'_\xi$，使得 $o(X',X'_\xi)=\xi$。

\medskip\noindent
按照引理 \ref{defos-lemma-NL-represent-ext-class}，为类 $\xi$ 选取
$\alpha:\mathcal{E}\to\mathcal{O}_X$。考虑满同态
$f^{-1}\mathcal{O}_S[\mathcal{E}]\to\mathcal{O}_X$，记其核为
$\mathcal{I}$；相应的朴素余切复形为
$\NL(\alpha) = (\mathcal{I}/\mathcal{I}^2 \to
\Omega_{f^{-1}\mathcal{O}_S[\mathcal{E}]/f^{-1}\mathcal{O}_S}
\otimes_{f^{-1}\mathcal{O}_S[\mathcal{E}]} \mathcal{O}_X)$.
由该引理，$\xi$ 是某个同态
$\delta:\mathcal{I}/\mathcal{I}^2\to\mathcal{G}$ 的类。把
$\mathcal{E}$ 替换为
$\mathcal{E}\times_{\mathcal{O}_X}\mathcal{O}_{X'}$ 后，还可假设
$\alpha$ 经映射 $\alpha':\mathcal{E}\to\mathcal{O}_{X'}$ 分解。

\medskip\noindent
这些选择确定 $f^{-1}\mathcal{O}_{S'}$-代数同态
$\varphi:f^{-1}\mathcal{O}_{S'}[\mathcal{E}]\to\mathcal{O}_{X'}$。
令
$\mathcal{I}'=\Ker(f^{-1}\mathcal{O}_{S'}[\mathcal{E}]
\to\mathcal{O}_X)$。注意，$\varphi$ 诱导同态
$\varphi|_{\mathcal{I}'}:\mathcal{I}'\to\mathcal{G}$，且
$\mathcal{O}_{X'}$ 是下图所示的推出：
$$
\xymatrix{
0 \ar[r] & \mathcal{G} \ar[r] & \mathcal{O}_{X'} \ar[r] &
\mathcal{O}_X \ar[r] & 0 \\
0 \ar[r] & \mathcal{I}' \ar[u]^{\varphi|_{\mathcal{I}'}} \ar[r] &
f^{-1}\mathcal{O}_{S'}[\mathcal{E}] \ar[u] \ar[r] &
\mathcal{O}_X \ar[u]_{=} \ar[r] & 0
}
$$
令 $\psi:\mathcal{I}'\to\mathcal{G}$ 为
$\varphi|_{\mathcal{I}'}$ 与下列复合同态之和：
$$
\mathcal{I}' \to \mathcal{I}'/(\mathcal{I}')^2 \to
\mathcal{I}/\mathcal{I}^2 \xrightarrow{\delta} \mathcal{G}.
$$
于是，沿 $\psi$ 所作的推出给出另一个环扩张
$\mathcal{O}_{X'_\xi}$，它可置于上述形式的图表中。计算（从略）表明
$o(X',X'_\xi)=\xi$，正如所需。
\end{proof}

\begin{lemma}
\label{defos-lemma-extensions-of-relative-ringed-spaces}
设 $f:(X,\mathcal{O}_X)\to(S,\mathcal{O}_S)$ 为环化空间的态射，
$\mathcal{G}$ 为 $\mathcal{O}_X$-模。所有如下
$f^{-1}\mathcal{O}_S$-代数扩张的同构类所成的集合
$$
0 \to \mathcal{G} \to \mathcal{O}_{X'} \to \mathcal{O}_X \to 0
$$
（其中 $\mathcal{G}$ 是平方零理想）\footnote{换言之，即 $S$ 上的一阶
增厚 $i:X\to X'$ 连同 $\mathcal{O}_X$-模同构
$\mathcal{G}\to\Ker(i^\sharp)$ 的同构类集合。}
与 $\Ext^1_{\mathcal{O}_X}(\NL_{X/S},\mathcal{G})$ 典范双射。
\end{lemma}

\begin{proof}
为证明这一点，把前述结果应用于
(\ref{defos-equation-to-solve-ringed-spaces}) 由下图给出的情形：
$$
\xymatrix{
0 \ar[r] &
\mathcal{G} \ar[r] &
{?} \ar[r] &
\mathcal{O}_X \ar[r] & 0 \\
0 \ar[r] &
0 \ar[u] \ar[r] &
\mathcal{O}_S \ar[u] \ar[r]^{\text{id}} &
\mathcal{O}_S \ar[u] \ar[r] & 0
}
$$
于是，本引理由引理 \ref{defos-lemma-choices-ringed-spaces} 以及解的存在性
推出；一个解就是 $\mathcal{G}\oplus\mathcal{O}_X$。（该双射的直接
构造见下述评注。）
\end{proof}

\begin{remark}
\label{defos-remark-extensions-of-relative-ringed-spaces}
设 $f:(X,\mathcal{O}_X)\to(S,\mathcal{O}_S)$ 与 $\mathcal{G}$
如引理 \ref{defos-lemma-extensions-of-relative-ringed-spaces} 所述。
考虑引理中的扩张
$0 \to \mathcal{G} \to \mathcal{O}_{X'} \to \mathcal{O}_X \to 0$
。可以选取集合层 $\mathcal{E}$ 以及交换图
$$
\xymatrix{
\mathcal{E} \ar[d]_{\alpha'} \ar[rd]^\alpha \\
\mathcal{O}_{X'} \ar[r] & \mathcal{O}_X
}
$$
使 $f^{-1}\mathcal{O}_S[\mathcal{E}]\to\mathcal{O}_X$ 是满的，核为
$\mathcal{J}$。（例如，可以取任意满射到 $\mathcal{O}_{X'}$ 的
集合层。）于是
$$
\NL_{X/S} \cong \NL(\alpha) = 
\left(
\mathcal{J}/\mathcal{J}^2
\longrightarrow
\Omega_{f^{-1}\mathcal{O}_S[\mathcal{E}]/f^{-1}\mathcal{O}_S}
\otimes_{f^{-1}\mathcal{O}_S[\mathcal{E}]} \mathcal{O}_X\right)
$$
见《模层》第 \ref{modules-section-netherlander} 节，特别是引理
\ref{modules-lemma-NL-up-to-qis}。当然，$\alpha'$ 确定同态
$f^{-1}\mathcal{O}_S[\mathcal{E}]\to\mathcal{O}_{X'}$，后者又确定同态
$$
\mathcal{J}/\mathcal{J}^2 \longrightarrow \mathcal{G}
$$
后者进一步确定
$\Ext^1_{\mathcal{O}_X}(\NL(\alpha), \mathcal{G}) =
\Ext^1_{\mathcal{O}_X}(\NL_{X/S}, \mathcal{G})$
中经本引理双射与 $\mathcal{O}_{X'}$ 对应的元素。
\end{remark}

\begin{lemma}
\label{defos-lemma-extensions-of-relative-ringed-spaces-functorial}
设 $f:(X,\mathcal{O}_X)\to(S,\mathcal{O}_S)$ 与
$g:(Y,\mathcal{O}_Y)\to(X,\mathcal{O}_X)$ 为环化空间的态射。
设 $\mathcal{F}$ 为 $\mathcal{O}_X$-模，$\mathcal{G}$ 为
$\mathcal{O}_Y$-模，$c:\mathcal{F}\to\mathcal{G}$ 为 $g$-映射。
最后，考虑
\begin{enumerate}
\item[(a)] 与 $\xi\in\Ext^1_{\mathcal{O}_X}(\NL_{X/S},\mathcal{F})$
对应的 $f^{-1}\mathcal{O}_S$-代数扩张
$0\to\mathcal{F}\to\mathcal{O}_{X'}\to\mathcal{O}_X\to0$；
\item[(b)] 与 $\zeta\in\Ext^1_{\mathcal{O}_Y}(\NL_{Y/S},\mathcal{G})$
对应的 $g^{-1}f^{-1}\mathcal{O}_S$-代数扩张
$0\to\mathcal{G}\to\mathcal{O}_{Y'}\to\mathcal{O}_Y\to0$。
\end{enumerate}
参见引理 \ref{defos-lemma-extensions-of-relative-ringed-spaces}。那么，存在与
$g$ 和 $c$ 相容的 $S$-态射 $g':Y'\to X'$，当且仅当 $\xi$ 与
$\zeta$ 在
$\Ext^1_{\mathcal{O}_Y}(Lg^*\NL_{X/S},\mathcal{G})$ 中的像相同。
\end{lemma}

\begin{proof}
这个陈述是有意义的，因为我们有同态
$$
\Ext^1_{\mathcal{O}_X}(\NL_{X/S}, \mathcal{F}) \to
\Ext^1_{\mathcal{O}_Y}(Lg^*\NL_{X/S}, Lg^*\mathcal{F}) \to
\Ext^1_{\mathcal{O}_Y}(Lg^*\NL_{X/S}, \mathcal{G})
$$
它使用同态
$Lg^*\mathcal{F}\to g^*\mathcal{F}\xrightarrow{c}\mathcal{G}$；还有
$$
\Ext^1_{\mathcal{O}_Y}(\NL_{Y/S}, \mathcal{G}) \to
\Ext^1_{\mathcal{O}_Y}(Lg^*\NL_{X/S}, \mathcal{G})
$$
它使用同态 $Lg^*\NL_{X/S}\to\NL_{Y/S}$。把引理
\ref{defos-lemma-huge-diagram-ringed-spaces} 应用于下图，即可推出本引理：
$$
\xymatrix{
& 0 \ar[r] &
\mathcal{G} \ar[r] &
\mathcal{O}_{Y'} \ar[r] &
\mathcal{O}_Y \ar[r] & 0 \\
& 0 \ar[r]|\hole & 0 \ar[u] \ar[r] &
\mathcal{O}_S \ar[u] \ar[r]|\hole &
\mathcal{O}_S \ar[u] \ar[r] & 0 \\
0 \ar[r] &
\mathcal{F} \ar[ruu] \ar[r] &
\mathcal{O}_{X'} \ar[r] &
\mathcal{O}_X \ar[ruu] \ar[r] & 0 \\
0 \ar[r] & 0 \ar[ruu]|\hole \ar[u] \ar[r] &
\mathcal{O}_S \ar[ruu]|\hole \ar[u] \ar[r] &
\mathcal{O}_S \ar[ruu]|\hole \ar[u] \ar[r] & 0
}
$$
同时还要使用引理 \ref{defos-lemma-extensions-of-relative-ringed-spaces} 与
\ref{defos-lemma-huge-diagram-ringed-spaces} 的证明中各项构造之间的一种
相容性；其陈述和证明从略。（直接论证见下述评注。）
\end{proof}

\begin{remark}
\label{defos-remark-extensions-of-relative-ringed-spaces-functorial}
设 $f:(X,\mathcal{O}_X)\to(S,\mathcal{O}_S)$、
$g:(Y,\mathcal{O}_Y)\to(X,\mathcal{O}_X)$、$\mathcal{F}$、
$\mathcal{G}$、$c:\mathcal{F}\to\mathcal{G}$、
$0\to\mathcal{F}\to\mathcal{O}_{X'}\to\mathcal{O}_X\to0$、
$\xi\in\Ext^1_{\mathcal{O}_X}(\NL_{X/S},\mathcal{F})$、
$0\to\mathcal{G}\to\mathcal{O}_{Y'}\to\mathcal{O}_Y\to0$ 以及
$\zeta\in\Ext^1_{\mathcal{O}_Y}(\NL_{Y/S},\mathcal{G})$
如引理 \ref{defos-lemma-extensions-of-relative-ringed-spaces-functorial}
所述。沿 $c:g^{-1}\mathcal{F}\to\mathcal{G}$ 作推出，可构造扩张
$$
\xymatrix{
0 \ar[r] &
\mathcal{G} \ar[r] &
\mathcal{O}'_1 \ar[r] &
g^{-1}\mathcal{O}_X \ar[r] & 0 \\
0 \ar[r] &
g^{-1}\mathcal{F} \ar[u]^c \ar[r] &
g^{-1}\mathcal{O}_{X'} \ar[u] \ar[r] &
g^{-1}\mathcal{O}_X \ar@{=}[u] \ar[r] & 0
}
$$
沿 $g^\sharp:g^{-1}\mathcal{O}_X\to\mathcal{O}_Y$ 作拉回，可构造扩张
$$
\xymatrix{
0 \ar[r] &
\mathcal{G} \ar[r] &
\mathcal{O}_{Y'} \ar[r] &
\mathcal{O}_Y \ar[r] & 0 \\
0 \ar[r] &
\mathcal{G} \ar@{=}[u] \ar[r] &
\mathcal{O}'_2 \ar[u] \ar[r] &
g^{-1}\mathcal{O}_X \ar[u] \ar[r] & 0
}
$$
追图表明，存在与 $g$ 和 $c$ 相容的 $S$-态射 $Y'\to X'$，当且仅当
把 $\mathcal{O}'_1$ 与 $\mathcal{O}'_2$ 都看作
$g^{-1}\mathcal{O}_X$ 被 $\mathcal{G}$ 扩张所得的
$g^{-1}f^{-1}\mathcal{O}_S$-代数时二者同构。由引理
\ref{defos-lemma-extensions-of-relative-ringed-spaces}，这些扩张由下式左端分类：
$$
\Ext^1_{g^{-1}\mathcal{O}_X}(
\NL_{g^{-1}\mathcal{O}_X/g^{-1}f^{-1}\mathcal{O}_S}, \mathcal{G}) =
\Ext^1_{\mathcal{O}_Y}(Lg^*\NL_{X/S}, \mathcal{G})
$$
这里的等号来自张量--Hom 伴随性以及等式
$$
\NL_{g^{-1}\mathcal{O}_X/g^{-1}f^{-1}\mathcal{O}_S} = g^{-1}\NL_{X/S}
\quad\text{且}\quad
Lg^*\NL_{X/S} =
g^{-1}\NL_{X/S} \otimes_{g^{-1}\mathcal{O}_X}^\mathbf{L} \mathcal{O}_Y
$$
第一个等式见《模层》引理 \ref{modules-lemma-pullback-NL}；第二个等式
由导出拉回的定义推出。因此，为证明引理
\ref{defos-lemma-extensions-of-relative-ringed-spaces-functorial}，
只需说明 $\mathcal{O}'_1$ 对应于 $\xi$ 的像，而
$\mathcal{O}'_2$ 对应于 $\zeta$ 的像。$\xi$ 与
$\mathcal{O}'_1$ 之间的对应由评注
\ref{defos-remark-extensions-of-relative-ringed-spaces} 中类 $\xi$ 的构造
立即得到。为说明 $\zeta$ 与 $\mathcal{O}'_2$ 之间的对应，先选取交换图
$$
\xymatrix{
\mathcal{E} \ar[d]_{\beta'} \ar[rd]^\beta \\
\mathcal{O}_{Y'} \ar[r] & \mathcal{O}_Y
}
$$
使 $g^{-1}f^{-1}\mathcal{O}_S[\mathcal{E}]\to\mathcal{O}_Y$
是满的，核为 $\mathcal{K}$。再选取交换图
$$
\xymatrix{
\mathcal{E} \ar[d]_{\beta'} &
\mathcal{E}' \ar[l]^\varphi \ar[d]_{\alpha'} \ar[rd]^\alpha \\
\mathcal{O}_{Y'} &
\mathcal{O}'_2 \ar[l] \ar[r] &
g^{-1}\mathcal{O}_X
}
$$
使 $g^{-1}f^{-1}\mathcal{O}_S[\mathcal{E}']\to
g^{-1}\mathcal{O}_X$ 是满的，核为 $\mathcal{J}$。（例如，只需把
$\mathcal{E}'=\mathcal{E}\amalg\mathcal{O}'_2$ 视为集合层。）
映射 $\varphi$ 诱导复形同态 $\NL(\alpha)\to\NL(\beta)$
（记号同《模层》第 \ref{modules-section-netherlander} 节），特别地诱导
$\bar\varphi:\mathcal{J}/\mathcal{J}^2\to
\mathcal{K}/\mathcal{K}^2$。于是 $\NL(\alpha)\cong\NL_{Y/S}$，且
$\NL(\beta) \cong \NL_{g^{-1}\mathcal{O}_X/g^{-1}f^{-1}\mathcal{O}_S}$
，而复形同态 $\NL(\alpha)\to\NL(\beta)$ 表示引理
\ref{defos-lemma-extensions-of-relative-ringed-spaces-functorial} 的陈述中
所用的同态 $Lg^*\NL_{X/S}\to\NL_{Y/S}$（见其证明的第一部分）。
现在，$\zeta$ 对应于由 $\beta'$ 诱导的同态
$\mathcal{K}/\mathcal{K}^2\to\mathcal{G}$ 的类，见评注
\ref{defos-remark-extensions-of-relative-ringed-spaces}。类似地，扩张
$\mathcal{O}'_2$ 对应于由 $\alpha'$ 诱导的同态
$\mathcal{J}/\mathcal{J}^2\to\mathcal{G}$。上述交换图表明，该同态是
由 $\beta'$ 诱导的同态
$\mathcal{K}/\mathcal{K}^2\to\mathcal{G}$ 与
$\bar\varphi:\mathcal{J}/\mathcal{J}^2\to
\mathcal{K}/\mathcal{K}^2$ 的复合。这证明了所需的相容性。
\end{remark}

\begin{lemma}
\label{defos-lemma-parametrize-solutions-ringed-spaces}
设 $t:(S,\mathcal{O}_S)\to(S',\mathcal{O}_{S'})$、
$\mathcal{J}=\Ker(t^\sharp)$、
$f:(X,\mathcal{O}_X)\to(S,\mathcal{O}_S)$、$\mathcal{G}$ 以及
$c:\mathcal{J}\to\mathcal{G}$ 如
(\ref{defos-equation-to-solve-ringed-spaces}) 所述。记
$\xi\in\Ext^1_{\mathcal{O}_S}(\NL_{S/S'},\mathcal{J})$ 为经引理
\ref{defos-lemma-extensions-of-relative-ringed-spaces} 与
$\mathcal{O}_S$ 被 $\mathcal{J}$ 扩张所得的 $\mathcal{O}_{S'}$
相对应的元素。解的同构类集合与同态
$$
\Ext^1_{\mathcal{O}_X}(\NL_{X/S'}, \mathcal{G}) \to
\Ext^1_{\mathcal{O}_X}(Lf^*\NL_{S/S'}, \mathcal{G})
$$
在 $\xi$ 的像上方的纤维典范双射。
\end{lemma}

\begin{proof}
把引理 \ref{defos-lemma-extensions-of-relative-ringed-spaces} 应用于
$X\to S'$ 和 $\mathcal{O}_X$-模 $\mathcal{G}$，可见
$\Ext^1_{\mathcal{O}_X}(\NL_{X/S'},\mathcal{G})$ 中的元素 $\zeta$
参数化 $f^{-1}\mathcal{O}_{S'}$-代数扩张
$0 \to \mathcal{G} \to \mathcal{O}_{X'} \to \mathcal{O}_X \to 0$
。把引理 \ref{defos-lemma-extensions-of-relative-ringed-spaces-functorial}
应用于 $X\to S\to S'$ 和 $c:\mathcal{J}\to\mathcal{G}$，可见存在与
$c$ 和 $f:X\to S$ 相容的 $S'$-态射 $X'\to S'$，当且仅当
$\zeta$ 映到 $\xi$。这当然等价于说 $\mathcal{O}_{X'}$ 是
(\ref{defos-equation-to-solve-ringed-spaces}) 的一个解。
\end{proof}

\begin{remark}
\label{defos-remark-parametrize-solutions-ringed-spaces}
在引理 \ref{defos-lemma-parametrize-solutions-ringed-spaces} 的情形中，
有复形同态
$$
Lf^*\NL_{S'/S} \to \NL_{X/S'} \to \NL_{X/S}
$$
% P11-E0003: frozen source complex retained; see control/ERRATA_CANDIDATES.jsonl.
这些同态几乎构成一个可区别三角，见《模层》引理
\ref{modules-lemma-exact-sequence-NL-ringed-topoi}。如果它确实是可区别
三角，便可得出：$\xi$ 在
$\Ext^2_{\mathcal{O}_X}(\NL_{X/S},\mathcal{G})$ 中的像，就是
(\ref{defos-equation-to-solve-ringed-spaces}) 存在解的障碍。
\end{remark}













\section{概形的形变}
\label{defos-section-deformations-schemes}

\noindent
本节具体说明第 \ref{defos-section-deformations-ringed-spaces} 节的结果对于
概形形变意味着什么。

\begin{lemma}
\label{defos-lemma-deform}
设 $S\subset S'$ 为概形的一阶增厚，$f:X\to S$ 为概形的平坦态射。
如果存在概形的平坦态射 $f':X'\to S'$ 以及 $S$ 上的同构
$a:X\to X'\times_{S'}S$，那么：
\begin{enumerate}
\item 二元组 $(f':X'\to S',a)$ 的同构类集合是
$\Ext^1_{\mathcal{O}_X}(\NL_{X/S},f^*\mathcal{C}_{S/S'})$
作用下的主齐性空间；
\item $S'$ 上在 $X'\times_{S'}S$ 上约化为恒等态射的自同构
$\varphi:X'\to X'$ 所成的集合为
$\Ext^0_{\mathcal{O}_X}(\NL_{X/S},f^*\mathcal{C}_{S/S'})$。
\end{enumerate}
\end{lemma}

\begin{proof}
先注意，《更多态射》第 \ref{more-morphisms-section-thickenings} 节所定义的
概形增厚，正是第 \ref{defos-section-thickenings-spaces} 节意义下为增厚的
概形态射。可以把 $X$ 视为 $X'$ 的闭子概形，使得
$(f,f'):(X\subset X')\to(S\subset S')$ 是一阶增厚的态射。于是，由
《更多态射》引理 \ref{more-morphisms-lemma-deform}（或更一般的引理
\ref{defos-lemma-deform-module}），$X$ 在 $X'$ 中的理想层等于
$f^*\mathcal{C}_{S/S'}$。因此有交换图
$$
\xymatrix{
0 \ar[r] & f^*\mathcal{C}_{S/S'} \ar[r] &
\mathcal{O}_{X'} \ar[r] &
\mathcal{O}_X \ar[r] & 0 \\
0 \ar[r] & \mathcal{C}_{S/S'} \ar[u] \ar[r] &
\mathcal{O}_{S'} \ar[u] \ar[r] &
\mathcal{O}_S \ar[u] \ar[r] & 0
}
$$
其中竖直箭头为 $f$-映射；请与
(\ref{defos-equation-to-solve-ringed-spaces}) 比较。因此，第 (1) 点由引理
\ref{defos-lemma-choices-ringed-spaces} 推出，第 (2) 点由引理
\ref{defos-lemma-huge-diagram-ringed-spaces} 的第 (2) 点推出。（注意，
《更多态射》第 \ref{more-morphisms-section-netherlander} 节为概形态射
所定义的 $\NL_{X/S}$，与第 \ref{defos-section-deformations-ringed-spaces} 节
所用的 $\NL_{X/S}$ 一致。）
\end{proof}










\section{环化拓扑斯的增厚}
\label{defos-section-thickenings-ringed-topoi}

\noindent
本节是第 \ref{defos-section-thickenings-spaces} 节对于环化拓扑斯的对应版本。
在下面几节中，我们将使用以下概念：
\begin{enumerate}
\item 环化拓扑斯 $(\Sh(\mathcal{D}),\mathcal{O}')$ 上的理想层
$\mathcal{I}\subset\mathcal{O}'$ 称为{\it 局部幂零的}，如果
$\mathcal{I}$ 的任意局部截面都是局部幂零的。
\item 环化拓扑斯的一个{\it 增厚}是环化拓扑斯的态射
$i : (\Sh(\mathcal{C}), \mathcal{O}) \to (\Sh(\mathcal{D}), \mathcal{O}')$
，满足：
\begin{enumerate}
\item $i_*$ 是等价 $\Sh(\mathcal{C})\to\Sh(\mathcal{D})$；
\item 同态 $i^\sharp:\mathcal{O}'\to i_*\mathcal{O}$ 是满的；
\item $i^\sharp$ 的核是局部幂零理想层。
\end{enumerate}
\item 环化拓扑斯的一个{\it 一阶增厚}是环化拓扑斯的增厚
$i : (\Sh(\mathcal{C}), \mathcal{O}) \to (\Sh(\mathcal{D}), \mathcal{O}')$
，且 $\Ker(i^\sharp)$ 平方为零。
\item 显然可以定义{\it 环化拓扑斯增厚的态射}、
{\it 基环化拓扑斯上环化拓扑斯增厚的态射}等概念。
\end{enumerate}
若
$i : (\Sh(\mathcal{C}), \mathcal{O}) \to (\Sh(\mathcal{D}), \mathcal{O}')$
是环化拓扑斯的增厚，我们便等同底层拓扑斯，并把 $\mathcal{O}$、
$\mathcal{O}'$ 以及 $\mathcal{I}=\Ker(i^\sharp)$ 都视为
$\mathcal{C}$ 上的层。由此得到 $\mathcal{O}'$-模的短正合列
$$
0 \to \mathcal{I} \to \mathcal{O}' \to \mathcal{O} \to 0
$$
由《位点上的模》引理 \ref{sites-modules-lemma-i-star-equivalence}，
$\mathcal{O}$-模范畴等价于被 $\mathcal{I}$ 零化的
$\mathcal{O}'$-模范畴。特别地，如果 $i$ 是一阶增厚，那么
$\mathcal{I}$ 是 $\mathcal{O}$-模。

\begin{situation}
\label{defos-situation-morphism-thickenings-ringed-topoi}
环化拓扑斯增厚的一个态射 $(f,f')$ 由如下交换图给出：
\begin{equation}
\label{defos-equation-morphism-thickenings-ringed-topoi}
\vcenter{
\xymatrix{
(\Sh(\mathcal{C}), \mathcal{O}) \ar[r]_i \ar[d]_f &
(\Sh(\mathcal{D}), \mathcal{O}') \ar[d]^{f'} \\
(\Sh(\mathcal{B}), \mathcal{O}_\mathcal{B}) \ar[r]^t &
(\Sh(\mathcal{B}'), \mathcal{O}_{\mathcal{B}'})
}
}
\end{equation}
其中各项为环化拓扑斯，水平箭头为增厚。在这一情形中，令
$\mathcal{I}=\Ker(i^\sharp)\subset\mathcal{O}'$、
$\mathcal{J}=\Ker(t^\sharp)\subset\mathcal{O}_{\mathcal{B}'}$。
由于在底层拓扑斯上 $f=f'$，我们等同逆像函子 $f^{-1}$ 与
$(f')^{-1}$。注意，
$(f')^\sharp : f^{-1}\mathcal{O}_{\mathcal{B}'} \to \mathcal{O}'$
特别地诱导同态 $f^{-1}\mathcal{J}\to\mathcal{I}$，从而诱导
$\mathcal{O}'$-模同态
$$
(f')^*\mathcal{J} \longrightarrow \mathcal{I}
$$
如果 $i$ 与 $t$ 都是一阶增厚，那么
$(f')^*\mathcal{J}=f^*\mathcal{J}$，上述同态就成为
$f^*\mathcal{J}\to\mathcal{I}$。
\end{situation}

\begin{definition}
\label{defos-definition-strict-morphism-thickenings-ringed-topoi}
在情形 \ref{defos-situation-morphism-thickenings-ringed-topoi} 中，如果同态
$(f')^*\mathcal{J}\longrightarrow\mathcal{I}$ 是满的，则称
$(f,f')$ 是一个{\it 增厚的严格态射}。
\end{definition}








\section{环化拓扑斯一阶增厚上的模}
\label{defos-section-modules-thickenings-ringed-topoi}

\noindent
本节讨论模的形变理论所需的一些预备知识。设
% P11-E0004: malformed frozen source object retained; see control/ERRATA_CANDIDATES.jsonl.
$i : (\Sh(\mathcal{C}, \mathcal{O}) \to (\Sh(\mathcal{D}), \mathcal{O}')$
为环化拓扑斯的一阶增厚。我们将随时使用第
\ref{defos-section-thickenings-ringed-topoi} 节引入的记号；特别地，我们等同
底层拓扑斯。本节考虑 $\mathcal{O}'$-模的短正合列
\begin{equation}
\label{defos-equation-extension-ringed-topoi}
0 \to \mathcal{K} \to \mathcal{F}' \to \mathcal{F} \to 0
\end{equation}
其中 $\mathcal{F}$、$\mathcal{K}$ 是 $\mathcal{O}$-模，而
$\mathcal{F}'$ 是 $\mathcal{O}'$-模。在这种情形中，有典范的
$\mathcal{O}$-模同态
$$
c_{\mathcal{F}'} :
\mathcal{I} \otimes_\mathcal{O} \mathcal{F}
\longrightarrow
\mathcal{K}
$$
其中 $\mathcal{I}=\Ker(i^\sharp)$。具体而言，给定 $\mathcal{I}$ 的
局部截面 $f$ 与 $\mathcal{F}$ 的局部截面 $s$，令
$c_{\mathcal{F}'}(f\otimes s)=fs'$，其中 $s'$ 是提升 $s$ 的
$\mathcal{F}'$ 的局部截面。

\begin{lemma}
\label{defos-lemma-inf-map-ringed-topoi}
设
$i:(\Sh(\mathcal{C}),\mathcal{O})\to
(\Sh(\mathcal{D}),\mathcal{O}')$ 是环化拓扑斯的一阶增厚，并给定扩张
$$
0 \to \mathcal{K} \to \mathcal{F}' \to \mathcal{F} \to 0
\quad\text{且}\quad
0 \to \mathcal{L} \to \mathcal{G}' \to \mathcal{G} \to 0
$$
如 (\ref{defos-equation-extension-ringed-topoi}) 所述，以及同态
$\varphi:\mathcal{F}\to\mathcal{G}$ 与
$\psi:\mathcal{K}\to\mathcal{L}$。
\begin{enumerate}
\item 如果存在与 $\varphi$ 和 $\psi$ 相容的 $\mathcal{O}'$-模同态
$\varphi':\mathcal{F}'\to\mathcal{G}'$，那么图表
$$
\xymatrix{
\mathcal{I} \otimes_\mathcal{O} \mathcal{F}
\ar[r]_-{c_{\mathcal{F}'}} \ar[d]_{1 \otimes \varphi} &
\mathcal{K} \ar[d]^\psi \\
\mathcal{I} \otimes_\mathcal{O} \mathcal{G}
\ar[r]^-{c_{\mathcal{G}'}} &
\mathcal{L}
}
$$
交换。
\item 与 $\varphi$ 和 $\psi$ 相容的 $\mathcal{O}'$-模同态
$\varphi':\mathcal{F}'\to\mathcal{G}'$ 的集合，如果非空，就是
$\Hom_\mathcal{O}(\mathcal{F},\mathcal{L})$ 作用下的主齐性空间。
\end{enumerate}
\end{lemma}

\begin{proof}
第 (1) 点由各同态的描述立即得到。对于第 (2) 点，若 $\varphi'$ 与
$\varphi''$ 是两个与 $\varphi$ 和 $\psi$ 相容的同态
$\mathcal{F}'\to\mathcal{G}'$，那么 $\varphi'-\varphi''$ 分解为
$$
\mathcal{F}' \to \mathcal{F} \to \mathcal{L} \to \mathcal{G}'
$$
由《位点上的模》引理 \ref{sites-modules-lemma-i-star-equivalence}，
中间的同态来自 $\Hom_\mathcal{O}(\mathcal{F},\mathcal{L})$ 中唯一的
元素。反过来，给定该群中的元素 $\alpha$，可以把上述复合（中间同态取
$\alpha$）加到 $\varphi'$ 上。若干细节从略。
\end{proof}

\begin{lemma}
\label{defos-lemma-inf-obs-map-ringed-topoi}
设
$i:(\Sh(\mathcal{C}),\mathcal{O})\to
(\Sh(\mathcal{D}),\mathcal{O}')$ 是环化拓扑斯的一阶增厚，并给定扩张
$$
0 \to \mathcal{K} \to \mathcal{F}' \to \mathcal{F} \to 0
\quad\text{且}\quad
0 \to \mathcal{L} \to \mathcal{G}' \to \mathcal{G} \to 0
$$
如 (\ref{defos-equation-extension-ringed-topoi}) 所述，以及同态
$\varphi:\mathcal{F}\to\mathcal{G}$ 与
$\psi:\mathcal{K}\to\mathcal{L}$。假设图表
$$
\xymatrix{
\mathcal{I} \otimes_\mathcal{O} \mathcal{F}
\ar[r]_-{c_{\mathcal{F}'}} \ar[d]_{1 \otimes \varphi} &
\mathcal{K} \ar[d]^\psi \\
\mathcal{I} \otimes_\mathcal{O} \mathcal{G}
\ar[r]^-{c_{\mathcal{G}'}} &
\mathcal{L}
}
$$
交换。那么存在元素
$$
o(\varphi, \psi) \in
\Ext^1_\mathcal{O}(\mathcal{F}, \mathcal{L})
$$
它为零，当且仅当存在与 $\varphi$ 和 $\psi$ 相容的同态
$\varphi':\mathcal{F}'\to\mathcal{G}'$。
\end{lemma}

\begin{proof}
可以显式构造扩张
$$
0 \to \mathcal{L} \to \mathcal{H} \to \mathcal{F} \to 0
$$
方法是令 $\mathcal{H}$ 为复形
$$
\mathcal{K}
\xrightarrow{1, - \psi}
\mathcal{F}' \oplus \mathcal{G}' \xrightarrow{\varphi, 1}
\mathcal{G}
$$
在中间项处的上同调层（记号含义显然）。利用引理中图表交换这一假设，
对局部截面作计算可知 $\mathcal{H}$ 被 $\mathcal{I}$ 零化。因此
$\mathcal{H}$ 在下述群中定义一个类：
$$
\Ext^1_\mathcal{O}(\mathcal{F}, \mathcal{L})
\subset
\Ext^1_{\mathcal{O}'}(\mathcal{F}, \mathcal{L})
$$
最后，$\mathcal{H}$ 的类等于沿 $\psi$ 对扩张 $\mathcal{F}'$ 所作的
推出与沿 $\varphi$ 对扩张 $\mathcal{G}'$ 所作的拉回之差（计算从略）。
因此，$\mathcal{H}$ 的类为零，等价于存在交换图
$$
\xymatrix{
0 \ar[r] &
\mathcal{K} \ar[r] \ar[d]_{\psi} &
\mathcal{F}' \ar[r] \ar[d]_{\varphi'} &
\mathcal{F} \ar[r] \ar[d]_\varphi & 0\\
0 \ar[r] &
\mathcal{L} \ar[r] &
\mathcal{G}' \ar[r] &
\mathcal{G} \ar[r] & 0
}
$$
这正是所需。
\end{proof}

\begin{lemma}
\label{defos-lemma-inf-ext-ringed-topoi}
设
$i:(\Sh(\mathcal{C}),\mathcal{O})\to
(\Sh(\mathcal{D}),\mathcal{O}')$ 是环化拓扑斯的一阶增厚。给定
$\mathcal{O}$-模 $\mathcal{F}$、$\mathcal{K}$ 以及
$\mathcal{O}$-线性同态
$c:\mathcal{I}\otimes_\mathcal{O}\mathcal{F}\to\mathcal{K}$。
如果存在满足 $c_{\mathcal{F}'}=c$ 的正合列
(\ref{defos-equation-extension-ringed-topoi})，那么这些扩张的同构类集合是
$\Ext^1_\mathcal{O}(\mathcal{F},\mathcal{K})$ 作用下的主齐性空间。
\end{lemma}

\begin{proof}
假设给定扩张
$$
0 \to \mathcal{K} \to \mathcal{F}'_1 \to \mathcal{F} \to 0
\quad\text{且}\quad
0 \to \mathcal{K} \to \mathcal{F}'_2 \to \mathcal{F} \to 0
$$
且 $c_{\mathcal{F}'_1}=c_{\mathcal{F}'_2}=c$。那么二者之差
（在扩张群中取差，见《同调代数》第 \ref{homology-section-extensions} 节）
是扩张
$$
0 \to \mathcal{K} \to \mathcal{E} \to \mathcal{F} \to 0
$$
其中 $\mathcal{E}$ 被 $\mathcal{I}$ 零化（局部计算从略）。所以该正合列
是 $\mathcal{O}$-模的扩张，见《位点上的模》引理
\ref{sites-modules-lemma-i-star-equivalence}。反过来，给定这样的扩张
$\mathcal{E}$，可将它加到 $\mathcal{O}'$-扩张 $\mathcal{F}'$ 上，
而不改变同态 $c_{\mathcal{F}'}$。若干细节从略。
\end{proof}

\begin{lemma}
\label{defos-lemma-inf-obs-ext-ringed-topoi}
设 $i : (\Sh(\mathcal{C}), \mathcal{O}) \to
(\Sh(\mathcal{D}), \mathcal{O}')$ 是环化拓扑斯的一阶增厚。给定
$\mathcal{O}$-模 $\mathcal{F}$、$\mathcal{K}$ 以及 $\mathcal{O}$-线性同态
$c : \mathcal{I} \otimes_\mathcal{O} \mathcal{F} \to \mathcal{K}$。
那么存在一个元素
$$
o(\mathcal{F}, \mathcal{K}, c) \in
\Ext^2_\mathcal{O}(\mathcal{F}, \mathcal{K})
$$
其消失是存在满足 $c_{\mathcal{F}'}=c$ 的正合列
(\ref{defos-equation-extension-ringed-topoi}) 的充要条件。
\end{lemma}

\begin{proof}
我们先证明：若 $\mathcal{K}$ 是内射 $\mathcal{O}$-模，
则确实存在满足 $c_{\mathcal{F}'}=c$ 的正合列
(\ref{defos-equation-extension-ringed-topoi})。为此，取一个平坦
$\mathcal{O}'$-模 $\mathcal{H}'$ 以及满射
$\mathcal{H}' \to \mathcal{F}$
（《位点上的模》引理 \ref{sites-modules-lemma-module-quotient-flat}）。
令 $\mathcal{J}\subset\mathcal{H}'$ 为其核。由于 $\mathcal{H}'$ 平坦，故有
$$
\mathcal{I} \otimes_{\mathcal{O}'} \mathcal{H}' =
\mathcal{I}\mathcal{H}'
\subset \mathcal{J} \subset \mathcal{H}'
$$
注意，同态
$$
\mathcal{I}\mathcal{H}' =
\mathcal{I} \otimes_{\mathcal{O}'} \mathcal{H}'
\longrightarrow
\mathcal{I} \otimes_{\mathcal{O}'} \mathcal{F} =
\mathcal{I} \otimes_\mathcal{O} \mathcal{F}
$$
零化 $\mathcal{I}\mathcal{J}$。事实上，若 $f$ 是 $\mathcal{I}$ 的局部截面，
$s$ 是 $\mathcal{H}$ 的局部截面，则 $fs$ 映为
$f\otimes\overline{s}$，其中 $\overline{s}$ 是 $s$ 在
$\mathcal{F}$ 中的像。因此得到 $\mathcal{O}$-模的图
$$
\xymatrix{
\mathcal{I}\mathcal{H}'/\mathcal{I}\mathcal{J}
\ar@{^{(}->}[r] \ar[d] &
\mathcal{J}/\mathcal{I}\mathcal{J} \ar@{..>}[d]_\gamma \\
\mathcal{I} \otimes_\mathcal{O} \mathcal{F} \ar[r]^-c &
\mathcal{K}
}
$$
若 $\mathcal{K}$ 作为 $\mathcal{O}$-模是内射的，就得到虚线箭头。
以 $\gamma':\mathcal{J}\to\mathcal{K}$ 记 $\gamma$ 与
$\mathcal{J}\to\mathcal{J}/\mathcal{I}\mathcal{J}$ 的复合。
局部计算表明，推出
$$
\xymatrix{
0 \ar[r] &
\mathcal{J} \ar[r] \ar[d]_{\gamma'} &
\mathcal{H}' \ar[r] \ar[d] &
\mathcal{F} \ar[r] \ar@{=}[d] &
0 \\
0 \ar[r] &
\mathcal{K} \ar[r] &
\mathcal{F}' \ar[r] &
\mathcal{F} \ar[r] &
0
}
$$
给出了本引理所述问题的一个解。

\medskip\noindent
一般情形。取嵌入 $\mathcal{K}\subset\mathcal{K}'$，其中
$\mathcal{K}'$ 是内射 $\mathcal{O}$-模。令 $\mathcal{Q}$ 为其商，
于是有正合列
$$
0 \to \mathcal{K} \to \mathcal{K}' \to \mathcal{Q} \to 0
$$
以 $c':\mathcal{I}\otimes_\mathcal{O}\mathcal{F}\to\mathcal{K}'$
记相应的复合同态。由上一段，存在正合列
$$
0 \to \mathcal{K}' \to \mathcal{E}' \to \mathcal{F} \to 0
$$
如 (\ref{defos-equation-extension-ringed-topoi}) 所示，且
$c_{\mathcal{E}'}=c'$。注意，$c'$ 与
$\mathcal{K}'\to\mathcal{Q}$ 的复合为零，故沿
$\mathcal{K}'\to\mathcal{Q}$ 对 $\mathcal{E}'$ 取推出所得扩张
$$
0 \to \mathcal{Q} \to \mathcal{D}' \to \mathcal{F} \to 0
$$
如 (\ref{defos-equation-extension-ringed-topoi}) 所示，且
$c_{\mathcal{D}'}=0$。这恰好意味着 $\mathcal{D}'$ 被
$\mathcal{I}$ 零化；换言之，$\mathcal{D}'$ 是 $\mathcal{O}$-模的扩张，
因而定义了元素
$$
o(\mathcal{F}, \mathcal{K}, c) \in
\Ext^1_\mathcal{O}(\mathcal{F}, \mathcal{Q}) =
\Ext^2_\mathcal{O}(\mathcal{F}, \mathcal{K})
$$
（该等式来自上述正合列所诱导的上同调长正合列，以及以内射模
$\mathcal{K}'$ 为第二变元的高阶 $\Ext$ 群消失）。若
$o(\mathcal{F},\mathcal{K},c)=0$，则可取分裂
$s:\mathcal{F}\to\mathcal{D}'$，并令
$$
\mathcal{F}' = \Ker(\mathcal{E}' \to \mathcal{D}'/s(\mathcal{F}))
$$
从而得到下图
$$
\xymatrix{
0 \ar[r] &
\mathcal{K} \ar[r] \ar[d] &
\mathcal{F}' \ar[r] \ar[d] &
\mathcal{F} \ar[r] \ar@{=}[d] &
0 \\
0 \ar[r] &
\mathcal{K}' \ar[r] &
\mathcal{E}' \ar[r] &
\mathcal{F} \ar[r] & 0
}
$$
其两行均正合，由此可见 $c_{\mathcal{F}'}=c$。反过来，若
$\mathcal{F}'$ 存在，则由引理 \ref{defos-lemma-inf-ext-ringed-topoi}
及以内射模 $\mathcal{K}'$ 为第二变元的高阶 $\Ext$ 群消失，
沿 $\mathcal{K}\to\mathcal{K}'$ 对 $\mathcal{F}'$ 取推出所得扩张
同构于 $\mathcal{E}'$。这给出如上图，从而说明
$\mathcal{D}'$ 作为扩张是分裂的，即类
$o(\mathcal{F},\mathcal{K},c)$ 为零。
\end{proof}

\begin{remark}
\label{defos-remark-trivial-thickening-ringed-topoi}
设 $(\Sh(\mathcal{C}),\mathcal{O})$ 是环化拓扑斯。若存在环化拓扑斯态射
$\pi : (\Sh(\mathcal{D}), \mathcal{O}') \to (\Sh(\mathcal{C}), \mathcal{O})$
作为 $i:(\Sh(\mathcal{C}),\mathcal{O})\to
(\Sh(\mathcal{D}),\mathcal{O}')$ 的左逆，就称这一阶增厚为
{\it 平凡的}。这样一个态射 $\pi$ 的选取称为该一阶增厚的一个
{\it 平凡化}。给定 $\pi$，利用 $\pi^\sharp$ 分裂满射
$\mathcal{O}'\to\mathcal{O}$，就得到 $\mathcal{C}$ 上代数层的分解
\begin{equation}
\label{defos-equation-splitting-ringed-topoi}
\mathcal{O}' = \mathcal{O} \oplus \mathcal{I}
\end{equation}
反过来，这样的分解决定一个态射 $\pi$。
$(\Sh(\mathcal{C}),\mathcal{O})$ 的带平凡化一阶增厚所成范畴
等价于 $\mathcal{O}$-模范畴。
\end{remark}

\begin{remark}
\label{defos-remark-trivial-extension-ringed-topoi}
设 $i:(\Sh(\mathcal{C}),\mathcal{O})\to
(\Sh(\mathcal{D}),\mathcal{O}')$ 是环化拓扑斯的平凡一阶增厚，
$\pi:(\Sh(\mathcal{D}),\mathcal{O}')\to
(\Sh(\mathcal{C}),\mathcal{O})$ 是一个平凡化。给定任意三元组
$(\mathcal{F},\mathcal{K},c)$，其中 $\mathcal{F}$、$\mathcal{K}$
是 $\mathcal{O}$-模，且有同态
$c : \mathcal{I} \otimes_\mathcal{O} \mathcal{F} \to \mathcal{K}$
可令
$$
\mathcal{F}'_{c, triv} = \mathcal{F} \oplus \mathcal{K}
$$
并利用与 $\pi$ 相伴的分解
(\ref{defos-equation-splitting-ringed-topoi}) 及同态 $c$ 定义
$\mathcal{O}'$-模结构，从而得到扩张
(\ref{defos-equation-extension-ringed-topoi})。我们称
$\mathcal{F}'_{c,triv}$ 为与 $c$ 和平凡化 $\pi$ 对应的、以
$\mathcal{K}$ 扩张 $\mathcal{F}$ 的{\it 平凡扩张}。
给定任意形如 (\ref{defos-equation-extension-ringed-topoi}) 的扩张
$\mathcal{F}'$，可利用 $\pi^\sharp:\mathcal{O}\to\mathcal{O}'$
把 $\mathcal{F}'$ 视为 $\mathcal{O}$-模扩张，因而得到
$\Ext^1_\mathcal{O}(\mathcal{F},\mathcal{K})$ 中的类
$\xi_{\mathcal{F}'}$。引理 \ref{defos-lemma-inf-ext-ringed-topoi}
保证 $\mathcal{F}'\mapsto\xi_{\mathcal{F}'}$ 诱导双射
$$
\left\{
\begin{matrix}
\text{满足下述条件的扩张 }\mathcal{F}'\text{ 的同构类}\\
\mathcal{F}'\text{ 如 (\ref{defos-equation-extension-ringed-topoi}) 所示，且 }
c=c_{\mathcal{F}'}
\end{matrix}
\right\}
\longrightarrow
\Ext^1_\mathcal{O}(\mathcal{F}, \mathcal{K})
$$
此外，平凡扩张 $\mathcal{F}'_{c,triv}$ 映到零类。
\end{remark}

\begin{remark}
\label{defos-remark-extension-functorial-ringed-topoi}
设 $(\Sh(\mathcal{C}),\mathcal{O})$ 是环化拓扑斯。设
$(\Sh(\mathcal{C}), \mathcal{O}) \to (\Sh(\mathcal{D}_i), \mathcal{O}'_i)$,
$i=1,2$ 是理想层为 $\mathcal{I}_i$ 的一阶增厚。设
$h:(\Sh(\mathcal{D}_1),\mathcal{O}'_1)\to
(\Sh(\mathcal{D}_2), \mathcal{O}'_2)$
是 $(\Sh(\mathcal{C}),\mathcal{O})$ 的一阶增厚之间的态射。图示为
$$
\xymatrix{
& (\Sh(\mathcal{C}), \mathcal{O}) \ar[ld] \ar[rd] & \\
(\Sh(\mathcal{D}_1), \mathcal{O}'_1) \ar[rr]^h & & 
(\Sh(\mathcal{D}_2), \mathcal{O}'_2)
}
$$
注意，$h^\sharp:\mathcal{O}'_2\to\mathcal{O}'_1$
特别诱导 $\mathcal{O}$-模同态 $\mathcal{I}_2\to\mathcal{I}_1$。
设 $\mathcal{F}$ 是 $\mathcal{O}$-模。设
$(\mathcal{K}_i,c_i)$，$i=1,2$，是两个二元组，
其中 $\mathcal{K}_i$ 是 $\mathcal{O}$-模，并有同态
$c_i : \mathcal{I}_i \otimes_\mathcal{O} \mathcal{F} \to
\mathcal{K}_i$。再给定 $\mathcal{O}$-模同态
$\mathcal{K}_2\to\mathcal{K}_1$，使图
$$
\xymatrix{
\mathcal{I}_2 \otimes_\mathcal{O} \mathcal{F}
\ar[r]_-{c_2} \ar[d] &
\mathcal{K}_2 \ar[d] \\
\mathcal{I}_1 \otimes_\mathcal{O} \mathcal{F}
\ar[r]^-{c_1} &
\mathcal{K}_1
}
$$
交换。那么存在典范的函子性映射
$$
\left\{
\begin{matrix}
\mathcal{F}'_2\text{ 如 (\ref{defos-equation-extension-ringed-topoi}) 所示，且 }\\
c_2=c_{\mathcal{F}'_2}\text{、}\mathcal{K}=\mathcal{K}_2
\end{matrix}
\right\}
\longrightarrow
\left\{
\begin{matrix}
\mathcal{F}'_1\text{ 如 (\ref{defos-equation-extension-ringed-topoi}) 所示，且 }\\
c_1=c_{\mathcal{F}'_1}\text{、}\mathcal{K}=\mathcal{K}_1
\end{matrix}
\right\}
$$
具体地，把 $\mathcal{O}$、$\mathcal{O}'_i$、$\mathcal{F}$、
$\mathcal{K}_i$ 等都视为 $\mathcal{C}$ 上的层；给定
$\mathcal{F}'_2$，令层 $\mathcal{F}'_1$ 等于推出，
即嵌入下列扩张图：
$$
\xymatrix{
0 \ar[r] &
\mathcal{K}_2 \ar[r] \ar[d] &
\mathcal{F}'_2 \ar[r] \ar[d] &
\mathcal{F} \ar@{=}[d] \ar[r] & 0 \\
0 \ar[r] &
\mathcal{K}_1 \ar[r] &
\mathcal{F}'_1 \ar[r] &
\mathcal{F} \ar[r] & 0
}
$$
推出上的 $\mathcal{O}'_1$-模结构之构造从略
（这里用到含 $c_1$、$c_2$ 的图的交换性）。
\end{remark}

\begin{remark}
\label{defos-remark-trivial-extension-functorial-ringed-topoi}
设 $(\Sh(\mathcal{C}),\mathcal{O})$、
$(\Sh(\mathcal{C}), \mathcal{O}) \to (\Sh(\mathcal{D}_i), \mathcal{O}'_i)$,
$\mathcal{I}_i$ 及 $h:(\Sh(\mathcal{D}_1),\mathcal{O}'_1)\to
(\Sh(\mathcal{D}_2),\mathcal{O}'_2)$ 如注
\ref{defos-remark-extension-functorial-ringed-topoi}。设给定平凡化
$\pi_i : (\Sh(\mathcal{D}_i), \mathcal{O}'_i) \to
(\Sh(\mathcal{C}), \mathcal{O})$，满足
$\pi_1=h\circ\pi_2$。换言之，设 $h$ 是
$(\Sh(\mathcal{C}),\mathcal{O})$ 的带平凡化一阶增厚之间的态射。
设 $(\mathcal{K}_i,c_i)$，$i=1,2$，是两个二元组，
其中 $\mathcal{K}_i$ 是 $\mathcal{O}$-模，并有同态
$c_i : \mathcal{I}_i \otimes_\mathcal{O} \mathcal{F} \to
\mathcal{K}_i$。再给定 $\mathcal{O}$-模同态
$\mathcal{K}_2\to\mathcal{K}_1$，使下图交换：
$$
\xymatrix{
\mathcal{I}_2 \otimes_\mathcal{O} \mathcal{F}
\ar[r]_-{c_2} \ar[d] &
\mathcal{K}_2 \ar[d] \\
\mathcal{I}_1 \otimes_\mathcal{O} \mathcal{F}
\ar[r]^-{c_1} &
\mathcal{K}_1
}
$$
在此情形，注 \ref{defos-remark-trivial-extension-ringed-topoi} 的构造
诱导交换图
$$
\xymatrix{
\{\mathcal{F}'_2\text{ 如 (\ref{defos-equation-extension-ringed-topoi}) 所示，且 }
c_2=c_{\mathcal{F}'_2}\text{、}\mathcal{K}=\mathcal{K}_2\}
\ar[d] \ar[rr] & &
\Ext^1_\mathcal{O}(\mathcal{F}, \mathcal{K}_2) \ar[d] \\
\{\mathcal{F}'_1\text{ 如 (\ref{defos-equation-extension-ringed-topoi}) 所示，且 }
c_1=c_{\mathcal{F}'_1}\text{、}\mathcal{K}=\mathcal{K}_1\}
\ar[rr] & &
\Ext^1_\mathcal{O}(\mathcal{F}, \mathcal{K}_1)
}
$$
其中右侧竖直箭头由 $\Ext$ 的函子性及同态
$\mathcal{K}_2\to\mathcal{K}_1$ 给出，左侧竖直箭头则是注
\ref{defos-remark-extension-functorial-ringed-topoi} 中的映射。
\end{remark}

\begin{remark}
\label{defos-remark-obstruction-extension-functorial-ringed-topoi}
设 $(\Sh(\mathcal{C}),\mathcal{O})$、
$(\Sh(\mathcal{C}), \mathcal{O}) \to (\Sh(\mathcal{D}_i), \mathcal{O}'_i)$,
$\mathcal{I}_i$ 及 $h:(\Sh(\mathcal{D}_1),\mathcal{O}'_1)\to
(\Sh(\mathcal{D}_2),\mathcal{O}'_2)$ 如注
\ref{defos-remark-extension-functorial-ringed-topoi}。注意，
$h^\sharp:\mathcal{O}'_2\to\mathcal{O}'_1$ 特别诱导
$\mathcal{O}$-模同态 $\mathcal{I}_2\to\mathcal{I}_1$。
设 $\mathcal{F}$ 是 $\mathcal{O}$-模。设
$(\mathcal{K}_i,c_i)$，$i=1,2$，是两个二元组，
其中 $\mathcal{K}_i$ 是 $\mathcal{O}$-模，并有同态
$c_i : \mathcal{I}_i \otimes_\mathcal{O} \mathcal{F} \to
\mathcal{K}_i$。再给定 $\mathcal{O}$-模同态
$\mathcal{K}_2\to\mathcal{K}_1$，使下图交换：
$$
\xymatrix{
\mathcal{I}_2 \otimes_\mathcal{O} \mathcal{F}
\ar[r]_-{c_2} \ar[d] &
\mathcal{K}_2 \ar[d] \\
\mathcal{I}_1 \otimes_\mathcal{O} \mathcal{F}
\ar[r]^-{c_1} &
\mathcal{K}_1
}
$$
我们{\bf 断言}，同态
$$
\Ext^2_\mathcal{O}(\mathcal{F}, \mathcal{K}_2)
\longrightarrow
\Ext^2_\mathcal{O}(\mathcal{F}, \mathcal{K}_1)
$$
把 $o(\mathcal{F},\mathcal{K}_2,c_2)$ 映到
$o(\mathcal{F},\mathcal{K}_1,c_1)$。

\medskip\noindent
为证明此断言，取嵌入 $j_2:\mathcal{K}_2\to\mathcal{K}_2'$，
其中 $\mathcal{K}_2'$ 是内射 $\mathcal{O}$-模。
如引理 \ref{defos-lemma-inf-obs-ext-ringed-topoi} 的证明所示，
可取一个 $\mathcal{O}_2$-模扩张
% P11-E0005: unprimed O_i symbols retained from frozen source; see control/ERRATA_CANDIDATES.jsonl.
$$
0 \to \mathcal{K}_2' \to \mathcal{E}_2 \to \mathcal{F} \to 0
$$
使 $c_{\mathcal{E}_2}=j_2\circ c_2$。引理
\ref{defos-lemma-inf-obs-ext-ringed-topoi} 的证明把
$o(\mathcal{F},\mathcal{K}_2,c_2)$ 构造为下列
$\mathcal{O}$-模正合列的 Yoneda 扩张类（其含义见《导出范畴》第
\ref{derived-section-ext} 节）：
$$
0 \to
\mathcal{K}_2 \to \mathcal{K}_2' \to
\mathcal{E}_2/\mathcal{K}_2 \to
\mathcal{F} \to 0
$$
令 $\mathcal{K}_1'$ 为
$\mathcal{K}_2\to\mathcal{K}_1\oplus\mathcal{K}_2'$ 的余核。
存在单射 $j_1:\mathcal{K}_1\to\mathcal{K}_1'$ 以及同态
$\mathcal{K}_2'\to\mathcal{K}_1'$，二者组成交换方块。取推出：
$$
\xymatrix{
0 \ar[r] &
\mathcal{K}_2' \ar[r] \ar[d] &
\mathcal{E}_2 \ar[r] \ar[d] &
\mathcal{F} \ar[r] \ar[d] & 0 \\
0 \ar[r] &
\mathcal{K}_1' \ar[r] &
\mathcal{E}_1 \ar[r] &
\mathcal{F} \ar[r] & 0
}
$$
$\mathcal{E}_1$ 上有典范的 $\mathcal{O}_1$-模结构，且在此结构下
$c_{\mathcal{E}_1}=j_1\circ c_1$（这里用到上面含
$c_1$、$c_2$ 的图的交换性）。引理
\ref{defos-lemma-inf-obs-ext-ringed-topoi} 的步骤告诉我们，
$o(\mathcal{F},\mathcal{K}_1,c_1)$ 是下列
$\mathcal{O}$-模正合列的 Yoneda 扩张类：
$$
0 \to
\mathcal{K}_1 \to
\mathcal{K}_1' \to
\mathcal{E}_1/\mathcal{K}_1 \to
\mathcal{F} \to 0
$$
由于存在正合列之间的同态
% P11-E0006: duplicated lower row retained from frozen source; see control/ERRATA_CANDIDATES.jsonl.
$$
\xymatrix{
0 \ar[r] &
\mathcal{K}_2 \ar[d] \ar[r] &
\mathcal{K}_2' \ar[d] \ar[r] &
\mathcal{E}_2/\mathcal{K}_2 \ar[r] \ar[d] &
\mathcal{F} \ar[r] \ar@{=}[d] &
0 \\
0 \ar[r] &
\mathcal{K}_2 \ar[r] &
\mathcal{K}_2' \ar[r] &
\mathcal{E}_2/\mathcal{K}_2 \ar[r] &
\mathcal{F} \ar[r] &
0
}
$$
故断言成立。
\end{remark}

\begin{remark}
\label{defos-remark-short-exact-sequence-thickenings-ringed-topoi}
设 $(\Sh(\mathcal{C}),\mathcal{O})$ 是环化拓扑斯。给定其一阶增厚的态射列
$$
(\Sh(\mathcal{D}_1), \mathcal{O}'_1) \to
(\Sh(\mathcal{D}_2), \mathcal{O}'_2) \to
(\Sh(\mathcal{D}_3), \mathcal{O}'_3)
$$
若理想层 $\mathcal{I}_i$ 之间的相应同态给出 $\mathcal{O}$-模复形
$\mathcal{I}_3 \to \mathcal{I}_2 \to \mathcal{I}_1$
（即复合为零），就称上述态射列为一个{\it 复形}。在此情形，复合
$(\Sh(\mathcal{D}_1), \mathcal{O}'_1) \to
(\Sh(\mathcal{D}_3), \mathcal{O}'_3)$ 分解通过
$(\Sh(\mathcal{C}), \mathcal{O}) \to
(\Sh(\mathcal{D}_3), \mathcal{O}'_3)$；换言之，
$(\Sh(\mathcal{D}_1),\mathcal{O}'_1)$ 作为
$(\Sh(\mathcal{C}),\mathcal{O})$ 的一阶增厚是平凡的，并带有典范平凡化
$\pi : (\Sh(\mathcal{D}_1), \mathcal{O}'_1) \to
(\Sh(\mathcal{C}), \mathcal{O})$。

\medskip\noindent
若 $(\Sh(\mathcal{C}),\mathcal{O})$ 的一阶增厚的态射列
$$
(\Sh(\mathcal{D}_1), \mathcal{O}'_1) \to
(\Sh(\mathcal{D}_2), \mathcal{O}'_2) \to
(\Sh(\mathcal{D}_3), \mathcal{O}'_3)
$$
所对应的理想层同态构成 $\mathcal{O}$-模短正合列
$$
0 \to \mathcal{I}_3 \to \mathcal{I}_2 \to \mathcal{I}_1 \to 0
$$
就称这一态射列为{\it 短正合列}。
\end{remark}

\begin{remark}
\label{defos-remark-complex-thickenings-and-ses-modules-ringed-topoi}
设 $(\Sh(\mathcal{C}),\mathcal{O})$ 是环化拓扑斯，
$\mathcal{F}$ 是 $\mathcal{O}$-模。设
$$
(\Sh(\mathcal{D}_1), \mathcal{O}'_1) \to
(\Sh(\mathcal{D}_2), \mathcal{O}'_2) \to
(\Sh(\mathcal{D}_3), \mathcal{O}'_3)
$$
是 $(\Sh(\mathcal{C}),\mathcal{O})$ 的一阶增厚复形，见注
\ref{defos-remark-short-exact-sequence-thickenings-ringed-topoi}。
设 $(\mathcal{K}_i,c_i)$，$i=1,2,3$，是三个二元组，
其中 $\mathcal{K}_i$ 是 $\mathcal{O}$-模，并有同态
$c_i : \mathcal{I}_i \otimes_\mathcal{O} \mathcal{F} \to
\mathcal{K}_i$。设给定 $\mathcal{O}$-模短正合列
$$
0 \to \mathcal{K}_3 \to \mathcal{K}_2 \to \mathcal{K}_1 \to 0
$$
使得
$$
\vcenter{
\xymatrix{
\mathcal{I}_2 \otimes_\mathcal{O} \mathcal{F}
\ar[r]_-{c_2} \ar[d] &
\mathcal{K}_2 \ar[d] \\
\mathcal{I}_1 \otimes_\mathcal{O} \mathcal{F}
\ar[r]^-{c_1} &
\mathcal{K}_1
}
}
\quad\text{且}\quad
\vcenter{
\xymatrix{
\mathcal{I}_3 \otimes_\mathcal{O} \mathcal{F}
\ar[r]_-{c_3} \ar[d] &
\mathcal{K}_3 \ar[d] \\
\mathcal{I}_2 \otimes_\mathcal{O} \mathcal{F}
\ar[r]^-{c_2} &
\mathcal{K}_2
}
}
$$
均为交换图。最后，设给定 $\mathcal{O}'_2$-模扩张
$$
0 \to \mathcal{K}_2 \to \mathcal{F}'_2 \to \mathcal{F} \to 0
$$
如 (\ref{defos-equation-extension-ringed-topoi}) 所示，其中
$\mathcal{K}=\mathcal{K}_2$，且 $c_{\mathcal{F}'_2}=c_2$。
在此情形，可应用注 \ref{defos-remark-extension-functorial-ringed-topoi}
的函子性，得到 $\mathcal{O}'_1$-模扩张 $\mathcal{F}'_1$
（下面会描述这一特殊情形中的 $\mathcal{F}'_1$）。由注
\ref{defos-remark-trivial-extension-ringed-topoi}，利用注
\ref{defos-remark-short-exact-sequence-thickenings-ringed-topoi} 中的典范分裂
$\pi : (\Sh(\mathcal{D}_1), \mathcal{O}'_1) \to
(\Sh(\mathcal{C}), \mathcal{O})$，得到
$\xi_{\mathcal{F}'_1} \in
\Ext^1_\mathcal{O}(\mathcal{F}, \mathcal{K}_1)$.
最后，还有障碍
$$
o(\mathcal{F}, \mathcal{K}_3, c_3) \in
\Ext^2_\mathcal{O}(\mathcal{F}, \mathcal{K}_3)
$$
见引理 \ref{defos-lemma-inf-obs-ext-ringed-topoi}。在此情形，
我们{\bf 断言}，由短正合列
$0\to\mathcal{K}_3\to\mathcal{K}_2\to\mathcal{K}_1\to0$
给出的典范同态
$$
\partial :
\Ext^1_\mathcal{O}(\mathcal{F}, \mathcal{K}_1)
\longrightarrow
\Ext^2_\mathcal{O}(\mathcal{F}, \mathcal{K}_3)
$$
把 $\xi_{\mathcal{F}'_1}$ 映到障碍类
$o(\mathcal{F},\mathcal{K}_3,c_3)$。

\medskip\noindent
为证明此断言，取嵌入 $j:\mathcal{K}_3\to\mathcal{K}$，
其中 $\mathcal{K}$ 是内射 $\mathcal{O}$-模。可把 $j$ 提升为同态
$j':\mathcal{K}_2\to\mathcal{K}$。令
$\mathcal{E}'_2=j'_*\mathcal{F}'_2$ 为沿 $j'$ 对
$\mathcal{F}'_2$ 取的推出，使
$c_{\mathcal{E}'_2}=j'\circ c_2$。图示如下：
$$
\xymatrix{
0 \ar[r] &
\mathcal{K}_2 \ar[r] \ar[d]_{j'} &
\mathcal{F}'_2 \ar[r] \ar[d] &
\mathcal{F} \ar[r] \ar[d] & 0 \\
0 \ar[r] &
\mathcal{K} \ar[r] &
\mathcal{E}'_2 \ar[r] &
\mathcal{F} \ar[r] & 0
}
$$
令 $\mathcal{E}'_3=\mathcal{E}'_2$，但通过
$\mathcal{O}'_3\to\mathcal{O}'_2$ 把它视为 $\mathcal{O}'_3$-模。
于是 $c_{\mathcal{E}'_3}=j\circ c_3$。引理
\ref{defos-lemma-inf-obs-ext-ringed-topoi} 的证明把
$o(\mathcal{F},\mathcal{K}_3,c_3)$ 构造为下列
$\mathcal{O}$-模扩张之类的边界：
$$
0 \to
\mathcal{K}/\mathcal{K}_3 \to
\mathcal{E}'_3/\mathcal{K}_3 \to
\mathcal{F} \to 0
$$
另一方面，注意 $\mathcal{F}'_1=\mathcal{F}'_2/\mathcal{K}_3$，
故类 $\xi_{\mathcal{F}'_1}$ 是扩张
$$
0 \to \mathcal{K}_2/\mathcal{K}_3 \to \mathcal{F}'_2/\mathcal{K}_3
\to \mathcal{F} \to 0
$$
之类；这里借助 $\pi^\sharp$ 把它视为 $\mathcal{O}$-模列，
其中 $\pi:(\Sh(\mathcal{D}_1),\mathcal{O}'_1)\to
(\Sh(\mathcal{C}),\mathcal{O})$ 是典范分裂。最后，存在交换图
$$
\xymatrix{
0 \ar[r] &
\mathcal{K}_2/\mathcal{K}_3 \ar[r] \ar[d] &
\mathcal{F}'_2/\mathcal{K}_3 \ar[r] \ar[d] &
\mathcal{F} \ar[r] \ar[d] & 0 \\
0 \ar[r] &
\mathcal{K}/\mathcal{K}_3 \ar[r] &
\mathcal{E}'_3/\mathcal{K}_3 \ar[r] &
\mathcal{F} \ar[r] & 0
}
$$
且该图关于上述 $\mathcal{O}$-模结构是 $\mathcal{O}$-线性的，
故断言成立。
\end{remark}







\section{环化拓扑斯上模的无穷小形变}
\label{defos-section-deformation-modules-ringed-topoi}

\noindent
设 $i:(\Sh(\mathcal{C}),\mathcal{O})\to
(\Sh(\mathcal{D}),\mathcal{O}')$ 是环化拓扑斯的一阶增厚。
我们沿用第 \ref{defos-section-thickenings-ringed-topoi} 节引入的记号。
设 $\mathcal{F}'$ 是 $\mathcal{O}'$-模，并令
$\mathcal{F}=i^*\mathcal{F}'$。此时有 $\mathcal{O}'$-模短正合列
$$
0 \to \mathcal{I}\mathcal{F}' \to \mathcal{F}' \to \mathcal{F} \to 0
$$
由于 $\mathcal{I}^2=0$，$\mathcal{I}\mathcal{F}'$ 上的
$\mathcal{O}'$-模结构来自唯一的 $\mathcal{O}$-模结构。
因此上述正合列是形如 (\ref{defos-equation-extension-ringed-topoi}) 的扩张。
特别地，若 $\mathcal{F}'=\mathcal{O}'$，则
$i^*\mathcal{O}'=\mathcal{O}$ 且
$\mathcal{I}\mathcal{O}'=\mathcal{I}$，于是恢复结构层正合列
$$
0 \to \mathcal{I} \to \mathcal{O}' \to \mathcal{O} \to 0
$$

\begin{lemma}
\label{defos-lemma-inf-map-special-ringed-topoi}
设 $i:(\Sh(\mathcal{C}),\mathcal{O})\to
(\Sh(\mathcal{D}),\mathcal{O}')$ 是环化拓扑斯的一阶增厚。
设 $\mathcal{F}'$、$\mathcal{G}'$ 是 $\mathcal{O}'$-模，令
$\mathcal{F}=i^*\mathcal{F}'$、$\mathcal{G}=i^*\mathcal{G}'$。
设 $\varphi:\mathcal{F}\to\mathcal{G}$ 是 $\mathcal{O}$-线性同态。
若把 $\varphi$ 提升为 $\mathcal{O}'$-线性同态
$\varphi':\mathcal{F}'\to\mathcal{G}'$ 的提升集合非空，
则它是 $\Hom_\mathcal{O}(\mathcal{F},\mathcal{I}\mathcal{G}')$
作用下的主齐性空间。
\end{lemma}

\begin{proof}
这是引理 \ref{defos-lemma-inf-map-ringed-topoi} 的特殊情形，
但我们也给出直接证明。有模的短正合列
$$
0 \to \mathcal{I} \to \mathcal{O}' \to \mathcal{O} \to 0
\quad\text{且}\quad
0 \to \mathcal{I}\mathcal{G}' \to \mathcal{G}' \to \mathcal{G} \to 0
$$
$\mathcal{F}'$ 也有类似的正合列。由于 $\mathcal{I}$ 平方为零，
$\mathcal{I}$ 与 $\mathcal{I}\mathcal{G}'$ 上的 $\mathcal{O}'$-模结构
都来自唯一的 $\mathcal{O}$-模结构。因此
$$
\Hom_{\mathcal{O}'}(\mathcal{F}', \mathcal{I}\mathcal{G}') =
\Hom_\mathcal{O}(\mathcal{F}, \mathcal{I}\mathcal{G}')
\quad\text{且}\quad
\Hom_{\mathcal{O}'}(\mathcal{F}', \mathcal{G}) =
\Hom_\mathcal{O}(\mathcal{F}, \mathcal{G})
$$
现在本引理由正合列
$$
0 \to \Hom_{\mathcal{O}'}(\mathcal{F}', \mathcal{I}\mathcal{G}') \to
\Hom_{\mathcal{O}'}(\mathcal{F}', \mathcal{G}') \to
\Hom_{\mathcal{O}'}(\mathcal{F}', \mathcal{G})
$$
得出，见《同调代数》引理 \ref{homology-lemma-check-exactness}。
\end{proof}

\begin{lemma}
\label{defos-lemma-deform-module-ringed-topoi}
设 $(f,f')$ 是情形 \ref{defos-situation-morphism-thickenings-ringed-topoi}
中的环化拓扑斯一阶增厚之间的态射。设 $\mathcal{F}'$ 是
$\mathcal{O}'$-模，并令 $\mathcal{F}=i^*\mathcal{F}'$。
假设 $\mathcal{F}$ 在 $\mathcal{O}_\mathcal{B}$ 上平坦，且
$(f,f')$ 是增厚的严格态射（定义
\ref{defos-definition-strict-morphism-thickenings-ringed-topoi}）。
那么下列条件等价：
\begin{enumerate}
\item $\mathcal{F}'$ 在 $\mathcal{O}_{\mathcal{B}'}$ 上平坦；
\item 典范同态
$f^*\mathcal{J} \otimes_\mathcal{O} \mathcal{F} \to
\mathcal{I}\mathcal{F}'$
是同构。
\end{enumerate}
此外，在此情形，同态
$$
f^*\mathcal{J} \otimes_\mathcal{O} \mathcal{F} \to
\mathcal{I} \otimes_\mathcal{O} \mathcal{F} \to
\mathcal{I}\mathcal{F}'
$$
均为同构。
\end{lemma}

\begin{proof}
由于 $(f,f')$ 是增厚的严格态射，同态
$f^*\mathcal{J}\to\mathcal{I}$ 是满射。因此最后一个断言由 (2) 得出。

\medskip\noindent
(1) 与 (2) 等价的证明。按定义，在 $\mathcal{O}_\mathcal{B}$ 上平坦
意指在 $f^{-1}\mathcal{O}_\mathcal{B}$ 上平坦；
$f^{-1}\mathcal{O}_{\mathcal{B}'}$ 上的平坦性同理。
注意，$(f,f')$ 的严格性及假设
$\mathcal{F}=i^*\mathcal{F}'$ 蕴含
$$
\mathcal{F} = \mathcal{F}'/(f^{-1}\mathcal{J})\mathcal{F}'
$$
（视为 $\mathcal{C}$ 上的层）。此外，注意
$f^*\mathcal{J} \otimes_\mathcal{O} \mathcal{F} =
f^{-1}\mathcal{J} \otimes_{f^{-1}\mathcal{O}_\mathcal{B}} \mathcal{F}$.
因此 (1) 与 (2) 的等价性由《位点上的模》引理
\ref{sites-modules-lemma-flat-over-thickening} 得出。
\end{proof}

\begin{lemma}
\label{defos-lemma-deform-fp-module-ringed-topoi}
设 $(f,f')$ 是情形 \ref{defos-situation-morphism-thickenings-ringed-topoi}
中的环化拓扑斯一阶增厚之间的态射。设 $\mathcal{F}'$ 是
$\mathcal{O}'$-模，并令 $\mathcal{F}=i^*\mathcal{F}'$。
假设 $\mathcal{F}'$ 在 $\mathcal{O}_{\mathcal{B}'}$ 上平坦，
且 $(f,f')$ 是增厚的严格态射。那么下列条件等价：
\begin{enumerate}
\item $\mathcal{F}'$ 是有限表示 $\mathcal{O}'$-模；
\item $\mathcal{F}$ 是有限表示 $\mathcal{O}$-模。
\end{enumerate}
\end{lemma}

\begin{proof}
(1) $\Rightarrow$ (2) 由《位点上的模》引理
\ref{sites-modules-lemma-local-pullback} 得出。
反过来，假设 $\mathcal{F}$ 有限表示。可以并且确实假设
$\mathcal{C}=\mathcal{C}'$。由引理
\ref{defos-lemma-deform-module-ringed-topoi}，有短正合列
% P11-E0007: undefined O_X subscript retained from frozen source; see control/ERRATA_CANDIDATES.jsonl.
$$
0 \to \mathcal{I} \otimes_{\mathcal{O}_X} \mathcal{F} \to
\mathcal{F}' \to \mathcal{F} \to 0
$$
取 $\mathcal{C}$ 的对象 $U$，使 $\mathcal{F}|_U$ 有表示
$$
\mathcal{O}_U^{\oplus m} \to \mathcal{O}_U^{\oplus n} \to \mathcal{F}|_U \to 0
$$
把 $U$ 换成某一覆盖的各成员后，可以假设同态
$\mathcal{O}_U^{\oplus n}\to\mathcal{F}|_U$ 提升为同态
$(\mathcal{O}'_U)^{\oplus n}\to\mathcal{F}'|_U$。
由 $\otimes$ 的右正合性，所诱导的同态
$\mathcal{I}^{\oplus n}\to\mathcal{I}\otimes\mathcal{F}$ 是满射。
因此再次把 $U$ 换成某一覆盖的各成员后，可找到给定同态
$\mathcal{O}_U^{\oplus m}\to\mathcal{O}_U^{\oplus n}$ 的提升
$(\mathcal{O}'|_U)^{\oplus m} \to (\mathcal{O}'|_U)^{\oplus n}$
使得
$$
(\mathcal{O}'_U)^{\oplus m} \to (\mathcal{O}'_U)^{\oplus n} \to
\mathcal{F}'|_U \to 0
$$
构成复形。再次利用 $\otimes$ 的右正合性可见此复形正合。
\end{proof}

\begin{lemma}
\label{defos-lemma-inf-map-rel-ringed-topoi}
设 $(f,f')$ 是情形 \ref{defos-situation-morphism-thickenings-ringed-topoi}
中的一阶增厚态射。设 $\mathcal{F}'$、$\mathcal{G}'$ 是
$\mathcal{O}'$-模，并令
$\mathcal{F}=i^*\mathcal{F}'$、$\mathcal{G}=i^*\mathcal{G}'$。
设 $\varphi:\mathcal{F}\to\mathcal{G}$ 是 $\mathcal{O}$-线性同态。
假设 $\mathcal{G}'$ 在 $\mathcal{O}_{\mathcal{B}'}$ 上平坦，
且 $(f,f')$ 是增厚的严格态射。若把 $\varphi$ 提升为
$\mathcal{O}'$-线性同态
$\varphi':\mathcal{F}'\to\mathcal{G}'$ 的提升集合非空，
则它是下列群作用下的主齐性空间：
$$
\Hom_\mathcal{O}(\mathcal{F},
\mathcal{G} \otimes_\mathcal{O} f^*\mathcal{J})
$$
\end{lemma}

\begin{proof}
合用引理 \ref{defos-lemma-inf-map-special-ringed-topoi} 与
\ref{defos-lemma-deform-module-ringed-topoi} 即得。
\end{proof}

\begin{lemma}
\label{defos-lemma-inf-obs-map-special-ringed-topoi}
设 $i:(\Sh(\mathcal{C}),\mathcal{O})\to
(\Sh(\mathcal{D}),\mathcal{O}')$ 是环化拓扑斯的一阶增厚。
设 $\mathcal{F}'$、$\mathcal{G}'$ 是 $\mathcal{O}'$-模，并令
$\mathcal{F}=i^*\mathcal{F}'$、$\mathcal{G}=i^*\mathcal{G}'$。
设 $\varphi:\mathcal{F}\to\mathcal{G}$ 是 $\mathcal{O}$-线性同态。
存在元素
$$
o(\varphi) \in
\Ext^1_\mathcal{O}(Li^*\mathcal{F}', \mathcal{I}\mathcal{G}')
$$
其消失是存在把 $\varphi$ 提升为 $\mathcal{O}'$-线性同态
$\varphi':\mathcal{F}'\to\mathcal{G}'$ 的提升的充要条件。
\end{lemma}

\begin{proof}
由引理 \ref{defos-lemma-inf-map-special-ringed-topoi} 的证明可知，
$\varphi$ 在同态
$$
\Hom_\mathcal{O}(\mathcal{F}, \mathcal{G}) =
\Hom_{\mathcal{O}'}(\mathcal{F}', \mathcal{G}) \longrightarrow
\Ext^1_{\mathcal{O}'}(\mathcal{F}', \mathcal{I}\mathcal{G}')
$$
下的边界消失，正是存在提升的充要条件。由导出范畴上
$i_*=Ri_*$ 与 $Li^*$ 的伴随性（《位点上的上同调》引理
\ref{sites-cohomology-lemma-adjoint}），有
$$
\Ext^1_{\mathcal{O}'}(\mathcal{F}', \mathcal{I}\mathcal{G}') =
\Ext^1_\mathcal{O}(Li^*\mathcal{F}', \mathcal{I}\mathcal{G}')
$$
结论随即得出。
\end{proof}

\begin{lemma}
\label{defos-lemma-inf-obs-map-rel-ringed-topoi}
设 $(f,f')$ 是情形 \ref{defos-situation-morphism-thickenings-ringed-topoi}
中的一阶增厚态射。设 $\mathcal{F}'$、$\mathcal{G}'$ 是
$\mathcal{O}'$-模，并令
$\mathcal{F}=i^*\mathcal{F}'$、$\mathcal{G}=i^*\mathcal{G}'$。
设 $\varphi:\mathcal{F}\to\mathcal{G}$ 是 $\mathcal{O}$-线性同态。
假设 $\mathcal{F}'$、$\mathcal{G}'$ 在
$\mathcal{O}_{\mathcal{B}'}$ 上平坦，且 $(f,f')$ 是增厚的严格态射。
存在元素
$$
o(\varphi) \in
\Ext^1_\mathcal{O}(\mathcal{F},
\mathcal{G} \otimes_\mathcal{O} f^*\mathcal{J})
$$
其消失是存在把 $\varphi$ 提升为 $\mathcal{O}'$-线性同态
$\varphi':\mathcal{F}'\to\mathcal{G}'$ 的提升的充要条件。
\end{lemma}

\begin{proof}[第一种证明]
这由引理 \ref{defos-lemma-inf-obs-map-special-ringed-topoi} 得出，
因为我们断言，在本引理的假设下有
$$
\Ext^1_\mathcal{O}(Li^*\mathcal{F}', \mathcal{I}\mathcal{G}') =
\Ext^1_\mathcal{O}(\mathcal{F},
\mathcal{G} \otimes_\mathcal{O} f^*\mathcal{J})
$$
事实上，由引理 \ref{defos-lemma-deform-module-ringed-topoi}，有
$\mathcal{I}\mathcal{G}' =
\mathcal{G}\otimes_\mathcal{O}f^*\mathcal{J}$。
另一方面，注意
$$
H^{-1}(Li^*\mathcal{F}') =
\text{Tor}_1^{\mathcal{O}'}(\mathcal{F}', \mathcal{O})
$$
（局部计算从略）。利用短正合列
$$
0 \to \mathcal{I} \to \mathcal{O}' \to \mathcal{O} \to 0
$$
可见，此 $\text{Tor}_1$ 由同态
$\mathcal{I} \otimes_\mathcal{O} \mathcal{F} \to \mathcal{I}\mathcal{F}'$
的核计算；由引理 \ref{defos-lemma-deform-module-ringed-topoi}
的最后一个断言，该核为零。因此
$\tau_{\geq-1}Li^*\mathcal{F}'=\mathcal{F}$。另一方面，
由《导出范畴》引理 \ref{derived-lemma-negative-vanishing} 的对偶形式，有
$$
\Ext^1_\mathcal{O}(Li^*\mathcal{F}',
\mathcal{I}\mathcal{G}') =
\Ext^1_\mathcal{O}(\tau_{\geq -1}Li^*\mathcal{F}',
\mathcal{I}\mathcal{G}')
$$
\end{proof}

\begin{proof}[第二种证明]
可如下应用引理 \ref{defos-lemma-inf-obs-map-ringed-topoi}。注意，由引理
\ref{defos-lemma-deform-module-ringed-topoi}，有
$\mathcal{K}=\mathcal{I}\otimes_\mathcal{O}\mathcal{F}$、
$\mathcal{L} = \mathcal{I} \otimes_\mathcal{O} \mathcal{G}$
且 $c_{\mathcal{F}'}=1\otimes1$、
$c_{\mathcal{G}'}=1\otimes1$。取 $\psi=1\otimes\varphi$，
则该引理中的图交换。因此取
$o(\varphi)=o(\varphi,1\otimes\varphi)$ 即可。
\end{proof}

\begin{lemma}
\label{defos-lemma-inf-ext-rel-ringed-topoi}
设 $(f,f')$ 是情形 \ref{defos-situation-morphism-thickenings-ringed-topoi}
中的一阶增厚态射。设 $\mathcal{F}$ 是 $\mathcal{O}$-模。
假设 $(f,f')$ 是增厚的严格态射，且 $\mathcal{F}$ 在
$\mathcal{O}_\mathcal{B}$ 上平坦。若存在二元组
$(\mathcal{F}',\alpha)$，其中 $\mathcal{F}'$ 是在
$\mathcal{O}_{\mathcal{B}'}$ 上平坦的 $\mathcal{O}'$-模，且
$\alpha:i^*\mathcal{F}'\to\mathcal{F}$ 是同构，
则这类二元组的同构类集合是下列群作用下的主齐性空间：
$\Ext^1_\mathcal{O}(
\mathcal{F}, \mathcal{I} \otimes_\mathcal{O} \mathcal{F})$.
\end{lemma}

\begin{proof}
若假设存在一个这样的模，则由引理
\ref{defos-lemma-deform-module-ringed-topoi}，典范同态
$$
f^*\mathcal{J} \otimes_\mathcal{O} \mathcal{F} \to
\mathcal{I} \otimes_\mathcal{O} \mathcal{F}
$$
是同构。应用引理 \ref{defos-lemma-inf-ext-ringed-topoi}，其中
$\mathcal{K}=\mathcal{I}\otimes_\mathcal{O}\mathcal{F}$ 且 $c=1$。
由引理 \ref{defos-lemma-deform-module-ringed-topoi}，相应扩张
$\mathcal{F}'$ 均在 $\mathcal{O}_{\mathcal{B}'}$ 上平坦。
\end{proof}

\begin{lemma}
\label{defos-lemma-inf-obs-ext-rel-ringed-topoi}
设 $(f,f')$ 是情形 \ref{defos-situation-morphism-thickenings-ringed-topoi}
中的一阶增厚态射。设 $\mathcal{F}$ 是 $\mathcal{O}$-模。
假设 $(f,f')$ 是增厚的严格态射，且 $\mathcal{F}$ 在
$\mathcal{O}_\mathcal{B}$ 上平坦。存在在
$\mathcal{O}_{\mathcal{B}'}$ 上平坦的 $\mathcal{O}'$-模
$\mathcal{F}'$ 满足 $i^*\mathcal{F}'\cong\mathcal{F}$，
当且仅当
\begin{enumerate}
\item 典范同态
$f^*\mathcal{J} \otimes_\mathcal{O} \mathcal{F} \to
\mathcal{I} \otimes_\mathcal{O} \mathcal{F}$
是同构；并且
\item 引理 \ref{defos-lemma-inf-obs-ext-ringed-topoi} 中的类
$o(\mathcal{F}, \mathcal{I} \otimes_\mathcal{O} \mathcal{F}, 1)
\in \Ext^2_\mathcal{O}(
\mathcal{F}, \mathcal{I} \otimes_\mathcal{O} \mathcal{F})$
为零。
\end{enumerate}
\end{lemma}

\begin{proof}
这立即由引理 \ref{defos-lemma-deform-module-ringed-topoi}
对在 $\mathcal{O}_{\mathcal{B}'}$ 上平坦的 $\mathcal{O}'$-模的刻画，
以及引理 \ref{defos-lemma-inf-obs-ext-ringed-topoi} 得出。
\end{proof}






\section{在环化拓扑斯平坦增厚上的平坦模中的应用}
\label{defos-section-flat-ringed-topoi}

\noindent
考虑环化拓扑斯的交换图
$$
\xymatrix{
(\Sh(\mathcal{C}), \mathcal{O}) \ar[r]_i \ar[d]_f &
(\Sh(\mathcal{D}), \mathcal{O}') \ar[d]^{f'} \\
(\Sh(\mathcal{B}), \mathcal{O}_\mathcal{B}) \ar[r]^t &
(\Sh(\mathcal{B}'), \mathcal{O}_{\mathcal{B}'})
}
$$
其水平箭头是情形 \ref{defos-situation-morphism-thickenings-ringed-topoi}
中的一阶增厚。令
$\mathcal{I}=\Ker(i^\sharp)\subset\mathcal{O}'$，并令
$\mathcal{J}=\Ker(t^\sharp)\subset\mathcal{O}_{\mathcal{B}'}$。
设 $\mathcal{F}$ 是 $\mathcal{O}$-模。假设
\begin{enumerate}
\item $(f,f')$ 是增厚的严格态射；
\item $f'$ 平坦；
\item $\mathcal{F}$ 在 $\mathcal{O}_\mathcal{B}$ 上平坦。
\end{enumerate}
注意，(1) $+$ (2) 蕴含 $\mathcal{I}=f^*\mathcal{J}$
（把引理 \ref{defos-lemma-deform-module-ringed-topoi} 应用于
$\mathcal{O}'$）。在这些假设下，上一节的理论格外简洁。
下列引理汇总已经得到的结果。

\begin{lemma}
\label{defos-lemma-flat-ringed-topoi}
在上述情形中：
\begin{enumerate}
\item 存在在 $\mathcal{O}_{\mathcal{B}'}$ 上平坦且满足
$i^*\mathcal{F}'\cong\mathcal{F}$ 的 $\mathcal{O}'$-模
$\mathcal{F}'$，当且仅当引理
\ref{defos-lemma-inf-obs-ext-ringed-topoi} 中的类
$o(\mathcal{F}, f^*\mathcal{J} \otimes_\mathcal{O} \mathcal{F}, 1)
\in \Ext^2_\mathcal{O}(
\mathcal{F}, f^*\mathcal{J} \otimes_\mathcal{O} \mathcal{F})$
为零。
\item 若这样的模存在，则提升的同构类集合是下列群作用下的主齐性空间：
$\Ext^1_\mathcal{O}(
\mathcal{F}, f^*\mathcal{J} \otimes_\mathcal{O} \mathcal{F})$.
\item 给定提升 $\mathcal{F}'$，$\mathcal{F}'$ 的所有拉回为
$\text{id}_\mathcal{F}$ 的自同构所成集合典范同构于
$\Ext^0_\mathcal{O}(
\mathcal{F}, f^*\mathcal{J} \otimes_\mathcal{O} \mathcal{F})$.
\end{enumerate}
\end{lemma}

\begin{proof}
由上面已见的 $\mathcal{I}=f^*\mathcal{J}$，(1) 由引理
\ref{defos-lemma-inf-obs-ext-rel-ringed-topoi} 得出。
(2) 由引理 \ref{defos-lemma-inf-ext-rel-ringed-topoi} 得出。
(3) 由引理 \ref{defos-lemma-inf-map-rel-ringed-topoi} 得出。
\end{proof}

\begin{situation}
\label{defos-situation-morphism-flat-thickenings-ringed-topoi}
设 $f:(\Sh(\mathcal{C}),\mathcal{O})\to
(\Sh(\mathcal{B}),\mathcal{O}_\mathcal{B})$ 是环化拓扑斯态射。
考虑交换图
$$
\xymatrix{
(\Sh(\mathcal{C}'_1), \mathcal{O}'_1) \ar[r]_h \ar[d]_{f'_1} &
(\Sh(\mathcal{C}'_2), \mathcal{O}'_2) \ar[d]_{f'_2} \\
(\Sh(\mathcal{B}'_1), \mathcal{O}_{\mathcal{B}'_1}) \ar[r] &
(\Sh(\mathcal{B}'_2), \mathcal{O}_{\mathcal{B}'_2})
}
$$
其中 $h$ 是 $(\Sh(\mathcal{C}),\mathcal{O})$ 的一阶增厚之间的态射，
下方水平箭头是 $(\Sh(\mathcal{B}),\mathcal{O}_\mathcal{B})$
的一阶增厚之间的态射，每个 $f'_i$ 都限制为 $f$，
两对 $(f,f_i')$ 都是增厚的严格态射，且两个 $f'_i$ 都平坦。
最后，设 $\mathcal{F}$ 是在 $\mathcal{O}_\mathcal{B}$ 上平坦的
$\mathcal{O}$-模。
\end{situation}

\begin{lemma}
\label{defos-lemma-functorial-ringed-topoi}
在情形 \ref{defos-situation-morphism-flat-thickenings-ringed-topoi} 中，
障碍类
$o(\mathcal{F}, f^*\mathcal{J}_2 \otimes_\mathcal{O} \mathcal{F}, 1)$
在典范同态
$$
\Ext^2_\mathcal{O}(
\mathcal{F}, f^*\mathcal{J}_2 \otimes_\mathcal{O} \mathcal{F})
\to \Ext^2_\mathcal{O}(
\mathcal{F}, f^*\mathcal{J}_1 \otimes_\mathcal{O} \mathcal{F})
$$
下映到障碍类
$o(\mathcal{F}, f^*\mathcal{J}_1 \otimes_\mathcal{O} \mathcal{F}, 1)$
。
\end{lemma}

\begin{proof}
由注 \ref{defos-remark-obstruction-extension-functorial-ringed-topoi} 得出。
\end{proof}

\begin{situation}
\label{defos-situation-ses-flat-thickenings-ringed-topoi}
设 $f:(\Sh(\mathcal{C}),\mathcal{O})\to
(\Sh(\mathcal{B}),\mathcal{O}_\mathcal{B})$ 是环化拓扑斯态射。
考虑交换图
$$
\xymatrix{
(\Sh(\mathcal{C}'_1), \mathcal{O}'_1) \ar[r]_h \ar[d]_{f'_1} &
(\Sh(\mathcal{C}'_2), \mathcal{O}'_2) \ar[r] \ar[d]_{f'_2} &
(\Sh(\mathcal{C}'_3), \mathcal{O}'_3) \ar[d]_{f'_3} \\
(\Sh(\mathcal{B}'_1), \mathcal{O}_{\mathcal{B}'_1}) \ar[r] &
(\Sh(\mathcal{B}'_2), \mathcal{O}_{\mathcal{B}'_2}) \ar[r] &
(\Sh(\mathcal{B}'_3), \mathcal{O}_{\mathcal{B}'_3})
}
$$
其中：(a) 上行是 $(\Sh(\mathcal{C}),\mathcal{O})$
的一阶增厚短正合列；(b) 下行是
$(\Sh(\mathcal{B}),\mathcal{O}_\mathcal{B})$ 的一阶增厚短正合列；
(c) 每个 $f'_i$ 都限制为 $f$；(d) 每一对 $(f,f_i')$
都是增厚的严格态射；(e) 每个 $f'_i$ 都平坦。最后，设
$\mathcal{F}'_2$ 是在 $\mathcal{O}_{\mathcal{B}'_2}$ 上平坦的
$\mathcal{O}'_2$-模，并令
$\mathcal{F}=\mathcal{F}'_2\otimes\mathcal{O}$。令
$\pi:(\Sh(\mathcal{C}'_1),\mathcal{O}'_1)\to
(\Sh(\mathcal{C}),\mathcal{O})$ 为典范分裂（注
\ref{defos-remark-short-exact-sequence-thickenings-ringed-topoi}）。
\end{situation}

\begin{lemma}
\label{defos-lemma-verify-iv-ringed-topoi}
在情形 \ref{defos-situation-ses-flat-thickenings-ringed-topoi} 中，
$\pi^*\mathcal{F}$ 与 $h^*\mathcal{F}'_2$ 是在
$\mathcal{O}_{\mathcal{B}'_1}$ 上平坦的 $\mathcal{O}'_1$-模，
且在 $(\Sh(\mathcal{C}),\mathcal{O})$ 上都限制为 $\mathcal{F}$。
二者之差（引理 \ref{defos-lemma-flat-ringed-topoi}）是
$\Ext^1_\mathcal{O}(\mathcal{F},
f^*\mathcal{J}_1 \otimes_\mathcal{O} \mathcal{F})$
中的元素 $\theta$，它在
$\Ext^2_\mathcal{O}(\mathcal{F},
f^*\mathcal{J}_3 \otimes_\mathcal{O} \mathcal{F})$
中的边界，等于把 $\mathcal{F}$ 提升为在
$\mathcal{O}_{\mathcal{B}'_3}$ 上平坦的 $\mathcal{O}'_3$-模的障碍
（引理 \ref{defos-lemma-flat-ringed-topoi}）。
\end{lemma}

\begin{proof}
注意，$\pi^*\mathcal{F}$ 与 $h^*\mathcal{F}'_2$ 在
$(\Sh(\mathcal{C}),\mathcal{O})$ 上都限制为 $\mathcal{F}$，
而 $\pi^*\mathcal{F}\to\mathcal{F}$ 与
$h^*\mathcal{F}'_2\to\mathcal{F}$ 的核均由
$f^*\mathcal{J}_1\otimes_\mathcal{O}\mathcal{F}$ 给出。
因此由引理 \ref{defos-lemma-deform-module-ringed-topoi} 可得平坦性。
取边界是有意义的，因为模列
$$
0 \to f^*\mathcal{J}_3 \otimes_\mathcal{O} \mathcal{F} \to
f^*\mathcal{J}_2 \otimes_\mathcal{O} \mathcal{F} \to
f^*\mathcal{J}_1 \otimes_\mathcal{O} \mathcal{F} \to 0
$$
由情形 \ref{defos-situation-ses-flat-thickenings-ringed-topoi} 的假设及
$\mathcal{F}$ 在 $\mathcal{O}_\mathcal{B}$ 上平坦这一事实而短正合。
关于障碍类的断言，正是把注
\ref{defos-remark-complex-thickenings-and-ses-modules-ringed-topoi}
的结果直接转译到这一特定情形。
\end{proof}







\section{环化拓扑斯的形变与朴素余切复形}
\label{defos-section-deformations-ringed-topoi}

\noindent
本节利用朴素余切复形研究一些形变理论。先取环化拓扑斯的一阶增厚
$t : (\Sh(\mathcal{B}), \mathcal{O}_\mathcal{B}) \to
(\Sh(\mathcal{B}'), \mathcal{O}_{\mathcal{B}'})$。
记 $\mathcal{J}=\Ker(t^\sharp)$，并等同 $\mathcal{B}$ 与
$\mathcal{B}'$ 的底层拓扑斯。再设给定环化拓扑斯态射
$f : (\Sh(\mathcal{C}), \mathcal{O}) \to
(\Sh(\mathcal{B}), \mathcal{O}_\mathcal{B})$、一个 $\mathcal{O}$-模
$\mathcal{G}$，以及 $f^{-1}\mathcal{O}_\mathcal{B}$-模层同态
$f^{-1}\mathcal{J}\to\mathcal{G}$。本节要问：能否找到下图问号处的对象
\begin{equation}
\label{defos-equation-to-solve-ringed-topoi}
\vcenter{
\xymatrix{
0 \ar[r] & \mathcal{G} \ar[r] & {?} \ar[r] & \mathcal{O} \ar[r] & 0 \\
0 \ar[r] & f^{-1}\mathcal{J} \ar[u]^c \ar[r] &
f^{-1}\mathcal{O}_{\mathcal{B}'} \ar[u] \ar[r] &
f^{-1}\mathcal{O}_\mathcal{B} \ar[u] \ar[r] & 0
}
}
\end{equation}
以及解（若存在）在何种意义下唯一。更确切地说，我们寻找一阶增厚
$i : (\Sh(\mathcal{C}), \mathcal{O}) \to (\Sh(\mathcal{C}'), \mathcal{O}')$
以及形如 (\ref{defos-equation-morphism-thickenings-ringed-topoi}) 的增厚态射
$(f,f')$，其中把 $\Ker(i^\sharp)$ 等同于 $\mathcal{G}$，
且 $(f')^\sharp$ 诱导给定同态 $c$。我们称
$(\Sh(\mathcal{C}'),\mathcal{O}')$ 是
(\ref{defos-equation-to-solve-ringed-topoi}) 的一个{\it 解}。


\begin{lemma}
\label{defos-lemma-huge-diagram-ringed-topoi}
设给定环化拓扑斯态射的交换图
\begin{equation}
\label{defos-equation-huge-1-ringed-topoi}
\vcenter{
\xymatrix{
& (\Sh(\mathcal{C}_2), \mathcal{O}_2) \ar[r]_{i_2} \ar[d]_{f_2} \ar[ddl]_g &
(\Sh(\mathcal{C}'_2), \mathcal{O}'_2) \ar[d]^{f'_2} \\
&
(\Sh(\mathcal{B}_2), \mathcal{O}_{\mathcal{B}_2}) \ar[r]^{t_2} \ar[ddl]|\hole &
(\Sh(\mathcal{B}'_2), \mathcal{O}_{\mathcal{B}'_2}) \ar[ddl] \\
(\Sh(\mathcal{C}_1), \mathcal{O}_1) \ar[r]_{i_1} \ar[d]_{f_1} &
(\Sh(\mathcal{C}'_1), \mathcal{O}'_1) \ar[d]^{f'_1} \\
(\Sh(\mathcal{B}_1), \mathcal{O}_{\mathcal{B}_1}) \ar[r]^{t_1} &
(\Sh(\mathcal{B}'_1), \mathcal{O}_{\mathcal{B}'_1})
}
}
\end{equation}
其水平箭头均为一阶增厚。令 $\mathcal{G}_j=\Ker(i_j^\sharp)$，
并设给定 $g^{-1}\mathcal{O}_1$-模同态
$\nu : g^{-1}\mathcal{G}_1 \to \mathcal{G}_2$
使得下图交换：
\begin{equation}
\label{defos-equation-huge-2-ringed-topoi}
\vcenter{
\xymatrix{
& 0 \ar[r] & \mathcal{G}_2 \ar[r] &
\mathcal{O}'_2 \ar[r] &
\mathcal{O}_2 \ar[r] & 0 \\
& 0 \ar[r]|\hole &
f_2^{-1}\mathcal{J}_2 \ar[u]_{c_2} \ar[r] &
f_2^{-1}\mathcal{O}_{\mathcal{B}'_2} \ar[u] \ar[r]|\hole &
f_2^{-1}\mathcal{O}_{\mathcal{B}_2} \ar[u] \ar[r] & 0 \\
0 \ar[r] &
\mathcal{G}_1 \ar[ruu] \ar[r] &
\mathcal{O}'_1 \ar[r] &
\mathcal{O}_1 \ar[ruu] \ar[r] & 0 \\
0 \ar[r] &
f_1^{-1}\mathcal{J}_1 \ar[ruu]|\hole \ar[u]^{c_1} \ar[r] &
f_1^{-1}\mathcal{O}_{\mathcal{B}'_1} \ar[ruu]|\hole \ar[u] \ar[r] &
f_1^{-1}\mathcal{O}_{\mathcal{B}_1} \ar[ruu]|\hole \ar[u] \ar[r] & 0
}
}
\end{equation}
其前后两个面都是 (\ref{defos-equation-to-solve-ringed-topoi}) 的解。
（西北偏北方向的箭头，是先对源施行 $g^{-1}$ 后得到的
$\mathcal{C}_2$ 上的同态。）
\begin{enumerate}
\item 存在典范元素
$\Ext^1_{\mathcal{O}_2}(
Lg^*\NL_{\mathcal{O}_1/\mathcal{O}_{\mathcal{B}_1}}, \mathcal{G}_2)$
其消失是存在环化拓扑斯态射
$(\Sh(\mathcal{C}'_2), \mathcal{O}'_2) \to
(\Sh(\mathcal{C}'_1), \mathcal{O}'_1)$ 使图
(\ref{defos-equation-huge-1-ringed-topoi})，并与 $\nu$ 相容的充要条件。
\item 若存在态射
$(\Sh(\mathcal{C}'_2), \mathcal{O}'_2) \to
(\Sh(\mathcal{C}'_1), \mathcal{O}'_1)$
使图 (\ref{defos-equation-huge-1-ringed-topoi}) 成立并与 $\nu$ 相容，
则所有这类态射所成集合是下列群作用下的主齐性空间：
$$
\Hom_{\mathcal{O}_1}(
\Omega_{\mathcal{O}_1/\mathcal{O}_{\mathcal{B}_1}}, g_*\mathcal{G}_2) =
\Hom_{\mathcal{O}_2}(
g^*\Omega_{\mathcal{O}_1/\mathcal{O}_{\mathcal{B}_1}}, \mathcal{G}_2) =
\Ext^0_{\mathcal{O}_2}(
Lg^*\NL_{\mathcal{O}_1/\mathcal{O}_{\mathcal{B}_1}}, \mathcal{G}_2).
$$
\end{enumerate}
\end{lemma}

\begin{proof}
本引理的证明与引理 \ref{defos-lemma-huge-diagram-ringed-spaces} 的证明相同。
我们建议读者阅读后者而非此处证明。我们将等同所涉每个增厚的底层拓扑斯
（陈述中已经采用这一约定）。引理最后一个断言中的等式由定义立即得出。
因此在余下的证明中，我们使用群
$\Ext^k_{\mathcal{O}_2}(
Lg^*\NL_{\mathcal{O}_1/\mathcal{O}_{\mathcal{B}_1}}, \mathcal{G}_2)$,
$k=0,1$。我们先说明，可归约到本引理所有环化拓扑斯的底层拓扑斯相同的情形。

\medskip\noindent
为此，注意
$g^{-1}\NL_{\mathcal{O}_1/\mathcal{O}_{\mathcal{B}_1}}$ 等于环层同态
$g^{-1}f_1^{-1}\mathcal{O}_{\mathcal{B}_1} \to g^{-1}\mathcal{O}_1$
的朴素余切复形，见《位点上的模》引理
\ref{sites-modules-lemma-pullback-differentials}。此外，
$\NL_{\mathcal{O}_1/\mathcal{O}_{\mathcal{B}_1}}$
的 $0$ 次项是平坦 $\mathcal{O}_1$-模，故典范同态
$$
Lg^*\NL_{\mathcal{O}_1/\mathcal{O}_{\mathcal{B}_1}}
\longrightarrow
g^{-1}\NL_{\mathcal{O}_1/\mathcal{O}_{\mathcal{B}_1}}
\otimes_{g^{-1}\mathcal{O}_1} \mathcal{O}_2
$$
在 $0$ 次与 $-1$ 次上同调层上诱导同构。因此可把引理中的
$\Ext$ 群替换为
$$
\Ext^k_{g^{-1}\mathcal{O}_1}(
g^{-1}\NL_{\mathcal{O}_1/\mathcal{O}_{\mathcal{B}_1}}, \mathcal{G}_2) =
\Ext^k_{g^{-1}\mathcal{O}_1}(
\NL_{g^{-1}\mathcal{O}_1/g^{-1}f_1^{-1}\mathcal{O}_{\mathcal{B}_1}},
\mathcal{G}_2)
$$
使图 (\ref{defos-equation-huge-1-ringed-topoi}) 成立并与 $\nu$ 相容的
环化拓扑斯态射
$(\Sh(\mathcal{C}'_2), \mathcal{O}'_2) \to
(\Sh(\mathcal{C}'_1), \mathcal{O}'_1)$ 所成集合，与所有同
$f^\sharp$、$\nu$ 相容的
$g^{-1}f_1^{-1}\mathcal{O}_{\mathcal{B}'_1}$-algebras
$g^{-1}\mathcal{O}'_1\to\mathcal{O}'_2$ 的代数同态集合一一对应。
由此可假设我们有位点 $\mathcal{C}$ 上层的图
(\ref{defos-equation-huge-2-ringed-topoi})（在底层拓扑斯上
$f_1=f_2=\text{id}$），并要寻找使该图成立的环层同态
$\mathcal{O}'_1\to\mathcal{O}'_2$。

\medskip\noindent
在本引理余下的证明中，假设所有底层拓扑空间相同；即有位点
$\mathcal{C}$ 上层的图 (\ref{defos-equation-huge-2-ringed-topoi})
（在底层拓扑斯上 $f_1=f_2=\text{id}$），并寻找使该图成立的环层同态
$\mathcal{O}'_1\to\mathcal{O}'_2$。我们使用的 $\Ext$ 群是
$\Ext^k_{\mathcal{O}_1}(
\NL_{\mathcal{O}_1/\mathcal{O}_{\mathcal{B}_1}}, \mathcal{G}_2)$, $k = 0, 1$.

\medskip\noindent
步骤 1：构造障碍类。考虑集合层
$$
\mathcal{E} = \mathcal{O}'_1 \times_{\mathcal{O}_2} \mathcal{O}'_2
$$
它带有满射 $\alpha:\mathcal{E}\to\mathcal{O}_1$，
故可用 $\NL(\alpha)$ 代替
$\NL_{\mathcal{O}_1/\mathcal{O}_{\mathcal{B}_1}}$，见
《位点上的模》引理 \ref{sites-modules-lemma-NL-up-to-qis}。令
$$
\mathcal{I}' =
\Ker(\mathcal{O}_{\mathcal{B}'_1}[\mathcal{E}] \to \mathcal{O}_1)
\quad\text{且}\quad
\mathcal{I} =
\Ker(\mathcal{O}_{\mathcal{B}_1}[\mathcal{E}] \to \mathcal{O}_1)
$$
存在满射 $\mathcal{I}'\to\mathcal{I}$，其核为
$\mathcal{J}_1\mathcal{O}_{\mathcal{B}'_1}[\mathcal{E}]$。
得到两个 $\mathcal{O}_{\mathcal{B}'_2}$-代数同态
$$
a : \mathcal{O}_{\mathcal{B}'_1}[\mathcal{E}] \to \mathcal{O}'_1
\quad\text{且}\quad
b : \mathcal{O}_{\mathcal{B}'_1}[\mathcal{E}] \to \mathcal{O}'_2
$$
它们诱导同态 $a|_{\mathcal{I}'}:\mathcal{I}'\to\mathcal{G}_1$ 与
$b|_{\mathcal{I}'}:\mathcal{I}'\to\mathcal{G}_2$。$a$、$b$ 均零化
$(\mathcal{I}')^2$。此外，由于
(\ref{defos-equation-huge-2-ringed-topoi}) 的左侧方块交换，
$a$、$b$ 作为取值于 $\mathcal{G}_2$ 的同态，在
$\mathcal{J}_1\mathcal{O}_{\mathcal{B}'_1}[\mathcal{E}]$
上一致。因此差
$b|_{\mathcal{I}'} - \nu \circ a|_{\mathcal{I}'}$
诱导良定义的 $\mathcal{O}_1$-线性同态
$$
\xi : \mathcal{I}/\mathcal{I}^2 \longrightarrow \mathcal{G}_2
$$
% P11-E0008: ill-typed reversed expression retained from frozen source; see control/ERRATA_CANDIDATES.jsonl.
它把 $\mathcal{I}$ 的局部截面 $f$ 之类映到
$a(f') - \nu(b(f'))$，
其中 $f'$ 是 $f$ 到 $\mathcal{I}'$ 的局部截面的一个提升。令
$[\xi] \in \Ext^1_{\mathcal{O}_1}(\NL(\alpha), \mathcal{G}_2)$
为其像（见下文）。

\medskip\noindent
步骤 2：$[\xi]$ 消失是必要的。记 $\Omega=
\Omega_{\mathcal{O}_{\mathcal{B}_1}[\mathcal{E}]/\mathcal{O}_{\mathcal{B}_1}}
\otimes_{\mathcal{O}_{\mathcal{B}_1}[\mathcal{E}]} \mathcal{O}_1$.
注意，$\NL(\alpha)=(\mathcal{I}/\mathcal{I}^2\to\Omega)$
嵌入可区别三角
$$
\Omega[0] \to
\NL(\alpha) \to
\mathcal{I}/\mathcal{I}^2[1] \to
\Omega[1]
$$
因此 $[\xi]$ 为零，当且仅当对于某个同态
$\Omega\to\mathcal{G}_2$，$\xi$ 是复合
$\mathcal{I}/\mathcal{I}^2\to\Omega\to\mathcal{G}_2$。
假设存在使图 (\ref{defos-equation-huge-2-ringed-topoi}) 成立的环层同态
$\varphi : \mathcal{O}'_1 \to \mathcal{O}'_2$
。在此情形，考虑同态
$\mathcal{O}'_1[\mathcal{E}] \to \mathcal{G}_2$,
$f'\mapsto b(f')-\varphi(a(f'))$。计算表明，它零化
$\mathcal{J}_1\mathcal{O}_{\mathcal{B}'_1}[\mathcal{E}]$，
并诱导导子 $\mathcal{O}_{\mathcal{B}_1}[\mathcal{E}]\to\mathcal{G}_2$。
所得线性同态 $\Omega\to\mathcal{G}_2$ 说明此时 $[\xi]=0$。

\medskip\noindent
步骤 3：$[\xi]$ 消失是充分的。设
$\theta:\Omega\to\mathcal{G}_2$ 是 $\mathcal{O}_1$-线性同态，
使 $\xi=\theta\circ(\mathcal{I}/\mathcal{I}^2\to\Omega)$。
计算表明
$$
b + \theta \circ d :
\mathcal{O}_{\mathcal{B}'_1}[\mathcal{E}]
\longrightarrow
\mathcal{O}'_2
$$
零化 $\mathcal{I}'$，因而定义使图
(\ref{defos-equation-huge-2-ringed-topoi}) 成立的同态
$\mathcal{O}'_1\to\mathcal{O}'_2$。

\medskip\noindent
上述特殊情形中 (2) 的证明从略。提示：这与引理
\ref{defos-lemma-huge-diagram} 中 (2) 的证明完全相同。
\end{proof}

\begin{lemma}
\label{defos-lemma-NL-represent-ext-class-ringed-topoi}
设 $\mathcal{C}$ 是位点，$\mathcal{A}\to\mathcal{B}$ 是
$\mathcal{C}$ 上的环层同态，$\mathcal{G}$ 是 $\mathcal{B}$-模。
设 $\xi\in\Ext^1_\mathcal{B}
(\NL_{\mathcal{B}/\mathcal{A}},\mathcal{G})$。
存在集合层同态 $\alpha:\mathcal{E}\to\mathcal{B}$，使
$\xi\in\Ext^1_\mathcal{B}(\NL(\alpha),\mathcal{G})$
是某个同态 $\mathcal{I}/\mathcal{I}^2\to\mathcal{G}$ 的类
（记号见证明）。
\end{lemma}

\begin{proof}
回忆：给定 $\alpha:\mathcal{E}\to\mathcal{B}$，使
$\mathcal{A}[\mathcal{E}]\to\mathcal{B}$ 是核为 $\mathcal{I}$ 的满射，
则复形
$\NL(\alpha) = (\mathcal{I}/\mathcal{I}^2 \to 
\Omega_{\mathcal{A}[\mathcal{E}]/\mathcal{A}}
\otimes_{\mathcal{A}[\mathcal{E}]} \mathcal{B})$
典范同构于 $\NL_{\mathcal{B}/\mathcal{A}}$，见《位点上的模》引理
\ref{sites-modules-lemma-NL-up-to-qis}。还要注意，
$\Omega = \Omega_{\mathcal{A}[\mathcal{E}]/\mathcal{A}}
\otimes_{\mathcal{A}[\mathcal{E}]} \mathcal{B}$ 是预层
$U \mapsto \bigoplus_{e \in \mathcal{E}(U)} \mathcal{B}(U)$.
的伴随层。换言之，$\Omega$ 是集合层 $\mathcal{E}$ 上的自由
$\mathcal{B}$-模；特别地，存在典范同态 $\mathcal{E}\to\Omega$。

\medskip\noindent
现在取某个 $\mathcal{E}$（例如像朴素余切复形的定义那样取
$\mathcal{E}=\mathcal{B}$）。把 $\xi$ 写成同态
$\mathcal{I}/\mathcal{I}^2\to\mathcal{G}$ 之类的障碍，是
$\Ext^1_\mathcal{B}(\Omega,\mathcal{G})$ 中的元素。
设它由 $\mathcal{B}$-模扩张
$0\to\mathcal{G}\to\mathcal{H}\to\Omega\to0$ 表示。
考虑集合层
$\mathcal{E}' = \mathcal{E} \times_\Omega \mathcal{H}$
它带有诱导同态 $\alpha':\mathcal{E}'\to\mathcal{B}$。令
$\mathcal{I}'=\Ker(\mathcal{A}[\mathcal{E}']\to\mathcal{B})$ 以及
$\Omega'=\Omega_{\mathcal{A}[\mathcal{E}']/\mathcal{A}}
\otimes_{\mathcal{A}[\mathcal{E}']}\mathcal{B}$。
$\xi$ 在拟同构 $\NL(\alpha')\to\NL(\alpha)$ 下的拉回，
在
$\Ext^1_\mathcal{B}(\Omega', \mathcal{G})$
中的像为零。这是因为，沿 $\Omega'\to\Omega$ 拉回扩张
$\mathcal{H}$ 所得扩张是分裂的：$\Omega'$ 是集合层
$\mathcal{E}'$ 上的自由 $\mathcal{B}$-模，而且按构造有交换图
$$
\xymatrix{
\mathcal{E}' \ar[r] \ar[d] & \mathcal{E} \ar[d] \\
\mathcal{H} \ar[r] & \Omega
}
$$
证明完毕。
\end{proof}

\begin{lemma}
\label{defos-lemma-choices-ringed-topoi}
若 (\ref{defos-equation-to-solve-ringed-topoi}) 存在解，则解的同构类集合是
$\Ext^1_\mathcal{O}
(\NL_{\mathcal{O}/\mathcal{O}_\mathcal{B}},\mathcal{G})$
作用下的主齐性空间。
\end{lemma}

\begin{proof}
立即注意到：给定 (\ref{defos-equation-to-solve-ringed-topoi}) 的两个解
$\mathcal{O}'_1$、$\mathcal{O}'_2$，由引理
\ref{defos-lemma-huge-diagram-ringed-topoi} 得到障碍元素
$o(\mathcal{O}'_1, \mathcal{O}'_2) \in \Ext^1_\mathcal{O}(
\NL_{\mathcal{O}/\mathcal{O}_\mathcal{B}}, \mathcal{G})$
它是存在同态 $\mathcal{O}'_1\to\mathcal{O}'_2$ 的障碍。
显然，此元素正是存在同构的障碍，因而区分同构类。
所以要完成证明，只需说明：给定一个解 $\mathcal{O}'$ 和元素
$\xi \in \Ext^1_\mathcal{O}(
\NL_{\mathcal{O}/\mathcal{O}_\mathcal{B}}, \mathcal{G})$
可找到第二个解 $\mathcal{O}'_\xi$，使
$o(\mathcal{O}',\mathcal{O}'_\xi)=\xi$。

\medskip\noindent
为类 $\xi$ 选取引理 \ref{defos-lemma-NL-represent-ext-class-ringed-topoi}
中的 $\alpha:\mathcal{E}\to\mathcal{O}$。考虑核为
$\mathcal{I}$ 的满射
$f^{-1}\mathcal{O}_\mathcal{B}[\mathcal{E}] \to \mathcal{O}$
及相应的朴素余切复形
$\NL(\alpha) = (\mathcal{I}/\mathcal{I}^2 \to
\Omega_{f^{-1}\mathcal{O}_\mathcal{B}[\mathcal{E}]/
f^{-1}\mathcal{O}_\mathcal{B}}
\otimes_{f^{-1}\mathcal{O}_\mathcal{B}[\mathcal{E}]} \mathcal{O})$.
由该引理，$\xi$ 是某个同态
$\delta:\mathcal{I}/\mathcal{I}^2\to\mathcal{G}$ 的类。
把 $\mathcal{E}$ 换成 $\mathcal{E}\times_\mathcal{O}\mathcal{O}'$
后，还可假设 $\alpha$ 分解通过同态
$\alpha':\mathcal{E}\to\mathcal{O}'$。

\medskip\noindent
这些选取决定 $f^{-1}\mathcal{O}_{\mathcal{B}'}$-代数同态
$\varphi:\mathcal{O}_{\mathcal{B}'}[\mathcal{E}]\to\mathcal{O}'$。
令 $\mathcal{I}'=\Ker(\varphi)$。注意，$\varphi$ 诱导同态
$\varphi|_{\mathcal{I}'}:\mathcal{I}'\to\mathcal{G}$，
且 $\mathcal{O}'$ 是下图中的推出：
$$
\xymatrix{
0 \ar[r] & \mathcal{G} \ar[r] & \mathcal{O}' \ar[r] &
\mathcal{O} \ar[r] & 0 \\
0 \ar[r] & \mathcal{I}' \ar[u]^{\varphi|_{\mathcal{I}'}} \ar[r] &
f^{-1}\mathcal{O}_{\mathcal{B}'}[\mathcal{E}] \ar[u] \ar[r] &
\mathcal{O} \ar[u]_{=} \ar[r] & 0
}
$$
令 $\psi:\mathcal{I}'\to\mathcal{G}$ 为同态
$\varphi|_{\mathcal{I}'}$ 与复合
$$
\mathcal{I}' \to \mathcal{I}'/(\mathcal{I}')^2 \to
\mathcal{I}/\mathcal{I}^2 \xrightarrow{\delta} \mathcal{G}.
$$
之和。沿 $\psi$ 取推出，得到另一个环扩张
$\mathcal{O}'_\xi$，它嵌入如上图。计算（从略）表明
$o(\mathcal{O}',\mathcal{O}'_\xi)=\xi$，正如所需。
\end{proof}

\begin{lemma}
\label{defos-lemma-extensions-of-relative-ringed-topoi}
设 $f:(\Sh(\mathcal{C}),\mathcal{O})\to
(\Sh(\mathcal{B}),\mathcal{O}_\mathcal{B})$ 是环化拓扑斯态射，
$\mathcal{G}$ 是 $\mathcal{O}$-模。所有
$f^{-1}\mathcal{O}_\mathcal{B}$-代数扩张
$$
0 \to \mathcal{G} \to \mathcal{O}' \to \mathcal{O} \to 0
$$
中 $\mathcal{G}$ 为平方零理想者的同构类集合\footnote{换言之，即
$(\Sh(\mathcal{B}),\mathcal{O}_\mathcal{B})$ 上一阶增厚
$i : (\Sh(\mathcal{C}), \mathcal{O}) \to (\Sh(\mathcal{C}), \mathcal{O}')$
连同 $\mathcal{O}$-模同构
$\mathcal{G}\to\Ker(i^\sharp)$ 的同构类集合。}，
与
$\Ext^1_\mathcal{O}(\NL_{\mathcal{O}/\mathcal{O}_\mathcal{B}}, \mathcal{G})$.
典范双射。
\end{lemma}

\begin{proof}
为证明此结论，把前述结果应用于由下图给出
(\ref{defos-equation-to-solve-ringed-topoi}) 的情形：
$$
\xymatrix{
0 \ar[r] &
\mathcal{G} \ar[r] &
{?} \ar[r] &
\mathcal{O} \ar[r] & 0 \\
0 \ar[r] &
0 \ar[u] \ar[r] &
f^{-1}\mathcal{O}_\mathcal{B} \ar[u] \ar[r]^{\text{id}} &
f^{-1}\mathcal{O}_\mathcal{B} \ar[u] \ar[r] & 0
}
$$
于是本引理由引理 \ref{defos-lemma-choices-ringed-topoi} 及解
$\mathcal{G}\oplus\mathcal{O}$ 的存在性得出。
（该双射的直接构造见下列注。）
\end{proof}

\begin{remark}
\label{defos-remark-extensions-of-relative-ringed-topoi}
% P11-E0009: missing Sh in frozen source codomain retained; see control/ERRATA_CANDIDATES.jsonl.
设 $f:(\Sh(\mathcal{C}),\mathcal{O})\to
(\mathcal{B},\mathcal{O}_\mathcal{B})$ 与 $\mathcal{G}$ 如引理
\ref{defos-lemma-extensions-of-relative-ringed-topoi}。考虑该引理中的扩张
$0 \to \mathcal{G} \to \mathcal{O}' \to \mathcal{O} \to 0$
。可以选取集合层 $\mathcal{E}$ 及交换图
$$
\xymatrix{
\mathcal{E} \ar[d]_{\alpha'} \ar[rd]^\alpha \\
\mathcal{O}' \ar[r] & \mathcal{O}
}
$$
使 $f^{-1}\mathcal{O}_\mathcal{B}[\mathcal{E}]\to\mathcal{O}$
是核为 $\mathcal{J}$ 的满射。（例如可取任意满射到
$\mathcal{O}'$ 的集合层。）于是
$$
\NL_{\mathcal{O}/\mathcal{O}_\mathcal{B}} \cong \NL(\alpha) = 
\left(
\mathcal{J}/\mathcal{J}^2
\longrightarrow
\Omega_{f^{-1}\mathcal{O}_\mathcal{B}[\mathcal{E}]/
f^{-1}\mathcal{O}_\mathcal{B}}
\otimes_{f^{-1}\mathcal{O}_\mathcal{B}[\mathcal{E}]} \mathcal{O}\right)
$$
见《位点上的模》第 \ref{sites-modules-section-netherlander} 节，
尤其是引理 \ref{sites-modules-lemma-NL-up-to-qis}。
当然，$\alpha'$ 决定同态
$f^{-1}\mathcal{O}_\mathcal{B}[\mathcal{E}]\to\mathcal{O}'$，
进而决定同态
$$
\mathcal{J}/\mathcal{J}^2 \longrightarrow \mathcal{G}
$$
后者又决定
$\Ext^1_\mathcal{O}(\NL(\alpha), \mathcal{G}) =
\Ext^1_\mathcal{O}(\NL_{\mathcal{O}/\mathcal{O}_\mathcal{B}}, \mathcal{G})$
中的元素；在本引理的双射下，该元素对应于 $\mathcal{O}'$。
\end{remark}

\begin{lemma}
\label{defos-lemma-extensions-of-relative-ringed-topoi-functorial}
设 $f:(\Sh(\mathcal{C}),\mathcal{O}_\mathcal{C})\to
(\Sh(\mathcal{B}),\mathcal{O}_\mathcal{B})$ 与
$g:(\Sh(\mathcal{D}),\mathcal{O}_\mathcal{D})\to
(\Sh(\mathcal{C}),\mathcal{O}_\mathcal{C})$ 是环化拓扑斯态射。
设 $\mathcal{F}$ 是 $\mathcal{O}_\mathcal{C}$-模，
$\mathcal{G}$ 是 $\mathcal{O}_\mathcal{D}$-模，且
$c:g^*\mathcal{F}\to\mathcal{G}$ 是 $\mathcal{O}_\mathcal{D}$-线性同态。
最后，考虑
\begin{enumerate}
\item[(a)]
$0 \to \mathcal{F} \to \mathcal{O}_{\mathcal{C}'} \to
\mathcal{O}_\mathcal{C} \to 0$
是对应于
$\xi\in\Ext^1_{\mathcal{O}_\mathcal{C}}
(\NL_{\mathcal{O}_\mathcal{C}/\mathcal{O}_\mathcal{B}},\mathcal{F})$
的 $f^{-1}\mathcal{O}_\mathcal{B}$-代数扩张；
\item[(b)]
$0 \to \mathcal{G} \to \mathcal{O}_{\mathcal{D}'} \to
\mathcal{O}_\mathcal{D} \to 0$
是对应于
$\zeta\in\Ext^1_{\mathcal{O}_\mathcal{D}}
(\NL_{\mathcal{O}_\mathcal{D}/\mathcal{O}_\mathcal{B}},\mathcal{G})$
的 $g^{-1}f^{-1}\mathcal{O}_\mathcal{B}$-代数扩张。
\end{enumerate}
见引理 \ref{defos-lemma-extensions-of-relative-ringed-topoi}。
存在环化拓扑斯态射
$$
g' :
(\Sh(\mathcal{D}), \mathcal{O}_{\mathcal{D}'})
\longrightarrow
(\Sh(\mathcal{C}), \mathcal{O}_{\mathcal{C}'})
$$
位于 $(\Sh(\mathcal{B}),\mathcal{O}_\mathcal{B})$ 上方且与
$g$、$c$ 相容，当且仅当 $\xi$、$\zeta$ 映到
$\Ext^1_{\mathcal{O}_\mathcal{D}}(
Lg^*\NL_{\mathcal{O}_\mathcal{C}/\mathcal{O}_\mathcal{B}}, \mathcal{G})$
中的同一元素。
\end{lemma}

\begin{proof}
该陈述有意义，因为利用同态
$Lg^*\mathcal{F}\to g^*\mathcal{F}\xrightarrow{c}\mathcal{G}$，
有同态
$$
\Ext^1_{\mathcal{O}_\mathcal{C}}(
\NL_{\mathcal{O}_\mathcal{C}/\mathcal{O}_\mathcal{B}}, \mathcal{F}) \to
\Ext^1_{\mathcal{O}_\mathcal{D}}(
Lg^*\NL_{\mathcal{O}_\mathcal{C}/\mathcal{O}_\mathcal{B}}, Lg^*\mathcal{F}) \to
\Ext^1_{\mathcal{O}_\mathcal{D}}
(Lg^*\NL_{\mathcal{O}_\mathcal{C}/\mathcal{O}_\mathcal{B}}, \mathcal{G})
$$
并且，利用同态
$Lg^*\NL_{\mathcal{O}_\mathcal{C}/\mathcal{O}_\mathcal{B}}\to
\NL_{\mathcal{O}_\mathcal{D}/\mathcal{O}_\mathcal{B}}$，有同态
% P11-E0010: undefined O_Y coefficient rings retained from frozen source; see control/ERRATA_CANDIDATES.jsonl.
$$
\Ext^1_{\mathcal{O}_Y}(
\NL_{\mathcal{O}_\mathcal{D}/\mathcal{O}_\mathcal{B}}, \mathcal{G}) \to
\Ext^1_{\mathcal{O}_Y}(
Lg^*\NL_{\mathcal{O}_\mathcal{C}/\mathcal{O}_\mathcal{B}}, \mathcal{G})
$$
把引理 \ref{defos-lemma-huge-diagram-ringed-topoi} 应用于图
$$
\xymatrix{
& 0 \ar[r] &
\mathcal{G} \ar[r] &
\mathcal{O}_{\mathcal{D}'} \ar[r] &
\mathcal{O}_\mathcal{D} \ar[r] & 0 \\
& 0 \ar[r]|\hole & 0 \ar[u] \ar[r] &
g^{-1}f^{-1}\mathcal{O}_\mathcal{B} \ar[u] \ar[r]|\hole &
g^{-1}f^{-1}\mathcal{O}_\mathcal{B} \ar[u] \ar[r] & 0 \\
0 \ar[r] &
\mathcal{F} \ar[ruu] \ar[r] &
\mathcal{O}_{\mathcal{C}'} \ar[r] &
\mathcal{O}_\mathcal{C} \ar[ruu] \ar[r] & 0 \\
0 \ar[r] & 0 \ar[ruu]|\hole \ar[u] \ar[r] &
f^{-1}\mathcal{O}_\mathcal{B} \ar[ruu]|\hole \ar[u] \ar[r] &
f^{-1}\mathcal{O}_\mathcal{B} \ar[ruu]|\hole \ar[u] \ar[r] & 0
}
$$
再结合引理 \ref{defos-lemma-extensions-of-relative-ringed-topoi} 与
\ref{defos-lemma-huge-diagram-ringed-topoi} 的证明中两个构造之间的相容性，
即可推出本引理。该相容性的陈述与证明从略。
（直接论证见下列注。）
\end{proof}

\begin{remark}
\label{defos-remark-extensions-of-relative-ringed-topoi-functorial}
设 $f : (\Sh(\mathcal{C}), \mathcal{O}_\mathcal{C}) \to
(\Sh(\mathcal{B}), \mathcal{O}_\mathcal{B})$,
$g : (\Sh(\mathcal{D}), \mathcal{O}_\mathcal{D}) \to
(\Sh(\mathcal{C}), \mathcal{O}_\mathcal{C})$,
$\mathcal{F}$,
$\mathcal{G}$,
$c : g^*\mathcal{F} \to \mathcal{G}$,
$0 \to \mathcal{F} \to \mathcal{O}_{\mathcal{C}'} \to
\mathcal{O}_\mathcal{C} \to 0$,
$\xi \in \Ext^1_{\mathcal{O}_\mathcal{C}}(
\NL_{\mathcal{O}_\mathcal{C}/\mathcal{O}_\mathcal{B}}, \mathcal{F})$,
$0 \to \mathcal{G} \to \mathcal{O}_{\mathcal{D}'} \to
\mathcal{O}_\mathcal{D} \to 0$，以及
$\zeta \in \Ext^1_{\mathcal{O}_\mathcal{D}}(
\NL_{\mathcal{O}_\mathcal{D}/\mathcal{O}_\mathcal{B}}, \mathcal{G})$
如引理 \ref{defos-lemma-extensions-of-relative-ringed-topoi-functorial}。
沿 $c:g^{-1}\mathcal{F}\to\mathcal{G}$ 取推出，可构造扩张
$$
\xymatrix{
0 \ar[r] &
\mathcal{G} \ar[r] &
\mathcal{O}'_1 \ar[r] &
g^{-1}\mathcal{O}_\mathcal{C} \ar[r] & 0 \\
0 \ar[r] &
g^{-1}\mathcal{F} \ar[u]^c \ar[r] &
g^{-1}\mathcal{O}_{\mathcal{C}'} \ar[u] \ar[r] &
g^{-1}\mathcal{O}_\mathcal{C} \ar@{=}[u] \ar[r] & 0
}
$$
沿 $g^\sharp:g^{-1}\mathcal{O}_\mathcal{C}\to
\mathcal{O}_\mathcal{D}$ 取拉回，可构造扩张
$$
\xymatrix{
0 \ar[r] &
\mathcal{G} \ar[r] &
\mathcal{O}_{\mathcal{D}'} \ar[r] &
\mathcal{O}_\mathcal{D} \ar[r] & 0 \\
0 \ar[r] &
\mathcal{G} \ar@{=}[u] \ar[r] &
\mathcal{O}'_2 \ar[u] \ar[r] &
g^{-1}\mathcal{O}_\mathcal{C} \ar[u] \ar[r] & 0
}
$$
追图可知，存在态射
$g' : (\Sh(\mathcal{D}), \mathcal{O}_{\mathcal{D}'}) \to
(\Sh(\mathcal{C}), \mathcal{O}_{\mathcal{C}'})$
位于 $(\Sh(\mathcal{B}),\mathcal{O}_\mathcal{B})$ 上方且与
$g$、$c$ 相容，当且仅当 $\mathcal{O}'_1$ 与 $\mathcal{O}'_2$
作为以 $\mathcal{G}$ 扩张 $g^{-1}\mathcal{O}_\mathcal{C}$ 的
$g^{-1}f^{-1}\mathcal{O}_\mathcal{B}$-代数扩张同构。
由引理 \ref{defos-lemma-extensions-of-relative-ringed-topoi}，这些扩张由下式左端分类：
$$
\Ext^1_{g^{-1}\mathcal{O}_\mathcal{C}}(
\NL_{g^{-1}\mathcal{O}_\mathcal{C}/g^{-1}f^{-1}\mathcal{O}_\mathcal{B}},
\mathcal{G}) =
\Ext^1_{\mathcal{O}_\mathcal{D}}(
Lg^*\NL_{\mathcal{O}_\mathcal{C}/\mathcal{O}_\mathcal{B}}, \mathcal{G})
$$
其中等式来自张量--Hom 伴随及等式
$$
\NL_{g^{-1}\mathcal{O}_\mathcal{C}/g^{-1}f^{-1}\mathcal{O}_\mathcal{B}} =
g^{-1}\NL_{\mathcal{O}_\mathcal{C}/\mathcal{O}_\mathcal{B}}
\quad\text{且}\quad
Lg^*\NL_{\mathcal{O}_\mathcal{C}/\mathcal{O}_\mathcal{B}} =
g^{-1}\NL_{\mathcal{O}_\mathcal{C}/\mathcal{O}_\mathcal{B}}
% P11-E0010: undefined O_X/O_Y coefficients retained from frozen source; see control/ERRATA_CANDIDATES.jsonl.
\otimes_{g^{-1}\mathcal{O}_X}^\mathbf{L} \mathcal{O}_Y
$$
第一个等式见《位点上的模》引理
\ref{sites-modules-lemma-pullback-NL}；第二个等式由导出拉回的定义得出。
因此，要说明引理
\ref{defos-lemma-extensions-of-relative-ringed-topoi-functorial} 成立，
只需证明 $\mathcal{O}'_1$ 对应于 $\xi$ 的像，且
$\mathcal{O}'_2$ 对应于 $\zeta$ 的像。$\xi$ 与
$\mathcal{O}'_1$ 的对应关系，立即来自注
\ref{defos-remark-extensions-of-relative-ringed-topoi} 中类 $\xi$ 的构造。
为说明 $\zeta$ 与 $\mathcal{O}'_2$ 的对应关系，先选取交换图
$$
\xymatrix{
\mathcal{E} \ar[d]_{\beta'} \ar[rd]^\beta \\
\mathcal{O}_{\mathcal{D}'} \ar[r] & \mathcal{O}_\mathcal{D}
}
$$
使 $g^{-1}f^{-1}\mathcal{O}_\mathcal{B}[\mathcal{E}]\to
\mathcal{O}_\mathcal{D}$ 是核为 $\mathcal{K}$ 的满射。再选取交换图
$$
\xymatrix{
\mathcal{E} \ar[d]_{\beta'} &
\mathcal{E}' \ar[l]^\varphi \ar[d]_{\alpha'} \ar[rd]^\alpha \\
\mathcal{O}_{\mathcal{D}'} &
\mathcal{O}'_2 \ar[l] \ar[r] &
g^{-1}\mathcal{O}_\mathcal{C}
}
$$
使 $g^{-1}f^{-1}\mathcal{O}_\mathcal{B}[\mathcal{E}']\to
g^{-1}\mathcal{O}_\mathcal{C}$ 是核为 $\mathcal{J}$ 的满射。
（例如，作为集合层直接取
$\mathcal{E}'=\mathcal{E}\amalg\mathcal{O}'_2$。）同态 $\varphi$
诱导复形同态 $\NL(\alpha)\to\NL(\beta)$
（记号见《模》第 \ref{modules-section-netherlander} 节），特别诱导
$\bar\varphi : \mathcal{J}/\mathcal{J}^2 \to \mathcal{K}/\mathcal{K}^2$.
% P11-E0011: reversed alpha/beta identifications retained from frozen source; see control/ERRATA_CANDIDATES.jsonl.
于是
$\NL(\alpha)\cong
\NL_{\mathcal{O}_\mathcal{D}/\mathcal{O}_\mathcal{B}}$
且
$\NL(\beta) \cong
\NL_{g^{-1}\mathcal{O}_\mathcal{C}/g^{-1}f^{-1}\mathcal{O}_\mathcal{B}}$
，而复形同态 $\NL(\alpha)\to\NL(\beta)$ 表示同态
$Lg^*\NL_{\mathcal{O}_\mathcal{C}/\mathcal{O}_\mathcal{B}} \to
\NL_{\mathcal{O}_\mathcal{D}/\mathcal{O}_\mathcal{B}}$
，即引理
\ref{defos-lemma-extensions-of-relative-ringed-topoi-functorial}
的陈述中所用的同态（见其证明第一部分）。现在，$\zeta$ 对应于
$\beta'$ 诱导的同态 $\mathcal{K}/\mathcal{K}^2\to\mathcal{G}$ 之类，
见注 \ref{defos-remark-extensions-of-relative-ringed-topoi}。类似地，扩张
$\mathcal{O}'_2$ 对应于 $\alpha'$ 诱导的同态
$\mathcal{J}/\mathcal{J}^2\to\mathcal{G}$。上面的交换图表明，
此同态是 $\beta'$ 诱导的同态
$\mathcal{K}/\mathcal{K}^2\to\mathcal{G}$ 与同态
$\bar\varphi:\mathcal{J}/\mathcal{J}^2\to\mathcal{K}/\mathcal{K}^2$
的复合。这证明了所需的相容性。
\end{remark}

\begin{lemma}
\label{defos-lemma-parametrize-solutions-ringed-topoi}
设 $t : (\Sh(\mathcal{B}), \mathcal{O}_\mathcal{B}) \to
(\Sh(\mathcal{B}'), \mathcal{O}_{\mathcal{B}'})$,
$\mathcal{J} = \Ker(t^\sharp)$,
$f : (\Sh(\mathcal{C}), \mathcal{O}) \to
(\Sh(\mathcal{B}), \mathcal{O}_\mathcal{B})$、$\mathcal{G}$ 以及
% P11-E0012: missing f^{-1} on J retained from frozen source; see control/ERRATA_CANDIDATES.jsonl.
$c : \mathcal{J} \to \mathcal{G}$ 如
(\ref{defos-equation-to-solve-ringed-topoi})。以
$\xi \in \Ext^1_{\mathcal{O}_\mathcal{B}}(
\NL_{\mathcal{O}_\mathcal{B}/\mathcal{O}_{\mathcal{B}'}}, \mathcal{J})$
记经引理 \ref{defos-lemma-extensions-of-relative-ringed-topoi}
对应于以 $\mathcal{J}$ 扩张 $\mathcal{O}_\mathcal{B}$ 所得扩张
$\mathcal{O}_{\mathcal{B}'}$ 的元素。解的同构类集合典范双射于同态
$$
\Ext^1_\mathcal{O}(\NL_{\mathcal{O}/\mathcal{O}_{\mathcal{B}'}},
\mathcal{G})\to
\Ext^1_\mathcal{O}(
Lf^*\NL_{\mathcal{O}_\mathcal{B}/\mathcal{O}_{\mathcal{B}'}}, \mathcal{G})
$$
在 $\xi$ 的像上方的纤维。
\end{lemma}

\begin{proof}
把引理 \ref{defos-lemma-extensions-of-relative-ringed-topoi}
应用于 $t\circ f:(\Sh(\mathcal{C}),\mathcal{O})\to
(\Sh(\mathcal{B}'), \mathcal{O}_{\mathcal{B}'})$
及 $\mathcal{O}$-模 $\mathcal{G}$，可见
$\Ext^1_\mathcal{O}(\NL_{\mathcal{O}/\mathcal{O}_{\mathcal{B}'}},
\mathcal{G})$
中的元素 $\zeta$ 参数化
$0 \to \mathcal{G} \to \mathcal{O}' \to \mathcal{O} \to 0$
这样的 $f^{-1}\mathcal{O}_{\mathcal{B}'}$-代数扩张。
把引理 \ref{defos-lemma-extensions-of-relative-ringed-topoi-functorial} 应用于
$$
(\Sh(\mathcal{C}), \mathcal{O}) \xrightarrow{f}
(\Sh(\mathcal{B}), \mathcal{O}_\mathcal{B}) \xrightarrow{t}
(\Sh(\mathcal{B}'), \mathcal{O}_{\mathcal{B}'})
$$
及 $c:\mathcal{J}\to\mathcal{G}$，可见存在态射
$$
f' : 
(\Sh(\mathcal{C}), \mathcal{O}')
\longrightarrow
(\Sh(\mathcal{B}'), \mathcal{O}_{\mathcal{B}'})
$$
位于 $(\Sh(\mathcal{B}'),\mathcal{O}_{\mathcal{B}'})$ 上方且与
$c$、$f$ 相容，当且仅当 $\zeta$ 映到 $\xi$。
这当然等价于说 $\mathcal{O}'$ 是
(\ref{defos-equation-to-solve-ringed-topoi}) 的一个解。
\end{proof}






















\section{代数空间的形变}
\label{defos-section-deformations-spaces}

\noindent
本节具体说明第 \ref{defos-section-deformations-ringed-topoi} 节的结果
对于代数空间的形变意味着什么。

\begin{lemma}
\label{defos-lemma-match-thickenings}
设 $S$ 是概形，$i:Z\to Z'$ 是 $S$ 上代数空间的态射。
下列条件等价：
\begin{enumerate}
\item $i$ 是《空间态射进阶》第
\ref{spaces-more-morphisms-section-thickenings} 节所定义的代数空间增厚；
\item 相伴的环化拓扑斯态射
$i_{small} : (\Sh(Z_\etale), \mathcal{O}_Z) \to
(\Sh(Z'_\etale), \mathcal{O}_{Z'})$
（《空间的性质》引理
\ref{spaces-properties-lemma-morphism-ringed-topoi}）
是第 \ref{defos-section-thickenings-ringed-topoi} 节意义下的增厚。
\end{enumerate}
\end{lemma}

\begin{proof}
我们强调，这并非显然之事。

\medskip\noindent
假设 (1)。由《空间态射进阶》引理
\ref{spaces-more-morphisms-lemma-thickening-equivalence}，
态射 $i$ 诱导小 étale 位点之间的等价，因而特别诱导拓扑斯之间的等价。
当然，按增厚的定义，$i^\sharp$ 是核局部幂零的满射。

\medskip\noindent
假设 (2)。（此方向不那么重要，更多是一则趣闻。）对任意 étale 态射
$Y'\to Z'$，可见 $Y=Z\times_{Z'}Y'$ 与 $Y'$ 有相同的 étale 拓扑斯。
特别地，$Y'$ 拟紧当且仅当 $Y$ 拟紧，因为拟紧性是拓扑斯论概念
（《位点》引理 \ref{sites-lemma-quasi-compact}）。由此还可见，
$Y'$ 拟紧且拟分离，当且仅当 $Y$ 拟紧且拟分离
（因为可用下述条件刻画 $Y'$ 拟分离：对于所有在 $Y'$ 上 étale 的
拟紧代数空间 $Y'_1,Y'_2$，纤维积
$Y'_1\times_{Y'}Y'_2$ 拟紧）。取 $Y'$ 仿射。
则代数空间 $Y$ 拟紧且拟分离。对于任意拟凝聚
$\mathcal{O}_Y$-模 $\mathcal{F}$，有
$H^q(Y, \mathcal{F}) = H^q(Y', (Y \to Y')_*\mathcal{F})$
，因为二者的 étale 拓扑斯相同。于是
$H^q(Y',(Y\to Y')_*\mathcal{F})=0$，因为推出是拟凝聚的
（《空间的态射》引理 \ref{spaces-morphisms-lemma-pushforward}），
且 $Y$ 仿射。
% P11-E0013: circular/reversed affine-space claims retained from frozen source; see control/ERRATA_CANDIDATES.jsonl.
由《空间的上同调》命题
\ref{spaces-cohomology-proposition-vanishing-affine}
可知 $Y'$ 仿射（本引理这一方向肯定有避开该引用的证明）。
因此 $i$ 是仿射态射。在仿射情形，由第
\ref{defos-section-thickenings-ringed-topoi} 节的条件容易得出
$i$ 是代数空间的增厚。
\end{proof}

\begin{lemma}
\label{defos-lemma-deform-spaces}
设 $S$ 是概形，$Y\subset Y'$ 是 $S$ 上代数空间的一阶增厚，
$f:X\to Y$ 是 $S$ 上代数空间的平坦态射。若存在 $S$ 上代数空间的
平坦态射 $f':X'\to Y'$ 以及 $Y$ 上的同构
$a:X\to X'\times_{Y'}Y$，则
\begin{enumerate}
\item 二元组 $(f':X'\to Y',a)$ 的同构类集合是
$\Ext^1_{\mathcal{O}_X}(\NL_{X/Y},f^*\mathcal{C}_{Y/Y'})$
作用下的主齐性空间；
\item 所有 $Y'$ 上自同构 $\varphi:X'\to X'$ 中，在
$X'\times_{Y'}Y$ 上约化为恒等者所成集合是
$\Ext^0_{\mathcal{O}_X}(\NL_{X/Y},f^*\mathcal{C}_{Y/Y'})$。
\end{enumerate}
\end{lemma}

\begin{proof}
我们把环化拓扑斯形变的材料应用于本引理各代数空间的小 étale 拓扑斯。
可把 $X$ 视为 $X'$ 的闭子空间，使
$(f,f'):(X\subset X')\to(Y\subset Y')$ 是一阶增厚之间的态射。
由引理 \ref{defos-lemma-match-thickenings}，这转化为环化拓扑斯增厚之间的态射。
由《空间态射进阶》引理
\ref{spaces-more-morphisms-lemma-deform}
（或更一般的引理 \ref{defos-lemma-deform-module-ringed-topoi}）可见，
% P11-E0014: reversed conormal indices retained from frozen source prose; see control/ERRATA_CANDIDATES.jsonl.
$X$ 在 $X'$ 中的理想层等于 $f^*\mathcal{C}_{Y'/Y}$，
而这实际上等价于 $X'$ 在 $Y'$ 上平坦。因此有交换图
$$
\xymatrix{
0 \ar[r] & f^*\mathcal{C}_{Y/Y'} \ar[r] &
\mathcal{O}_{X'} \ar[r] &
\mathcal{O}_X \ar[r] & 0 \\
0 \ar[r] &
f_{small}^{-1}\mathcal{C}_{Y/Y'} \ar[u] \ar[r] &
f_{small}^{-1}\mathcal{O}_{Y'} \ar[u] \ar[r] &
f_{small}^{-1}\mathcal{O}_Y \ar[u] \ar[r] & 0
}
$$
请与 (\ref{defos-equation-to-solve-ringed-topoi}) 比较。注意，(2) 中的自同构
$\varphi$ 给出嵌入上图的自同构
$\varphi^\sharp:\mathcal{O}_{X'}\to\mathcal{O}_{X'}$。
反过来，由《空间态射进阶》引理
\ref{spaces-more-morphisms-lemma-first-order-thickening-maps}.
，任意嵌入上方层图的自同构
$\alpha:\mathcal{O}_{X'}\to\mathcal{O}_{X'}$，
都等于 (2) 中某个自同构 $\varphi$ 的 $\varphi^\sharp$。
最后，由《空间态射进阶》引理
\ref{spaces-more-morphisms-lemma-first-order-thickening}
，若找到 $X_\etale$ 上另一个嵌入下图的环层 $\mathcal{A}$
$$
\xymatrix{
0 \ar[r] & f^*\mathcal{C}_{Y/Y'} \ar[r] &
\mathcal{A} \ar[r] &
\mathcal{O}_X \ar[r] & 0 \\
0 \ar[r] &
f_{small}^{-1}\mathcal{C}_{Y/Y'} \ar[u] \ar[r] &
f_{small}^{-1}\mathcal{O}_{Y'} \ar[u] \ar[r] &
f_{small}^{-1}\mathcal{O}_Y \ar[u] \ar[r] & 0
}
$$
，则存在一阶增厚 $X\subset X''$，满足
$\mathcal{O}_{X''}=\mathcal{A}$；再次应用《空间态射进阶》引理
\ref{spaces-more-morphisms-lemma-first-order-thickening-maps}，
就得到具备所有所需性质的态射
$(f,f''):(X\subset X'')\to(Y\subset Y')$。
因此 (1) 由引理 \ref{defos-lemma-choices-ringed-topoi} 得出，
(2) 由引理 \ref{defos-lemma-huge-diagram-ringed-topoi} 的 (2) 得出。
（注意，《空间态射进阶》第
\ref{spaces-more-morphisms-section-netherlander}
节为代数空间态射定义的 $\NL_{X/Y}$，与第
\ref{defos-section-deformations-ringed-topoi} 节所用的 $\NL_{X/Y}$ 一致。）
\end{proof}

\noindent
设 $S$ 是概形，$f:X\to B$ 是 $S$ 上代数空间的态射。
设 $\mathcal{F}\to\mathcal{G}$ 是 $\mathcal{O}_X$-模同态
（未必拟凝聚）。考虑由推出给出的函子
$$
F :
\left\{
\begin{matrix}
f^{-1}\mathcal{O}_B\text{-代数扩张}\\
0 \to \mathcal{F} \to \mathcal{O}' \to \mathcal{O}_X \to 0\\
\text{其中 }\mathcal{F}\text{ 是平方零理想}
\end{matrix}
\right\}
\longrightarrow
\left\{
\begin{matrix}
f^{-1}\mathcal{O}_B\text{-代数扩张}\\
0 \to \mathcal{G} \to \mathcal{O}' \to \mathcal{O}_X \to 0\\
\text{其中 }\mathcal{G}\text{ 是平方零理想}
\end{matrix}
\right\}
$$

\begin{lemma}
\label{defos-lemma-thickening-space-quasi-coherent}
在上述情形中，假设 $X$ 拟紧且拟分离，并且
$DQ_X(\mathcal{F})\to DQ_X(\mathcal{G})$
（《空间的导出范畴》第
\ref{spaces-perfect-section-better-coherator} 节）
是同构。那么函子 $F$ 是范畴等价。
\end{lemma}

\begin{proof}
回忆，$\NL_{X/B}$ 是 $D_\QCoh(\mathcal{O}_X)$ 的对象，见
《空间态射进阶》引理
\ref{spaces-more-morphisms-lemma-netherlander-quasi-coherent}.
因此我们的假设蕴含对每个 $i$，同态
$$
\Ext^i_X(\NL_{X/B}, \mathcal{F}) \longrightarrow
\Ext^i_X(\NL_{X/B}, \mathcal{G})
$$
都是同构。由引理 \ref{defos-lemma-huge-diagram-ringed-topoi}，
这蕴含该函子全忠实。另一方面，由引理
\ref{defos-lemma-choices-ringed-topoi}，该函子本质满射，
因为两个范畴中分别有解
$\mathcal{O}_X\oplus\mathcal{F}$ 与
$\mathcal{O}_X\oplus\mathcal{G}$。
\end{proof}

\noindent
设 $S$ 是概形，$B\subset B'$ 是 $S$ 上代数空间的一阶增厚，
其理想层为 $\mathcal{J}$。可把 $\mathcal{J}$ 视为拟凝聚
$\mathcal{O}_B$-模，也可视为 $B'$ 上的拟凝聚理想层；见
《空间态射进阶》第 \ref{spaces-more-morphisms-section-thickenings} 节。
设 $f:X\to B$ 是 $S$ 上代数空间的态射，
$\mathcal{F}\to\mathcal{G}$ 是 $\mathcal{O}_X$-模同态
（未必拟凝聚）。设
$c:f^{-1}\mathcal{J}\to\mathcal{F}$ 是
$f^{-1}\mathcal{O}_B$-模同态，并以
$c':f^{-1}\mathcal{J}\to\mathcal{G}$ 记相应复合。
考虑由推出给出的函子
$$
FT :
\{(\ref{defos-equation-to-solve-ringed-topoi})
\text{ 关于 }\mathcal{F}\text{ 与 }c\text{ 的解}\}
\longrightarrow
\{(\ref{defos-equation-to-solve-ringed-topoi})
\text{ 关于 }\mathcal{G}\text{ 与 }c'\text{ 的解}\}
$$

\begin{lemma}
\label{defos-lemma-thickening-over-thickening-space-quasi-coherent}
在上述情形中，假设 $X$ 拟紧且拟分离，并且
$DQ_X(\mathcal{F})\to DQ_X(\mathcal{G})$
（《空间的导出范畴》第
\ref{spaces-perfect-section-better-coherator} 节）
是同构。那么函子 $FT$ 是范畴等价。
\end{lemma}

\begin{proof}
(\ref{defos-equation-to-solve-ringed-topoi}) 关于 $\mathcal{F}$ 的一个解，
特别给出 $f^{-1}\mathcal{O}_{B'}$-代数扩张
$$
0 \to \mathcal{F} \to \mathcal{O}' \to \mathcal{O}_X \to 0
$$
其中 $\mathcal{F}$ 是平方零理想。$\mathcal{G}$ 的情形类似。
此外，给定这样的扩张，得到同态
$c_{\mathcal{O}'}:f^{-1}\mathcal{J}\to\mathcal{F}$。
因此所考察的是这类 $f^{-1}\mathcal{O}_{B'}$-代数扩张中
满足 $c=c_{\mathcal{O}'}$ 者所成的全子范畴。显然，若
$\mathcal{O}''=F(\mathcal{O}')$，其中 $F$ 是引理
\ref{defos-lemma-thickening-space-quasi-coherent} 的等价
（这次应用于 $X\to B'$），则 $c_{\mathcal{O}''}$ 是
$c_{\mathcal{O}'}$ 与同态 $\mathcal{F}\to\mathcal{G}$ 的复合。
本引理得证。
\end{proof}





\section{复形的形变}
\label{defos-section-deformations-complexes}

\noindent
本节为下一节作预备。我们会尽量利用交换代数各章的材料。

\begin{lemma}
\label{defos-lemma-canonical-class-algebra}
设 $R'\to R$ 是环满射，其核为平方零理想 $I$。
对于每个 $K\in D^-(R)$，$D(R)$ 中存在典范同态
$$
\omega(K) : K \longrightarrow K \otimes_R^\mathbf{L} I[2]
$$
具有下列性质：
\begin{enumerate}
\item $\omega(K)=0$ 当且仅当存在 $K'\in D(R')$，满足
$K'\otimes_{R'}^\mathbf{L}R=K$；
\item 给定 $D^-(R)$ 中的同态 $K\to L$，下图交换：
$$
\xymatrix{
K \ar[d] \ar[rr]_-{\omega(K)} & &
K \otimes^\mathbf{L}_R I[2] \ar[d] \\
L \ar[rr]^-{\omega(L)} & &
L \otimes^\mathbf{L}_R I[2]
}
$$
\item $\omega(K)$ 的形成与环同态 $R'\to S'$ 相容
（准确陈述见证明）。
\end{enumerate}
\end{lemma}

\begin{proof}
取表示 $K$ 的有上界自由 $R$-模复形 $K^\bullet$。
可选取提升 $K^n$ 的自由 $R'$-模 $(K')^n$，并选取提升
$K^\bullet$ 的微分 $d^n_K:K^n\to K^{n+1}$ 的 $R'$-模同态
$(d')^n_K:(K')^n\to(K')^{n+1}$。复合
$$
(d')^{n + 1}_K \circ (d')^n_K : (K')^n \to (K')^{n + 2}
$$
虽未必为零，却确实分解为
$$
(K')^n \to K^n \xrightarrow{\omega^n_K}
K^{n + 2} \otimes_R I = I(K')^{n + 2} \to (K')^{n + 2}
$$
，因为 $d^{n+1}\circ d^n=0$。计算表明，$\omega^n_K$
定义复形同态。此复形同态定义 $\omega(K)$。

\medskip\noindent
下面证明此构造与有上界自由 $R$-模复形同态
$\alpha^\bullet:K^\bullet\to L^\bullet$ 及给定的提升选取
$(K')^n,(L')^n,(d')^n_K,(d')^n_L$ 相容。具体地，取提升分量
$\alpha^n:K^n\to L^n$ 的同态
$(\alpha')^n:(K')^n\to(L')^n$。与前面一样，
$(d')^n_L\circ(\alpha')^n-(\alpha')^{n+1}\circ(d')_K^n$
分解为
% P11-E0015: wrong terminal cohomological degree retained from frozen source; see control/ERRATA_CANDIDATES.jsonl.
$$
(K')^n \to K^n \xrightarrow{h^n}
L^{n + 1} \otimes_R I = I(L')^{n + 1} \to (L')^{n + 2}
$$
。容易计算得到
$$
\omega^n_L \circ \alpha^n =
(d_L^{n + 1} \otimes \text{id}_I) \circ h^n + h^{n + 1} \circ d_K^n +
(\alpha^{n + 2} \otimes \text{id}_I) \circ \omega^n_K
$$
这证明了本引理 (2) 中的图在此特殊情形下交换。
把它应用于表示 $K$ 的两种不同有上界自由复形选取，
可知 $\omega(K)$ 良定义；当然，(2) 在一般情形下也成立。

\medskip\noindent
若 $K$ 提升为 $D^-(R')$ 中的 $K'$，则可用有上界自由
$R'$-模复形表示 $K'$，立即可见 $\omega(K)=0$。
反过来，回到选取 $K^\bullet,(K')^n,(d')^n_K$；若
$\omega(K)=0$，则可找到同态
$g^n:K^n\to K^{n+1}\otimes_R I$，满足
$$
\omega^n = (d_K^{n + 1} \otimes \text{id}_I) \circ g^n +
g^{n + 1} \circ d_K^n
$$
这意味着以 $(d')^n_K-g^n:(K')^n\to(K')^{n+1}$ 为微分，
得到提升 $K^\bullet$ 的自由 $R'$-模复形。这证明了 (1)。

\medskip\noindent
最后，(3) 的含义如下。设 $R'\to S'$ 是环同态，令
$S=S'\otimes_{R'}R$，并以 $J=IS'\subset S'$ 记
$S'\to S$ 的平方零核。给定 $K\in D^-(R)$，所得下图交换：
$$
\xymatrix{
K \otimes_R^\mathbf{L} S \ar[d] \ar[rr]_-{\omega(K) \otimes \text{id}} & &
(K \otimes^\mathbf{L}_R I[2]) \otimes_R^\mathbf{L} S \ar[d] \\
K \otimes_R^\mathbf{L} S \ar[rr]^-{\omega(K \otimes_R^\mathbf{L} S)} & &
(K \otimes_R^\mathbf{L} S) \otimes^\mathbf{L}_S J[2]
}
$$
其中右侧竖直箭头来自
$$
(K \otimes^\mathbf{L}_R I[2]) \otimes_R^\mathbf{L} S =
(K \otimes_R^\mathbf{L} S) \otimes_S^\mathbf{L}
(I \otimes_R^\mathbf{L} S)[2] \longrightarrow
(K \otimes_R^\mathbf{L} S) \otimes_S^\mathbf{L} J[2]
$$
如上选取 $K^\bullet,(K')^n,(d')^n_K$。于是可用
$K^\bullet\otimes_R S$、$(K')^n\otimes_{R'}S'$ 以及
$(d')^n_K\otimes\text{id}_{S'}$ 构造
$\omega(K\otimes_R^\mathbf{L}S)$。
在这些选取下，立即可在复形同态层次验证交换性。
\end{proof}







\section{环化拓扑斯上复形的形变}
\label{defos-section-thickenings-complexes}

\noindent
本节材料取自 \cite{lieblich-complexes}。

\medskip\noindent
本节材料适用于第 \ref{defos-section-thickenings-ringed-topoi} 节所定义的
环化拓扑斯一阶增厚情形。不过，为简化记号，我们假设底层位点
$\mathcal{C}$ 与 $\mathcal{D}$ 相同。此外，把环层的满同态
$\mathcal{O}' \to \mathcal{O}$ 记作
$\mathcal{O} \to \mathcal{O}_0$；这种记法在文献中或许更常见。

\begin{lemma}
\label{defos-lemma-lift-complex}
设 $\mathcal{C}$ 是位点，$\mathcal{O} \to \mathcal{O}_0$
是环层的满射。假设给定以下数据：
\begin{enumerate}
\item 平坦 $\mathcal{O}$-模 $\mathcal{G}^n$；
\item $\mathcal{O}$-模同态 $\mathcal{G}^n \to \mathcal{G}^{n + 1}$；
\item $\mathcal{O}_0$-模复形 $\mathcal{K}_0^\bullet$；
\item $\mathcal{O}$-模同态 $\mathcal{G}^n \to \mathcal{K}_0^n$。
\end{enumerate}
并且满足：
\begin{enumerate}
\item[(a)] 当 $n \gg 0$ 时，$H^n(\mathcal{K}_0^\bullet) = 0$；
\item[(b)] 当 $n \gg 0$ 时，$\mathcal{G}^n = 0$；
\item[(c)] 令
$\mathcal{G}^n_0 = \mathcal{G}^n \otimes_\mathcal{O} \mathcal{O}_0$
，则诱导同态给出复形 $\mathcal{G}_0^\bullet$ 及复形同态
$\mathcal{G}_0^\bullet \to \mathcal{K}_0^\bullet$。
\end{enumerate}
则存在：
\begin{enumerate}
\item[(\romannumeral1)]
平坦 $\mathcal{O}$-模 $\mathcal{F}^n$；
\item[(\romannumeral2)]
$\mathcal{O}$-模同态 $\mathcal{F}^n \to \mathcal{F}^{n + 1}$；
\item[(\romannumeral3)]
$\mathcal{O}$-模同态 $\mathcal{F}^n \to \mathcal{K}_0^n$；
\item[(\romannumeral4)]
$\mathcal{O}$-模同态 $\mathcal{G}^n \to \mathcal{F}^n$。
\end{enumerate}
它们满足：当 $n \gg 0$ 时 $\mathcal{F}^n = 0$；对所有 $n$，图
$$
\xymatrix{
\mathcal{G}^n \ar[r] \ar[d] & \mathcal{G}^{n + 1} \ar[d] \\
\mathcal{F}^n \ar[r] & \mathcal{F}^{n + 1}
}
$$
交换；复合
$\mathcal{G}^n \to \mathcal{F}^n \to \mathcal{K}_0^n$
等于给定同态 $\mathcal{G}^n \to \mathcal{K}_0^n$；并且令
$\mathcal{F}^n_0 = \mathcal{F}^n \otimes_\mathcal{O} \mathcal{O}_0$
时，得到复形 $\mathcal{F}_0^\bullet$ 和拟同构的复形同态
$\mathcal{F}_0^\bullet \to \mathcal{K}_0^\bullet$。
\end{lemma}

\begin{proof}
我们对 $e$ 作递降归纳，证明可对 $n \geq e$ 找到 $\mathcal{F}^n$、
$\mathcal{G}^n \to \mathcal{F}^n$ 及
$\mathcal{F}^n \to \mathcal{F}^{n + 1}$，使其嵌入交换图
$$
\xymatrix{
\ldots \ar[r] &
\mathcal{G}^{e - 1} \ar[r] \ar@/_2pc/[dd] &
\mathcal{G}^e \ar[d] \ar[r] \ar@/_2pc/[dd] &
\mathcal{G}^{e + 1} \ar[d] \ar[r] \ar@/_2pc/[dd]|\hole &
\ldots \\
& &
\mathcal{F}^e \ar[d] \ar[r] &
\mathcal{F}^{e + 1} \ar[d] \ar[r] & \ldots \\
\ldots  \ar[r] &
\mathcal{K}_0^{e - 1} \ar[r] &
\mathcal{K}_0^e \ar[r] &
\mathcal{K}_0^{e + 1} \ar[r] & \ldots
}
$$
，其中 $\mathcal{F}_0^\bullet$ 是复形，且诱导同态
$\mathcal{F}_0^\bullet \to \mathcal{K}_0^\bullet$
在 $n>e$ 时于 $H^n$ 上诱导同构，在 $n=e$ 时诱导满射。
当 $e\gg0$ 时，由假设 (a)、(b)，可对 $n\geq e$ 取
$\mathcal{F}^n=0$，故结论成立。

\medskip\noindent
归纳步骤。我们需要构造 $\mathcal{F}^{e - 1}$ 以及同态
$\mathcal{G}^{e - 1} \to \mathcal{F}^{e - 1}$、
$\mathcal{F}^{e - 1} \to \mathcal{F}^e$ 与
$\mathcal{F}^{e - 1} \to \mathcal{K}_0^{e - 1}$。
我们把 $\mathcal{F}^{e - 1} = A \oplus B \oplus C$
选为三个部分的直和。

\medskip\noindent
第一部分取 $A = \mathcal{G}^{e - 1}$，并把同态
$\mathcal{G}^{e - 1} \to \mathcal{F}^{e - 1}$ 选为第一直和项的嵌入。
同态 $A \to \mathcal{K}^{e - 1}_0$ 与
$A \to \mathcal{F}^e$ 取为显然的同态。

\medskip\noindent
为选取 $B$，考虑由归纳假设得到的满射
$$
\gamma :
\Ker(\mathcal{F}^e_0 \to \mathcal{F}^{e + 1}_0)
\longrightarrow
\Ker(\mathcal{K}^e_0 \to \mathcal{K}^{e + 1}_0)/
\Im(\mathcal{K}^{e - 1}_0 \to \mathcal{K}^e_0)
$$
可选取集合 $I$，并对每个 $i\in I$ 选取 $\mathcal{C}$ 的对象
$U_i$ 以及截面
$s_i \in \mathcal{F}^e(U_i)$、$t_i \in \mathcal{K}^{e - 1}_0(U_i)$
，使得：
\begin{enumerate}
\item $s_i$ 映为 $\Ker(\gamma) \subset
\Ker(\mathcal{F}^e_0 \to \mathcal{F}^{e + 1}_0)$ 的截面；
\item $s_i$ 与 $t_i$ 映为 $\mathcal{K}^e_0$ 的同一截面；
\item 各截面 $s_i$ 作为 $\mathcal{O}_0$-模生成 $\Ker(\gamma)$。
\end{enumerate}
我们略去完整论证；这里用到
$\mathcal{F}^e \to \mathcal{F}^e_0$ 是集合层的满射。令
$$
B = \bigoplus\nolimits_{i \in I} j_{U_i!}\mathcal{O}_{U_i}
$$
并用 $s_i$ 与 $t_i$ 确定各直和项
$j_{U_i!}\mathcal{O}_{U_i}$ 的像，从而定义同态
$B \to \mathcal{F}^e$ 与 $B \to \mathcal{K}_0^{e - 1}$。

\medskip\noindent
取 $\mathcal{F}^{e - 1} = A \oplus B$ 及上述同态，便得到对应于
$e-1$ 的上述图，并且
$\mathcal{F}_0^\bullet \to \mathcal{K}_0^\bullet$
在 $n\geq e$ 时于 $H^n$ 上诱导同构。为了使该同态在
$H^{e-1}$ 上为满射，按如下方式选取直和项 $C$。
选取集合 $J$，并对每个 $j\in J$ 选取 $\mathcal{C}$ 的对象
$U_j$ 及 $U_j$ 上
$\Ker(\mathcal{K}^{e - 1}_0 \to \mathcal{K}^e_0)$ 的截面 $t_j$，
使这些截面作为 $\mathcal{O}_0$-模生成该核。令
$$
C = \bigoplus\nolimits_{j \in J} j_{U_j!}\mathcal{O}_{U_j}
$$
% P11-E0016: source uses s_j although only t_j was introduced; frozen symbol retained; see control/ERRATA_CANDIDATES.jsonl.
再取零同态 $C \to \mathcal{F}^e$，并用 $s_j$ 确定直和项
$j_{U_j!}\mathcal{O}_{U_j}$ 的像，从而定义同态
$C \to \mathcal{K}_0^{e - 1}$。取
$\mathcal{F}^{e - 1} = A \oplus B \oplus C$ 及上述各同态，
便完成归纳步骤。
\end{proof}

\begin{lemma}
\label{defos-lemma-canonical-class}
设 $\mathcal{C}$ 是位点，$\mathcal{O} \to \mathcal{O}_0$
是环层的满射，其核是平方零理想层 $\mathcal{I}$。
对 $D^-(\mathcal{O}_0)$ 中每个对象 $K_0$，存在典范同态
$$
\omega(K_0) :
K_0 \longrightarrow
K_0 \otimes_{\mathcal{O}_0}^\mathbf{L} \mathcal{I}[2]
$$
属于 $D(\mathcal{O}_0)$，并且对 $D^-(\mathcal{O}_0)$ 中任意同态
$K_0 \to L_0$，图
$$
\xymatrix{
K_0 \ar[d] \ar[rr]_-{\omega(K_0)} & &
(K_0 \otimes^\mathbf{L}_{\mathcal{O}_0} \mathcal{I})[2] \ar[d] \\
L_0 \ar[rr]^-{\omega(L_0)} & &
(L_0 \otimes^\mathbf{L}_{\mathcal{O}_0} \mathcal{I})[2]
}
$$
交换。
\end{lemma}

\begin{proof}
用任意 $\mathcal{O}_0$-模复形 $\mathcal{K}_0^\bullet$ 表示 $K_0$。
在所有 $n$ 上取 $\mathcal{G}^n=0$，应用引理
\ref{defos-lemma-lift-complex}。把该引理给出的同态记为
$d : \mathcal{F}^n \to \mathcal{F}^{n + 1}$。于是
$d \circ d : \mathcal{F}^n \to \mathcal{F}^{n + 2}$
模 $\mathcal{I}$ 为零。由于 $\mathcal{F}^n$ 平坦，有
$\mathcal{I}\mathcal{F}^n =
\mathcal{F}^n \otimes_{\mathcal{O}} \mathcal{I} =
\mathcal{F}^n_0 \otimes_{\mathcal{O}_0} \mathcal{I}$。
因此得到典范复形同态
$$
d \circ d : \mathcal{F}_0^\bullet
\longrightarrow
(\mathcal{F}_0^\bullet \otimes_{\mathcal{O}_0} \mathcal{I})[2]
$$
由于 $\mathcal{F}_0^\bullet$ 是平坦 $\mathcal{O}_0$-模的有上界复形，
它是 K-平坦的，因而可用于计算导出张量积。此外，按构造，复形同态
$\mathcal{F}_0^\bullet \to \mathcal{K}_0^\bullet$ 是拟同构。
因此，刚构造的同态的源与靶分别表示 $K_0$ 和
$K_0 \otimes_{\mathcal{O}_0}^\mathbf{L} \mathcal{I}[2]$
，从而得到同态 $\omega(K_0)$。

\medskip\noindent
下面证明此过程与复形同态相容。设 $\mathcal{L}_0^\bullet$
表示 $D^-(\mathcal{O}_0)$ 中另一对象，并假设
$$
\mathcal{K}_0^\bullet \longrightarrow \mathcal{L}_0^\bullet
$$
是复形同态。对复形 $\mathcal{L}_0^\bullet$、平坦模
$\mathcal{F}^n$、同态 $\mathcal{F}^n \to \mathcal{F}^{n + 1}$
以及复合
$\mathcal{F}^n \to \mathcal{K}_0^n \to \mathcal{L}_0^n$
应用引理 \ref{defos-lemma-lift-complex}（这里字母的使用次序颠倒了，谨致歉意）。
得到平坦模 $\mathcal{G}^n$，以及同态
$\mathcal{F}^n \to \mathcal{G}^n$、
$\mathcal{G}^n \to \mathcal{G}^{n + 1}$ 和
$\mathcal{G}^n \to \mathcal{L}_0^n$，它们具备该引理中的全部性质。
于是显然
$$
\xymatrix{
\mathcal{F}_0^\bullet \ar[d] \ar[r] &
(\mathcal{F}_0^\bullet \otimes_{\mathcal{O}_0} \mathcal{I})[2] \ar[d] \\
\mathcal{G}_0^\bullet \ar[r] &
(\mathcal{G}_0^\bullet \otimes_{\mathcal{O}_0} \mathcal{I})[2]
}
$$
是复形的交换图。

\medskip\noindent
为证明 $\omega(K_0)$ 良定义，假设有两个表示 $K_0$ 的
$\mathcal{O}_0$-模复形
$\mathcal{K}_0^\bullet$ 与 $(\mathcal{K}'_0)^\bullet$，以及上述类型的两组数据
$(\mathcal{F}^n, d : \mathcal{F}^n \to \mathcal{F}^{n + 1},
\mathcal{F}^n \to \mathcal{K}_0^n)$
以及
$((\mathcal{F}')^n, d : (\mathcal{F}')^n \to (\mathcal{F}')^{n + 1},
(\mathcal{F}')^n \to \mathcal{K}_0^n)$
。可以选取复形 $(\mathcal{K}''_0)^\bullet$ 和拟同构
$\mathcal{K}_0^\bullet  \to (\mathcal{K}''_0)^\bullet$
以及
$(\mathcal{K}'_0)^\bullet  \to (\mathcal{K}''_0)^\bullet$
，以实现这两个复形在导出范畴中均表示 $K_0$ 这一事实。然后把上一段的结果应用于
$$
(\mathcal{K}_0)^\bullet \oplus (\mathcal{K}'_0)^\bullet
\longrightarrow
(\mathcal{K}''_0)^\bullet
$$
，得到交换图
$$
\xymatrix{
\mathcal{F}_0^\bullet \oplus (\mathcal{F}'_0)^\bullet
\ar[d] \ar[r] &
(\mathcal{F}_0^\bullet \otimes_{\mathcal{O}_0} \mathcal{I})[2] \oplus
((\mathcal{F}'_0)^\bullet \otimes_{\mathcal{O}_0} \mathcal{I})[2] \ar[d] \\
\mathcal{G}_0^\bullet \ar[r] &
(\mathcal{G}_0^\bullet \otimes_{\mathcal{O}_0} \mathcal{I})[2]
}
$$
由于竖直箭头在各直和项上给出拟同构，便得到
$D(\mathcal{O}_0)$ 中所需的交换性。

\medskip\noindent
既已证明良定义性，与同态相容的断言即由第二段的结果推出。
\end{proof}

\begin{lemma}
\label{defos-lemma-induced-map}
设 $(\mathcal{C}, \mathcal{O})$ 是环化位点，
$\alpha : K \to L$ 是 $D^-(\mathcal{O})$ 中的同态，
$\mathcal{F}$ 是 $\mathcal{O}$-模层，且 $n \in \mathbf{Z}$。
\begin{enumerate}
\item 若对 $i \geq n$，$H^i(\alpha)$ 是同构，则对 $i \geq n$，
$H^i(\alpha \otimes_\mathcal{O}^\mathbf{L} \text{id}_\mathcal{F})$
也是同构。
\item 若对 $i > n$，$H^i(\alpha)$ 是同构，且对 $i = n$ 它是满射，
则对 $i > n$，
$H^i(\alpha \otimes_\mathcal{O}^\mathbf{L} \text{id}_\mathcal{F})$
是同构，且对 $i = n$ 它是满射。
\end{enumerate}
\end{lemma}

\begin{proof}
选取可区别三角
$$
K \to L \to C \to K[1]
$$
在情形 (2) 中，对 $i\geq n$ 有 $H^i(C)=0$。因此由
《导出范畴》引理 \ref{derived-lemma-negative-vanishing}
的对偶，对 $i\geq n$ 有
$H^i(C \otimes_\mathcal{O}^\mathbf{L} \mathcal{F}) = 0$。
这又表明
$H^i(\alpha \otimes_\mathcal{O}^\mathbf{L} \text{id}_\mathcal{F})$
在 $i>n$ 时是同构，在 $i=n$ 时是满射。
在情形 (1) 中，还可知 $H^{n - 1}(L) \to H^{n - 1}(C)$
是满射。考虑图
$$
\xymatrix{
H^{n - 1}(L) \otimes_\mathcal{O} \mathcal{F} \ar[r] \ar[d] &
H^{n - 1}(C) \otimes_\mathcal{O} \mathcal{F} \ar@{=}[d] \\
H^{n - 1}(L \otimes_\mathcal{O}^\mathbf{L} \mathcal{F}) \ar[r] &
H^{n - 1}(C \otimes_\mathcal{O}^\mathbf{L} \mathcal{F})
}
$$
，可知下方水平箭头是满射。结合前述结论，这表明
$H^n(\alpha \otimes_\mathcal{O}^\mathbf{L} \text{id}_\mathcal{F})$
是同构。
\end{proof}

\begin{lemma}
\label{defos-lemma-canonical-class-obstruction}
设 $\mathcal{C}$ 是位点，$\mathcal{O} \to \mathcal{O}_0$
是环层的满射，其核是平方零理想层 $\mathcal{I}$。
对 $D^-(\mathcal{O}_0)$ 中每个对象 $K_0$，以下条件等价：
\begin{enumerate}
\item 引理 \ref{defos-lemma-canonical-class} 中构造的类
$\omega(K_0) \in
\Ext^2_{\mathcal{O}_0}(K_0, K_0 \otimes_{\mathcal{O}_0} \mathcal{I})$
为零；
\item 存在 $K \in D^-(\mathcal{O})$，使得
$K \otimes_\mathcal{O}^\mathbf{L} \mathcal{O}_0 = K_0$
在 $D(\mathcal{O}_0)$ 中成立。
\end{enumerate}
\end{lemma}

\begin{proof}
设 $K$ 如 (2) 所述。可用平坦 $\mathcal{O}$-模的有上界复形
$\mathcal{F}^\bullet$ 表示 $K$。于是
$\mathcal{F}_0^\bullet =
\mathcal{F}^\bullet \otimes_{\mathcal{O}} \mathcal{O}_0$
在 $D(\mathcal{O}_0)$ 中表示 $K_0$。由于 $\mathcal{F}^\bullet$ 是复形，
$d_{\mathcal{F}^\bullet} \circ d_{\mathcal{F}^\bullet} = 0$；
从 $\omega(K_0)$ 的构造即可看出它为零。

\medskip\noindent
假设 (1) 成立。取构造 $\omega(K_0)$ 时的
$\mathcal{F}^n$ 与 $d : \mathcal{F}^n \to \mathcal{F}^{n + 1}$。
$\omega(K_0)$ 为零意味着同态
$$
\omega = d \circ d : \mathcal{F}_0^\bullet
\longrightarrow
(\mathcal{F}_0^\bullet \otimes_{\mathcal{O}_0} \mathcal{I})[2]
$$
在 $D(\mathcal{O}_0)$ 中为零。按导出范畴作为
$\mathcal{O}_0$-模复形同伦范畴之局部化的定义，存在拟同构
$\alpha : \mathcal{G}_0^\bullet \to \mathcal{F}_0^\bullet$
以及 $\mathcal{O}_0$-模同态
$h^n : \mathcal{G}_0^n \to
\mathcal{F}_0^{n + 1} \otimes_\mathcal{O} \mathcal{I}$
，使得
$$
\omega \circ \alpha =
d_{\mathcal{F}_0^\bullet \otimes \mathcal{I}} \circ h +
h \circ d_{\mathcal{G}_0^\bullet}
$$
令
$$
\mathcal{H}^n = \mathcal{F}^n \times_{\mathcal{F}^n_0} \mathcal{G}_0^n
$$
并定义
$$
d' : \mathcal{H}^n \longrightarrow \mathcal{H}^{n + 1},\quad
(f^n, g_0^n) \longmapsto (d(f^n) - h^n(g_0^n), d(g_0^n))
$$
这里使用显然的记号，并用到
$\mathcal{F}_0^{n + 1} \otimes_{\mathcal{O}_0} \mathcal{I} =
\mathcal{F}^{n + 1} \otimes_\mathcal{O} \mathcal{I} =
\mathcal{I}\mathcal{F}^{n + 1} \subset \mathcal{F}^{n + 1}$。
由 $h^n$ 的选取和 $\omega$ 的定义可核验 $d'\circ d'=0$。
因此 $\mathcal{H}^\bullet$ 定义 $D(\mathcal{O})$ 中的对象。
另一方面，有 $\mathcal{O}$-模复形的短正合列
$$
0 \to \mathcal{F}_0^\bullet \otimes_{\mathcal{O}_0} \mathcal{I} \to
\mathcal{H}^\bullet \to \mathcal{G}_0^\bullet \to 0
$$
尚需证明
$\mathcal{H}^\bullet \otimes_\mathcal{O}^\mathbf{L} \mathcal{O}_0$
与 $K_0$ 同构。选取拟同构
$\mathcal{E}^\bullet \to \mathcal{H}^\bullet$
，其中 $\mathcal{E}^\bullet$ 是平坦 $\mathcal{O}$-模的有上界复形。
得到交换图
$$
\xymatrix{
0 \ar[r] &
\mathcal{E}^\bullet \otimes_\mathcal{O} \mathcal{I} \ar[d]^\beta \ar[r] &
\mathcal{E}^\bullet \ar[d]^\gamma \ar[r] &
\mathcal{E}_0^\bullet \ar[d]^\delta \ar[r] &
0 \\
0 \ar[r] &
\mathcal{F}_0^\bullet \otimes_{\mathcal{O}_0} \mathcal{I} \ar[r] &
\mathcal{H}^\bullet \ar[r] &
\mathcal{G}_0^\bullet \ar[r] &
0
}
$$
我们断言 $\delta$ 是拟同构。由于当 $i\gg0$ 时 $H^i(\delta)$
是同构，可以对满足“当 $i\geq n$ 时 $H^i(\delta)$ 为同构”的
$n$ 作递降归纳。注意，
$\mathcal{E}^\bullet \otimes_\mathcal{O} \mathcal{I}$
表示
$\mathcal{E}_0^\bullet \otimes_{\mathcal{O}_0}^\mathbf{L} \mathcal{I}$，
而
$\mathcal{F}_0^\bullet \otimes_{\mathcal{O}_0} \mathcal{I}$
表示
$\mathcal{G}_0^\bullet \otimes_{\mathcal{O}_0}^\mathbf{L} \mathcal{I}$，
并且作为 $D(\mathcal{O}_0)$ 中的同态，有
$\beta = \delta \otimes_{\mathcal{O}_0}^\mathbf{L} \text{id}_\mathcal{I}$
。这是因为
$\beta =
(\alpha \otimes \text{id}_\mathcal{I})
\circ
(\delta \otimes \text{id}_\mathcal{I})$。
假设在次数 $\geq n$ 上 $H^i(\delta)$ 是同构。由刚才所述及引理
\ref{defos-lemma-induced-map}，$\beta$ 也具有同一性质。考虑图
$$
\xymatrix{
H^{n - 1}(\mathcal{E}^\bullet \otimes_\mathcal{O} \mathcal{I})
\ar[r] \ar[d]^{H^{n - 1}(\beta)} &
H^{n - 1}(\mathcal{E}^\bullet) \ar[r] \ar[d] &
H^{n - 1}(\mathcal{E}_0^\bullet) \ar[r] \ar[d]^{H^{n - 1}(\delta)} &
H^n(\mathcal{E}^\bullet \otimes_\mathcal{O} \mathcal{I})
\ar[r] \ar[d]^{H^n(\beta)} &
H^n(\mathcal{E}^\bullet) \ar[d] \\
H^{n - 1}(\mathcal{F}_0^\bullet \otimes_\mathcal{O} \mathcal{I}) \ar[r] &
H^{n - 1}(\mathcal{H}^\bullet) \ar[r] &
H^{n - 1}(\mathcal{G}_0^\bullet) \ar[r] &
H^n(\mathcal{F}_0^\bullet \otimes_\mathcal{O} \mathcal{I}) \ar[r] &
H^n(\mathcal{H}^\bullet)
}
$$
由《同调代数》引理 \ref{homology-lemma-four-lemma}，
$H^{n - 1}(\delta)$ 是满射。再由引理 \ref{defos-lemma-induced-map}，
$H^{n - 1}(\beta)$ 是满射。再次应用《同调代数》引理
\ref{homology-lemma-four-lemma}，可知 $H^{n - 1}(\delta)$ 是同构。
归纳证明了该断言，故 $\delta$ 是拟同构，这正是所需结论。
\end{proof}

\begin{lemma}
\label{defos-lemma-lift-map-complexes}
设 $\mathcal{C}$ 是位点，$\mathcal{O} \to \mathcal{O}_0$
是环层的满射。假设给定以下数据：
\begin{enumerate}
\item $\mathcal{O}$-模复形 $\mathcal{F}^\bullet$；
\item $\mathcal{O}_0$-模复形 $\mathcal{K}_0^\bullet$；
\item 拟同构 $\mathcal{K}_0^\bullet \to
\mathcal{F}^\bullet \otimes_\mathcal{O} \mathcal{O}_0$。
\end{enumerate}
则存在拟同构 $\mathcal{G}^\bullet \to \mathcal{F}^\bullet$，使得复形同态
$\mathcal{G}^\bullet  \otimes_\mathcal{O} \mathcal{O}_0 \to
\mathcal{F}^\bullet \otimes_\mathcal{O} \mathcal{O}_0$
在 $\mathcal{O}_0$-模复形的同伦范畴中经由
$\mathcal{K}_0^\bullet$ 分解。
\end{lemma}

\begin{proof}
令 $\mathcal{F}_0^\bullet =
\mathcal{F}^\bullet \otimes_\mathcal{O} \mathcal{O}_0$。
由《导出范畴》引理 \ref{derived-lemma-make-surjective}，
给定同态有分解
$$
\mathcal{K}_0^\bullet \to \mathcal{L}_0^\bullet \to \mathcal{F}_0^\bullet
$$
，其中第一箭头有精确到同伦的逆，第二箭头逐项为分裂满射。
因此可以假设 $\mathcal{K}_0^\bullet \to \mathcal{F}_0^\bullet$
逐项为满射。此时取
$$
\mathcal{G}^n = \mathcal{F}^n \times_{\mathcal{F}^n_0} \mathcal{K}_0^n
$$
，一切即显然。
\end{proof}

\begin{lemma}
\label{defos-lemma-inf-obs-map-defo-complex}
设 $\mathcal{C}$ 是位点，$\mathcal{O} \to \mathcal{O}_0$
是环层的满射，其核是平方零理想层 $\mathcal{I}$。
设 $K,L\in D^-(\mathcal{O})$，并在 $D^-(\mathcal{O}_0)$ 中令
$K_0 = K \otimes_\mathcal{O}^\mathbf{L} \mathcal{O}_0$ 及
$L_0 = L \otimes_\mathcal{O}^\mathbf{L} \mathcal{O}_0$。
给定 $D(\mathcal{O}_0)$ 中的 $\alpha_0 : K_0 \to L_0$，存在典范元素
$$
o(\alpha_0) \in \Ext^1_{\mathcal{O}_0}(K_0,
L_0 \otimes_{\mathcal{O}_0}^\mathbf{L} \mathcal{I})
$$
；它为零当且仅当 $D(\mathcal{O})$ 中存在同态
$\alpha : K \to L$，满足
$\alpha_0 = \alpha \otimes_\mathcal{O}^\mathbf{L} \text{id}$。
\end{lemma}

\begin{proof}
寻找提升 $\alpha_0$ 的 $\alpha:K\to L$，等价于寻找
$\alpha:K\to L$，使复合 $K \xrightarrow{\alpha} L \to L_0$
等于复合 $K \to K_0 \xrightarrow{\alpha_0} L_0$。短正合列
$0 \to \mathcal{I} \to \mathcal{O} \to \mathcal{O}_0 \to 0$
给出典范可区别三角
$$
L \otimes_\mathcal{O}^\mathbf{L} \mathcal{I} \to
L \to
L_0 \to
(L \otimes_\mathcal{O}^\mathbf{L} \mathcal{I})[1]
$$
属于 $D(\mathcal{O})$。由《导出范畴》引理
\ref{derived-lemma-representable-homological}，复合
$$
K \to K_0 \xrightarrow{\alpha_0} L_0 \to
(L \otimes_\mathcal{O}^\mathbf{L} \mathcal{I})[1]
$$
为零当且仅当可找到提升 $\alpha_0$ 的 $\alpha:K\to L$。
此复合是下列群的元素：
$$
\Hom_{D(\mathcal{O})}(K, (L \otimes_\mathcal{O}^\mathbf{L} \mathcal{I})[1]) =
\Hom_{D(\mathcal{O}_0)}(K_0,
(L \otimes_\mathcal{O}^\mathbf{L} \mathcal{I})[1]) =
\Ext^1_{\mathcal{O}_0}(K_0,
L_0 \otimes_{\mathcal{O}_0}^\mathbf{L} \mathcal{I})
$$
，其中等号来自伴随。
\end{proof}

\begin{lemma}
\label{defos-lemma-first-order-defos-complex}
设 $\mathcal{C}$ 是位点，$\mathcal{O} \to \mathcal{O}_0$
是环层的满射，其核是平方零理想层 $\mathcal{I}$。
% P11-E0017: K_0 is typed over O although the lift is identified in D(O_0); frozen category retained; see control/ERRATA_CANDIDATES.jsonl.
设 $K_0 \in D^-(\mathcal{O})$。$K_0$ 的一个提升是二元组
$(K,\alpha_0)$，其中 $K$ 是 $D^-(\mathcal{O})$ 中的对象，而
$\alpha_0 : K \otimes_\mathcal{O}^\mathbf{L} \mathcal{O}_0 \to K_0$
是 $D(\mathcal{O}_0)$ 中的同构。
\begin{enumerate}
\item 给定提升 $(K,\alpha)$，该二元组的自同构群典范地等于同态
$$
\Ext^{-1}_{\mathcal{O}_0}(K_0, K_0)
\longrightarrow
\Hom_{\mathcal{O}_0}(K_0, K_0 \otimes_{\mathcal{O}_0}^\mathbf{L} \mathcal{I})
$$
的余核。
\item 若存在提升，则提升的同构类之集合在
$\Ext^1_{\mathcal{O}_0}(K_0,
K_0 \otimes_{\mathcal{O}_0}^\mathbf{L} \mathcal{I})$
的作用下为主齐性空间。
\end{enumerate}
\end{lemma}

\begin{proof}
$(K,\alpha)$ 的一个自同构是 $D(\mathcal{O})$ 中满足
$\varphi \otimes_\mathcal{O} \text{id}_{\mathcal{O}_0} = \text{id}$
的同态 $\varphi:K\to K$。这等价于说复合
$$
K \xrightarrow{\varphi - \text{id}} K \to
K \otimes_\mathcal{O}^\mathbf{L} \mathcal{O}_0
$$
为零。由此可知，自同构群是同态
$$
\Hom_\mathcal{O}(K, K_0[-1])
\longrightarrow
\Hom_\mathcal{O}(K, K_0 \otimes_{\mathcal{O}_0}^\mathbf{L} \mathcal{I})
$$
的余核；这是由 $D(\mathcal{O})$ 中的可区别三角
$$
K \otimes_\mathcal{O}^\mathbf{L} \mathcal{I} \to
K \to
K \otimes_\mathcal{O}^\mathbf{L} \mathcal{O}_0 \to
(K \otimes_\mathcal{O}^\mathbf{L} \mathcal{I})[1]
$$
及《导出范畴》引理 \ref{derived-lemma-representable-homological}
得出的。要把它改写为引理中的群，使用限制函子
$D(\mathcal{O}_0) \to D(\mathcal{O})$ 与
$- \otimes_\mathcal{O} \mathcal{O}_0 : D(\mathcal{O}) \to D(\mathcal{O}_0)$
之间的伴随。这证明了 (1)。

\medskip\noindent
证明 (2)。假设在 $D(\mathcal{O})$ 中
$K_0 = K \otimes_\mathcal{O}^\mathbf{L} \mathcal{O}_0$。
由引理 \ref{defos-lemma-inf-obs-map-defo-complex}，把提升
$(K',\alpha_0)$ 映到提升 $\alpha_0$ 的障碍 $o(\alpha_0)$ 的同态，
定义了从二元组的同构类集合到
$\Ext^1_{\mathcal{O}_0}(K_0,
K_0 \otimes_{\mathcal{O}_0}^\mathbf{L} \mathcal{I})$
的典范单射。为完成证明，只需证明它为满射。
取本引理的 $\Ext^1$ 中的
$\xi : K_0 \to (K_0 \otimes_{\mathcal{O}_0}^\mathbf{L} \mathcal{I})[1]$。
选取表示 $K$ 的平坦 $\mathcal{O}$-模的有上界复形
$\mathcal{F}^\bullet$。同态 $\xi$ 可表示为 $t\circ s^{-1}$，其中
$s : \mathcal{K}_0^\bullet \to \mathcal{F}_0^\bullet$ 是拟同构，而
$t : \mathcal{K}_0^\bullet \to
\mathcal{F}_0^\bullet \otimes_{\mathcal{O}_0} \mathcal{I}[1]$
是复形同态。由引理 \ref{defos-lemma-lift-map-complexes}，可以假设存在
$\mathcal{O}$-模复形间的拟同构
$\mathcal{G}^\bullet \to \mathcal{F}^\bullet$，使得
$\mathcal{G}_0^\bullet \to \mathcal{F}_0^\bullet$
在精确到同伦的意义下经由 $s$ 分解。我们可以并且确实把
$\mathcal{G}^\bullet$ 换成平坦 $\mathcal{O}$-模的有上界复形
（选取从这样一个复形到 $\mathcal{G}^\bullet$ 的拟同构并作替换）。
于是 $\xi$ 由复形同态
$t : \mathcal{G}_0^\bullet \to
\mathcal{F}_0^\bullet \otimes_{\mathcal{O}_0} \mathcal{I}[1]$
及拟同构 $\mathcal{G}_0^\bullet \to \mathcal{F}_0^\bullet$ 表示。令
$$
\mathcal{H}^n = \mathcal{F}^n \times_{\mathcal{F}_0^n} \mathcal{G}_0^n
$$
，其微分为
$$
\mathcal{H}^n \to \mathcal{H}^{n + 1},\quad
(f^n, g_0^n) \mapsto (d(f^n) + t(g_0^n), d(g_0^n))
$$
这是有意义的，因为
$\mathcal{F}_0^{n + 1} \otimes_{\mathcal{O}_0} \mathcal{I} =
\mathcal{F}^{n + 1} \otimes_\mathcal{O} \mathcal{I} =
\mathcal{I}\mathcal{F}^{n + 1} \subset \mathcal{F}^{n + 1}$。
我们略去证明 $\mathcal{H}^\bullet$ 是 $\mathcal{O}$-模复形的计算。
按构造有短正合列
$$
0 \to \mathcal{F}_0^\bullet \otimes_{\mathcal{O}_0} \mathcal{I} \to
\mathcal{H}^\bullet \to \mathcal{G}_0^\bullet \to 0
$$
，它由 $\mathcal{O}$-模复形构成。与引理
\ref{defos-lemma-canonical-class-obstruction} 的证明完全相同，
可证明此列诱导同构
$\alpha_0 :
\mathcal{H}^\bullet \otimes_\mathcal{O}^\mathbf{L} \mathcal{O}_0 \to
\mathcal{G}_0^\bullet$，属于 $D(\mathcal{O}_0)$。
换言之，我们构造出了二元组 $(\mathcal{H}^\bullet,\alpha_0)$。
我们略去 $o(\alpha_0)=\xi$ 的核验；提示：在此情形下，
由于存在提升 $\alpha_0$ 各分量的同态
$\mathcal{H}^n \to \mathcal{F}^n$（它们与微分不相容），
可以显式计算 $o(\alpha_0)$。证明完毕。
\end{proof}
